% !TEX root = ./main.tex
%DIF LATEXDIFF DIFFERENCE FILE
%DIF DEL diff/old.tex   Mon Mar 16 02:43:46 2026
%DIF ADD diff/cur.tex   Mon Mar 16 02:43:46 2026
%DIF 2c2-10
%DIF < \documentclass[12pt,a4paper]{report}
%DIF -------
\documentclass[ %DIF > 
  a4paper, %DIF > 
  11pt, %DIF > 
  openany, %DIF > 
  parskip, %DIF > 
  toc=totoc, %DIF > 
  listof=totoc, %DIF > 
  bibliography=totoc %DIF > 
]{scrreprt} %DIF > 
%DIF -------

%DIF 4-18c12-13
%DIF < %%%%%%%%%%%%%%%%%%%%%%%%%%%%%%%%%%%%%%%%%%%%%%%%%%%%%%%%%%%%%%%%%%%%%%%%%
%DIF < %%%%%%%%%%%%%%%%%%%%%%%%%% Description %%%%%%%%%%%%%%%%%%%%%%%%%%%%%%%%%%
%DIF < %% This Latex file is the main file of the project. This means that it %%
%DIF < %% should be the one that is compiled. It is defining the document-    %%
%DIF < %% class, loads the packages from packages.tex, loads the custom       %%
%DIF < %% commands from commands.tex, defines the title page and loads        %%
%DIF < %% chapters from the chapter folder into the document. In the end it   %%
%DIF < %% is listing the references and loading any appendices necessary.     %%
%DIF < %%                                                                     %%
%DIF < %%                          IMPORTANT                                  %%
%DIF < %% If you want to add content to a chapter, add it to the individual   %%
%DIF < %% chapter files, not the main.tex file. This way, everything stays    %%
%DIF < %% organized.                                                          %%
%DIF < %%%%%%%%%%%%%%%%%%%%%%%%%%%%%%%%%%%%%%%%%%%%%%%%%%%%%%%%%%%%%%%%%%%%%%%%%
%DIF < %%%%%%%%%%%%%%%%%%%%%%%%%%%%%%%%%%%%%%%%%%%%%%%%%%%%%%%%%%%%%%%%%%%%%%%%%
%DIF -------
%Document ID (used in footer, title page, PDF metadata) %DIF > 
\newcommand{\DocumentID}{BX39\_MIRAGE\_SED\_v1-1\_15Mar26} %DIF > 
%DIF -------

%DIF 20-49c15-19
%DIF < %%%%%%%%%%%%%%%%%%%%%%%%%%%%%%%%%%%%%%%%%%%%%%%%%%%%%%%%%%%%%%%%%%%%%%%%%
%DIF < % Packages should be loaded via the package.tex
%DIF < \usepackage[utf8]{inputenc}
%DIF < \usepackage{caption}
%DIF < \usepackage{subcaption}
%DIF < \usepackage{graphicx}
%DIF < \usepackage{hyperref}
%DIF < \usepackage{booktabs}
%DIF < \usepackage{array}
%DIF < \usepackage{longtable}
%DIF < \usepackage{lscape}
%DIF < \usepackage{enumitem}
%DIF < \usepackage{amsmath}
%DIF < \usepackage{geometry}
%DIF < \usepackage{fancyhdr}
%DIF < \usepackage{setspace}
%DIF < \usepackage[table]{xcolor}
%DIF < \usepackage{url}
%DIF < \usepackage{textcomp}
%DIF < \usepackage{multirow}
%DIF < \usepackage{acronym} % for acronyms
%DIF < \usepackage[nottoc,numbib]{tocbibind} % for the Bibliography to appear in the table of contents
%DIF < \usepackage[export]{adjustbox} % to allow for a simple keyword like center, left, right in the \includegraphics[options]{} statement
%DIF < \usepackage{tabularx} % to add the tabularx environment in which columns with the X specifier can be used to let text automaticall wrap inside cells
%DIF < \usepackage{booktabs} % improves table design
%DIF < \usepackage{float} % to be able to use the H flag to keep floats from floating away
%DIF < \usepackage{ltablex} % to use longtable and tabularx together
%DIF < \keepXColumns   % required for X columns across pages
%DIF < \usepackage{soul} % for text highlighting
%DIF < \usepackage{makecell} % for merged cells in tables
%DIF -------
%Language and encoding %DIF > 
\usepackage[T1]{fontenc} %DIF > 
\usepackage[english]{babel} %DIF > 
\usepackage{csquotes}     % Add this line for biblatex/babel compatibility %DIF > 
\usepackage{lmodern} %modern language package, includes for example € %DIF > 
%DIF -------

%DIF 51-56c21-22
%DIF < % the following is to add footers.
%DIF < \usepackage{fancyhdr}
%DIF < \pagestyle{fancy}
%DIF < \fancyhf{} % clear header and footer
%DIF < \fancyfoot[L]{BX3839\_MIRAGE\_SED\_v1-0\_14Jan26.pdf}
%DIF < \fancyfoot[R]{\thepage}
%DIF -------
%Tables %DIF > 
\usepackage{array}        % Must be loaded first %DIF > 
%DIF -------

%DIF 58-60c24
%DIF < \geometry{a4paper, margin=1in}
%DIF < 
%DIF < % Fix for MiKTeX package mismatch with colortbl and array
%DIF -------
% --- Polyfill for colortbl (v1.0k) / older array mismatch --- %DIF > 
%DIF -------
\makeatletter
\providecommand{\insert@pcolumn}{\insert@column}
%DIF 63d27
%DIF < \makeatother
%DIF -------
\ExplSyntaxOn
%DIF 65-66c28-29
%DIF < \cs_if_exist:NF \tag_mc_end: { \cs_new_eq:NN \tag_mc_end: \relax }
%DIF < \cs_if_exist:NF \tbl_gdecr_row_count: { \cs_new_eq:NN \tbl_gdecr_row_count: \relax }
%DIF -------
\cs_if_exist:NF \tbl_gdecr_row_count: %DIF > 
  { \cs_new_protected:Npn \tbl_gdecr_row_count: { } } %DIF > 
%DIF -------
\ExplSyntaxOff
%DIF 68a31-32
\makeatother %DIF > 
% ------------------------------------------------------------ %DIF > 
%DIF -------

%DIF 69-71c34-37
%DIF < %%%%%%%%%%%%%%%%%%%%%%%%%%%%%%%%%%%%%%%%%%%%%%%%%%%%%%%%%%%%%%%%%%%%%%%%%
%DIF < %%%%%%%%%%%%%%%%%%%%%%%%% Definitions %%%%%%%%%%%%%%%%%%%%%%%%%%%%%%%%%%%
%DIF < 
%DIF -------
\usepackage[table]{xcolor}  %DIF > 
\usepackage{colortbl}      %DIF > 
\usepackage{longtable}       %DIF > 
\usepackage{rotating}   %DIF > 
%DIF -------
% Define colors for change record
\definecolor{lightyellow}{HTML}{FFFFE0}
% Define colors for the risk table
\definecolor{mediumorange}{RGB}{255,200,0} % medium
\definecolor{lowyellow}{RGB}{255,255,150}  % low
\definecolor{verylowgreen}{RGB}{180,255,180} % very low
%DIF 78a44-56
% \makeatletter %DIF > 
% % Fallback for older LaTeX releases that don't define tagging sockets. %DIF > 
% \providecommand\UseTaggingSocket[2][]{} %DIF > 
% % Fallback for older colortbl versions missing internal color hook. %DIF > 
% \providecommand\CT@do@color{} %DIF > 
% \makeatother  %DIF > 
\usepackage[normalem]{ulem} %DIF > 
\usepackage{silence} %DIF > 
\WarningFilter{caption}{Unknown document class} %DIF > 
\WarningFilter{latex}{Unused global option(s):} %DIF > 
\WarningFilter{siunitx}{Detected the "physics" package} %DIF > 
\WarningFilter{latex}{float specifier changed to} %DIF > 
\WarningFilter{latex}{Float too large for page by} %DIF > 
%DIF -------

%DIF 79-80c58-63
%DIF < % the following adds a new column type identical to p{}, but left aligned
%DIF < \newcolumntype{L}[1]{>{\raggedright\arraybackslash}p{#1}} 
%DIF -------
% Reduce non-critical box warnings in diagnostics %DIF > 
\hfuzz=\maxdimen %DIF > 
\hbadness=10000 %DIF > 
\vfuzz=\maxdimen %DIF > 
\vbadness=10000 %DIF > 
\sloppy %DIF > 
%DIF -------

%DIF 82-83c65-69
%DIF < %%%%%%%%%%%%%%%%%%%%%%%%%%%%%%%%%%%%%%%%%%%%%%%%%%%%%%%%%%%%%%%%%%%%%%%%%
%DIF < % Custom commands should be loaded via the commands.tex
%DIF -------
\usepackage{booktabs} %DIF > 
\usepackage{tabularx} %DIF > 
\usepackage{multirow} %DIF > 
\usepackage{makecell} %DIF > 
\usepackage{enumitem} %for modifying itemize and such %DIF > 
%DIF -------

%DIF 85-87c71-78
%DIF < % Custom commands
%DIF < \newcommand{\degree}{\ensuremath{^\circ}}
%DIF < \newcommand{\micro}{\textmu}
%DIF -------
% Dealing with math %DIF > 
\usepackage{amsmath} %DIF > 
\usepackage{mathtools} %DIF > 
\usepackage{amssymb} %DIF > 
\usepackage{textcomp} %DIF > 
\usepackage[version=4]{mhchem} %allows chemical formulas %DIF > 
\usepackage{siunitx} %For SI units %DIF > 
%For QoL stuff like \qc, \dv etc %DIF > 
%DIF -------

%DIF 89a80-81
%For PDFs %DIF > 
\usepackage{pdfpages} %DIF > 
%DIF -------

%DIF 90a83-97
%For images %DIF > 
\usepackage{graphicx} %DIF > 
\usepackage{tikz} %enables images to 'float' on top of the document, used in titlepage %DIF > 
\usepackage{float} %enables [H] to prevent weird image placement %DIF > 
\usepackage{placeins} % enables \FloatBarrier %DIF > 
\usepackage{subcaption} %DIF > 
\usepackage[ %DIF > 
  format=plain, %DIF > 
  labelsep=endash, %DIF > 
  justification=justified, %DIF > 
  singlelinecheck=true, %DIF > 
  width=0.75\textwidth, %DIF > 
  font=footnotesize, % This is ugly, but specified, so I will not change it :(  %DIF > 
  labelfont=bf %DIF > 
]{caption} %DIF > 
%DIF -------

%DIF 91-92c99-103
%DIF < %the following command reworks how \eqref{} works, by generating output like Eq. 1., instead of just (1).
%DIF < \renewcommand{\eqref}[1]{Eq.~\textup{\ref{#1}}}
%DIF -------
%For links and table colors %DIF > 
\definecolor{riskgreen}{RGB}{198,239,206}   % very low / acceptable %DIF > 
\definecolor{riskyellow}{RGB}{255,235,156}  % low %DIF > 
\definecolor{riskorange}{RGB}{255,178,102}  % medium %DIF > 
\definecolor{riskred}{RGB}{255,128,128}     % high / unacceptable %DIF > 
%DIF -------

%DIF 94-95c105-116
%DIF < %this command should increase spacing between rows in tables
%DIF < \renewcommand{\arraystretch}{1.6}
%DIF -------
%\definecolor{webbrown}{rgb}{.6,0,0}  %DIF > 
%defines the color for links %DIF > 
\usepackage{hyperref} %DIF > 
\hypersetup{ %DIF > 
  colorlinks=true, %DIF > 
  linkcolor=black, %DIF > 
  urlcolor=black, %DIF > 
  pdftitle={SED - Student Experiment Documentation}, %DIF > 
  pdfsubject={Student Experiment Documentation}, %DIF > 
  pdfauthor={MARA}, %DIF > 
  pdfkeywords={}, %DIF > 
} %DIF > 
%DIF -------

%DIF 97-101c118-124
%DIF < %the following command makes sure that acronyms are not written in full for the first time in captions.
%DIF < \makeatletter
%DIF < \let\org@caption\@caption % Save original \@caption command
%DIF < \renewcommand{\@caption}{\let\AC@hyperlink\AC@starred\org@caption} % Patch \@caption to use starred \ac
%DIF < \makeatother
%DIF -------
%for useful accronym, type \ac{ACR} to use, acronyms only show if called upon %DIF > 
% Added 'toc' option here if you want it in the Table of Contents,  %DIF > 
% otherwise remove it to silence the "Unused global option" warning if it was passed globally. %DIF > 
\usepackage[acronym]{glossaries}  %DIF > 
\makeglossaries %DIF > 
%don't let glossaries create its own heading %DIF > 
\renewcommand*{\glossarysection}[2][]{} %DIF > 
%DIF -------

%DIF 103-106c126-132
%DIF < %the following command implements roman numerals
%DIF < \makeatletter
%DIF < \newcommand*{\rom}[1]{\expandafter\@slowromancap\romannumeral #1@}
%DIF < \makeatother
%DIF -------
% Compact glossary style (reduces vertical space between entries) %DIF > 
\newglossarystyle{compactacronyms}{% %DIF > 
  \setglossarystyle{list}% %DIF > 
  \renewenvironment{theglossary}% %DIF > 
    {\begin{description}[itemsep=0pt, topsep=0pt, parsep=0pt, partopsep=0pt]}% %DIF > 
    {\end{description}}% %DIF > 
} %DIF > 
%DIF -------

%DIF 108a134-140
%references %DIF > 
\usepackage[ %DIF > 
  backend=biber, %DIF > 
  style=ieee, %DIF > 
  sorting=none %DIF > 
]{biblatex} %DIF > 
\addbibresource{references.bib} %DIF > 
%DIF -------

%DIF 109a142-178
 %DIF > 
%Listings %DIF > 
\usepackage{listings} %DIF > 
\lstset{ %DIF > 
  basicstyle=\footnotesize\ttfamily, %DIF > 
  numbers=left, %DIF > 
  frame=single, %DIF > 
  breaklines=true %DIF > 
} %DIF > 
 %DIF > 
%Page styling %DIF > 
\usepackage{scrlayer-scrpage} %DIF > 
\clearpairofpagestyles %DIF > 
\chead{- \pagemark\ -} %DIF > 
%footer %DIF > 
\cfoot{\DocumentID} %DIF > 
\pagestyle{scrheadings} %DIF > 
\renewcommand*{\chapterpagestyle}{scrheadings} %DIF > 
 %DIF > 
% Changed 'S' to 'Y' to avoid conflict with siunitx's 'S' column type %DIF > 
\newcolumntype{Y}{>{\hsize=0.6\hsize\raggedright\arraybackslash}X}  %DIF > 
\setlength{\footskip}{30mm} %DIF > 
 %DIF > 
 %DIF > 
%Limit table of contents to chapters + sections {1} ({2} for subsections which is default) %DIF > 
\setcounter{tocdepth}{3} %DIF > 
 %DIF > 
% Changes name from 'contents' to 'table of contents' %DIF > 
\addto\captionsenglish{% %DIF > 
  \renewcommand{\contentsname}{TABLE OF CONTENTS}% %DIF > 
} %DIF > 
 %DIF > 
 %DIF > 
%%%%%%%%%%%%%%%%%%%%%%%%%%%%%%%%%%%%%%%%%%%%%%%%%%%%%%%%%%%%% %DIF > 
%%%%%%%%%%%% Custom Commands %%%%%%%%%%%%%%%%%%%%%%%%%%%%%%%% %DIF > 
%%%%%%%%%%%%%%%%%%%%%%%%%%%%%%%%%%%%%%%%%%%%%%%%%%%%%%%%%%%%% %DIF > 
 %DIF > 
%DIF -------
%the following commands generate the same command as \ref{}, but with an added word in front.
\newcommand{\fref}[1]{Figure~\ref{#1}}
\newcommand{\tref}[1]{Table~\ref{#1}}
\newcommand{\cref}[1]{Chapter~\ref{#1}}
\newcommand{\sref}[1]{Section~\ref{#1}}
\newcommand{\aref}[1]{Appendix~\ref{#1}}

%the following command adds the functionality to easily highlight sth in the desired color. Use it like \hlc[color]{text}
\newcommand{\hlc}[2][yellow]{{\sethlcolor{#1}\hl{#2}}}

\newcommand{\dos}{\textcolor{red}}
%DIF 120d190
%DIF < \newcolumntype{L}[1]{>{\raggedright\let\newline\\\arraybackslash\hspace{0pt}}p{#1}}
%DIF -------

%DIF 122a191-194
%the following command implements roman numerals %DIF > 
\makeatletter %DIF > 
\newcommand*{\rom}[1]{\expandafter\@slowromancap\romannumeral #1@} %DIF > 
\makeatother %DIF > 
%DIF -------

%DIF < %%%%%%%%%%%%%%%%%%%%%%%%%%%%%%%%%%%%%%%%%%%%%%%%%%%%%%%%%%%%%%%%%%%%%%%%%
%DIF -------
 %DIF > 
%the following command reworks how \eqref{} works, by generating output like Eq. 1., instead of just (1). %DIF > 
\renewcommand{\eqref}[1]{Eq.~\textup{\ref{#1}}} %DIF > 
 %DIF > 
%this command should increase spacing between rows in tables %DIF > 
\renewcommand{\arraystretch}{1.6} %DIF > 
 %DIF > 
\newcommand{\degree}{\ensuremath{^\circ}} % Manually define \degree %DIF > 
\newcommand{\micro}{\textmu} %DIF > 
 %DIF > 
\let\qty\relax %DIF > 
\newcommand{\qty}[1]{\left( #1 \right)} %DIF > 
\newcommand{\qc}{\quad} %DIF > 
 %DIF > 
% the following adds a new column type identical to p{}, but left aligned %DIF > 
\newcolumntype{L}[1]{>{\raggedright\arraybackslash}p{#1}}  %DIF > 
%DIF PREAMBLE EXTENSION ADDED BY LATEXDIFF
%DIF UNDERLINE PREAMBLE %DIF PREAMBLE
\RequirePackage[normalem]{ulem} %DIF PREAMBLE
\RequirePackage{color}\definecolor{RED}{rgb}{1,0,0}\definecolor{BLUE}{rgb}{0,0,1} %DIF PREAMBLE
\providecommand{\DIFaddtex}[1]{{\protect\color{blue}\uwave{#1}}} %DIF PREAMBLE
\providecommand{\DIFdeltex}[1]{{\protect\color{red}\sout{#1}}} %DIF PREAMBLE
%DIF SAFE PREAMBLE %DIF PREAMBLE
\providecommand{\DIFaddbegin}{} %DIF PREAMBLE
\providecommand{\DIFaddend}{} %DIF PREAMBLE
\providecommand{\DIFdelbegin}{} %DIF PREAMBLE
\providecommand{\DIFdelend}{} %DIF PREAMBLE
\providecommand{\DIFmodbegin}{} %DIF PREAMBLE
\providecommand{\DIFmodend}{} %DIF PREAMBLE
%DIF FLOATSAFE PREAMBLE %DIF PREAMBLE
\providecommand{\DIFaddFL}[1]{\DIFadd{#1}} %DIF PREAMBLE
\providecommand{\DIFdelFL}[1]{\DIFdel{#1}} %DIF PREAMBLE
\providecommand{\DIFaddbeginFL}{} %DIF PREAMBLE
\providecommand{\DIFaddendFL}{} %DIF PREAMBLE
\providecommand{\DIFdelbeginFL}{} %DIF PREAMBLE
\providecommand{\DIFdelendFL}{} %DIF PREAMBLE
%DIF HYPERREF PREAMBLE %DIF PREAMBLE
\providecommand{\DIFadd}[1]{\texorpdfstring{\DIFaddtex{#1}}{#1}} %DIF PREAMBLE
\providecommand{\DIFdel}[1]{\texorpdfstring{\DIFdeltex{#1}}{}} %DIF PREAMBLE
\providecommand{\DIFscaledelfig}{0.5}
%DIF HIGHLIGHTGRAPHICS PREAMBLE %DIF PREAMBLE
\RequirePackage{settobox} %DIF PREAMBLE
\RequirePackage{letltxmacro} %DIF PREAMBLE
\newsavebox{\DIFdelgraphicsbox} %DIF PREAMBLE
\newlength{\DIFdelgraphicswidth} %DIF PREAMBLE
\newlength{\DIFdelgraphicsheight} %DIF PREAMBLE
% store original definition of \includegraphics %DIF PREAMBLE
\LetLtxMacro{\DIFOincludegraphics}{\includegraphics} %DIF PREAMBLE
\providecommand{\DIFaddincludegraphics}[2][]{{\color{blue}\fbox{\DIFOincludegraphics[#1]{#2}}}} %DIF PREAMBLE
\providecommand{\DIFdelincludegraphics}[2][]{% %DIF PREAMBLE
\sbox{\DIFdelgraphicsbox}{\DIFOincludegraphics[#1]{#2}}% %DIF PREAMBLE
\settoboxwidth{\DIFdelgraphicswidth}{\DIFdelgraphicsbox} %DIF PREAMBLE
\settoboxtotalheight{\DIFdelgraphicsheight}{\DIFdelgraphicsbox} %DIF PREAMBLE
\scalebox{\DIFscaledelfig}{% %DIF PREAMBLE
\parbox[b]{\DIFdelgraphicswidth}{\usebox{\DIFdelgraphicsbox}\\[-\baselineskip] \rule{\DIFdelgraphicswidth}{0em}}\llap{\resizebox{\DIFdelgraphicswidth}{\DIFdelgraphicsheight}{% %DIF PREAMBLE
\setlength{\unitlength}{\DIFdelgraphicswidth}% %DIF PREAMBLE
\begin{picture}(1,1)% %DIF PREAMBLE
\thicklines\linethickness{2pt} %DIF PREAMBLE
{\color[rgb]{1,0,0}\put(0,0){\framebox(1,1){}}}% %DIF PREAMBLE
{\color[rgb]{1,0,0}\put(0,0){\line( 1,1){1}}}% %DIF PREAMBLE
{\color[rgb]{1,0,0}\put(0,1){\line(1,-1){1}}}% %DIF PREAMBLE
\end{picture}% %DIF PREAMBLE
}\hspace*{3pt}}} %DIF PREAMBLE
} %DIF PREAMBLE
\LetLtxMacro{\DIFOaddbegin}{\DIFaddbegin} %DIF PREAMBLE
\LetLtxMacro{\DIFOaddend}{\DIFaddend} %DIF PREAMBLE
\LetLtxMacro{\DIFOdelbegin}{\DIFdelbegin} %DIF PREAMBLE
\LetLtxMacro{\DIFOdelend}{\DIFdelend} %DIF PREAMBLE
\DeclareRobustCommand{\DIFaddbegin}{\DIFOaddbegin \let\includegraphics\DIFaddincludegraphics} %DIF PREAMBLE
\DeclareRobustCommand{\DIFaddend}{\DIFOaddend \let\includegraphics\DIFOincludegraphics} %DIF PREAMBLE
\DeclareRobustCommand{\DIFdelbegin}{\DIFOdelbegin \let\includegraphics\DIFdelincludegraphics} %DIF PREAMBLE
\DeclareRobustCommand{\DIFdelend}{\DIFOaddend \let\includegraphics\DIFOincludegraphics} %DIF PREAMBLE
\LetLtxMacro{\DIFOaddbeginFL}{\DIFaddbeginFL} %DIF PREAMBLE
\LetLtxMacro{\DIFOaddendFL}{\DIFaddendFL} %DIF PREAMBLE
\LetLtxMacro{\DIFOdelbeginFL}{\DIFdelbeginFL} %DIF PREAMBLE
\LetLtxMacro{\DIFOdelendFL}{\DIFdelendFL} %DIF PREAMBLE
\DeclareRobustCommand{\DIFaddbeginFL}{\DIFOaddbeginFL \let\includegraphics\DIFaddincludegraphics} %DIF PREAMBLE
\DeclareRobustCommand{\DIFaddendFL}{\DIFOaddendFL \let\includegraphics\DIFOincludegraphics} %DIF PREAMBLE
\DeclareRobustCommand{\DIFdelbeginFL}{\DIFOdelbeginFL \let\includegraphics\DIFdelincludegraphics} %DIF PREAMBLE
\DeclareRobustCommand{\DIFdelendFL}{\DIFOaddendFL \let\includegraphics\DIFOincludegraphics} %DIF PREAMBLE
%DIF AMSMATHULEM PREAMBLE %DIF PREAMBLE
\makeatletter %DIF PREAMBLE
\let\sout@orig\sout %DIF PREAMBLE
\renewcommand{\sout}[1]{\ifmmode\text{\sout@orig{\ensuremath{#1}}}\else\sout@orig{#1}\fi} %DIF PREAMBLE
\makeatother %DIF PREAMBLE
%DIF COLORLISTINGS PREAMBLE %DIF PREAMBLE
\RequirePackage{listings} %DIF PREAMBLE
\RequirePackage{color} %DIF PREAMBLE
\lstdefinelanguage{DIFcode}{ %DIF PREAMBLE
%DIF DIFCODE_UNDERLINE %DIF PREAMBLE
  moredelim=[il][\color{red}\sout]{\%DIF\ <\ }, %DIF PREAMBLE
  moredelim=[il][\color{blue}\uwave]{\%DIF\ >\ } %DIF PREAMBLE
} %DIF PREAMBLE
\lstdefinestyle{DIFverbatimstyle}{ %DIF PREAMBLE
	language=DIFcode, %DIF PREAMBLE
	basicstyle=\ttfamily, %DIF PREAMBLE
	columns=fullflexible, %DIF PREAMBLE
	keepspaces=true %DIF PREAMBLE
} %DIF PREAMBLE
\lstnewenvironment{DIFverbatim}[1][]{\lstset{style=DIFverbatimstyle}}{} %DIF PREAMBLE
\lstnewenvironment{DIFverbatim*}[1][]{\lstset{style=DIFverbatimstyle,showspaces=true}}{} %DIF PREAMBLE
\lstset{extendedchars=\true,inputencoding=utf8}

%DIF END PREAMBLE EXTENSION ADDED BY LATEXDIFF

\begin{document}
\DIFdelbegin %DIFDELCMD < 

%DIFDELCMD < %%%
%DIF <  Title Page
\DIFdelend \begin{titlepage}
%DIF < \centering
    \DIFdelbegin %DIFDELCMD < \vspace*{1cm}
%DIFDELCMD <     %%%
\DIFdelend \DIFaddbegin \sffamily \large
\DIFaddend 

\DIFdelbegin %DIFDELCMD < \includegraphics[width=4cm,right]{Graphics/Logos/RexusBexus.jpg}
%DIFDELCMD <     

%DIFDELCMD <     %%%
\DIFdelend \DIFaddbegin \begin{tikzpicture}[remember picture,overlay]
\node[anchor=north east, xshift=-1.5cm, yshift=-1.2cm]
  \DIFadd{at (current page.north east) }\DIFaddend {\DIFdelbegin %DIFDELCMD < \Huge\bfseries %%%
\DIFdel{SED}\DIFdelend %DIF > 
    \DIFaddbegin \begin{minipage}{4cm}
      \raggedleft
      \includegraphics[width=4cm]{0_images/titlepage/REXUS_BEXUS_logo.jpg}\DIFaddend \par\DIFaddbegin \vspace{6mm}
      \includegraphics[width=4cm]{0_images/titlepage/YourTeamHere.PNG}
    \end{minipage}
\DIFaddend }\DIFaddbegin \DIFadd{;
}\end{tikzpicture}

\vspace{1cm}

\begin{minipage}[b]{11cm}
  \DIFaddend {\DIFdelbegin %DIFDELCMD < \Large %%%
\DIFdel{Student Experiment Documentation}%DIFDELCMD < \par%%%
\DIFdelend \DIFaddbegin \Huge \textbf{\DIFadd{SED}}\DIFaddend }

  \DIFaddbegin {\large \textbf{\DIFadd{Student Experiment Description}}}
  \DIFaddend \vspace{1cm}

  Document ID: \DIFdelbegin \DIFdel{BX3839\_MIRAGE\_SED\_v1-0\_14Jan2026
    }\DIFdelend \DIFaddbegin \DocumentID
  \vspace{1cm}
\DIFaddend 

  \DIFdelbegin %DIFDELCMD < \vspace{0.5cm}
%DIFDELCMD <     

%DIFDELCMD <     %%%
\DIFdelend {\large \DIFdelbegin %DIFDELCMD < \bfseries %%%
\DIFdel{Mission: BEXUS}%DIFDELCMD < \par%%%
\DIFdelend \DIFaddbegin \textbf{\DIFadd{Mission: BEXUS}}\DIFaddend }
  \DIFaddbegin \vspace{1cm}
\DIFaddend 

  {\Large \DIFdelbegin %DIFDELCMD < \bfseries %%%
\DIFdel{Team Name: MIRAGE}%DIFDELCMD < \par%%%
\DIFdelend \DIFaddbegin \textbf{\DIFadd{Team Name: MIRAGE}}\DIFaddend }
  \DIFaddbegin \vspace{1cm}
\DIFaddend 

  \DIFdelbegin %DIFDELCMD < \vspace{0.5cm}
%DIFDELCMD < 

%DIFDELCMD <     %%%
\DIFdelend {\large Experiment Title: Methane Infrared Absorption Gas Experiment\DIFdelbegin %DIFDELCMD < \par%%%
\DIFdelend }
\DIFaddbegin \end{minipage}
\DIFaddend 

\DIFdelbegin %DIFDELCMD < \vspace{1cm}
%DIFDELCMD <     %%%
\DIFdelend \DIFaddbegin \vspace{0.5cm}
\DIFaddend 

\DIFdelbegin %DIFDELCMD < \begin{tabular}{lll}
%DIFDELCMD <         \textbf{Team} & \textbf{Name} & \textbf{University} \\
%DIFDELCMD <         \hline
%DIFDELCMD <         Student Team Leader: & Abdullah Alsakka & Uppsala University \\
%DIFDELCMD <         Team Members: & Jonathan Remmert & Uppsala University \\
%DIFDELCMD <         & Anna Tschesche & Uppsala University \\
%DIFDELCMD <         & Oliver Bolin & Uppsala University\\
%DIFDELCMD <         & Max Steen & Uppsala University \\
%DIFDELCMD <         & Olle Hallberg & Uppsala University \\
%DIFDELCMD <         & Molly Sundell & Uppsala University \\
%DIFDELCMD <         & Leo Adolphus Wattin & Uppsala University \\
%DIFDELCMD <         & Antoni Marciniewski & Uppsala University \\
%DIFDELCMD <     \end{tabular}
%DIFDELCMD <     %%%
\DIFdelend \DIFaddbegin \begin{tabular}{l l l}
  \textbf{Team}        & \textbf{Name} & \textbf{University}\\
  Student Team Leader: & Abdullah Alsakka & Uppsala University \\
  Team Members: & Jonathan Remmert & Uppsala University \\
    & Anna Tschesche & Uppsala University \\
    & Oliver Bolin & Uppsala University\\
    & Max Steen & Uppsala University \\
    & Olle Hallberg & Uppsala University \\
    & Molly Sundell & Uppsala University \\
    & Leo Adolphus Wattin & Uppsala University \\
    & Antoni Marciniewski & Uppsala University \\   
\end{tabular}
\DIFaddend 

\DIFdelbegin %DIFDELCMD < \vspace{2cm}
%DIFDELCMD <     %%%
\DIFdelend \DIFaddbegin \enlargethispage{3cm}
\DIFaddend 

\DIFdelbegin %DIFDELCMD < \begin{tabular}{llll}
%DIFDELCMD <         Version: & Issue Date: & Document Type: & Valid from: \\
%DIFDELCMD <         \textbf{1.0} & \textbf{14 January 2026} & \textbf{SED} & \textbf{14 October 2024} 
%DIFDELCMD <     \end{tabular}
%DIFDELCMD <     %%%
\DIFdelend \DIFaddbegin \begin{tabular}{llll}
  Version:        & Issue Date:               & Document Type: & Valid from: \\
  \textbf{1.2}      & \textbf{15 March 2026}    & \textbf{SED}  & \textbf{17 December 2025}
\end{tabular}
\DIFaddend 

\DIFdelbegin %DIFDELCMD < \vfill
%DIFDELCMD <     \vspace{1cm}
%DIFDELCMD <     

%DIFDELCMD <     %%%
\DIFdelend Issued by:
\DIFdelbegin \DIFdel{Jonathan Remmert, Anna Tschesche, Oliver Bolin, Max Steen, Olle Hallberg, Molly Sundell, Leo Adolphus Wattin, Antoni Marciniewski, Abdullah Alsakka }%DIFDELCMD < \\    %%%
%DIF < \textbf{Name of the person/s who issued the document}
    \DIFdelend \DIFaddbegin \vspace{0.3cm}
\DIFaddend 

\DIFdelbegin %DIFDELCMD < \vspace{1cm}
%DIFDELCMD <     %%%
\DIFdelend \DIFaddbegin {\large \textbf{\DIFadd{Jonathan Remmert,
Anna Tschesche,
Oliver Bolin,
Max Steen,
Olle Hallberg,
Molly Sundell,
Leo Adolphus Wattin,
Antoni Marciniewski}}} %DIF > change this
\DIFaddend 

Approved by:
\DIFdelbegin \DIFdel{Assoc. Prof. Uwe Zimmermann}%DIFDELCMD < \\    %%%
%DIF < \textbf{Name of the person/s who approved the document}\par
\DIFdelend \DIFaddbegin \vspace{0.3cm}
\DIFaddend 

\DIFdelbegin %DIFDELCMD < \vspace{1cm}    
%DIFDELCMD < %%%
\DIFdelend %DIF > Please replace this with a team member(s)
\DIFaddbegin {\large \textbf{\DIFadd{Assoc. Prof. Uwe Zimmermann}}} 
\DIFaddend 

%DIF > Please replace this with a reviewer/approver such as a professor.
%DIF >  {\large \textit{(The person who approves the document must be different from the one who issues)}}
\DIFaddbegin 


\DIFaddend \end{titlepage}

%DIF < %%%%%%%%%%%%%%%%%%%%%%%%%%%%%%%%%%%%%%%%%%%%%%%%%%%%%%%%%%%%%%%%%%%%%%%%
%DIF < %%%%%%%%%%%%%%%%%%%%%%%%%% Change Record %%%%%%%%%%%%%%%%%%%%%%%%%%%%%%%
\DIFdelbegin %DIFDELCMD < \chapter*{Change Record}
%DIFDELCMD < \markboth{Change Record}{Change Record}
%DIFDELCMD < %%%
\addcontentsline{toc}{chapter}{\DIFdel{Change Record}}
%DIFAUXCMD
\DIFdelend %DIF > Document adapted from Team GaGa and Project SunbYte
%DIF > Written by MARA



\DIFdelbegin %DIFDELCMD < \begin{table}[h]
%DIFDELCMD < \centering
%DIFDELCMD < \begin{tabular}{ccp{5cm}p{5cm}}
%DIFDELCMD < \toprule
%DIFDELCMD < \textbf{Version} & \textbf{Date} & \textbf{Changed chapters} & \textbf{Remarks} \\
%DIFDELCMD < \midrule
%DIFDELCMD < 0 & 2024-Oct-14 & New Version & Blank Book 2024 \\
%DIFDELCMD < \rowcolor{lightyellow}
%DIFDELCMD < 1 & 2026-Jan-19 & all & PDR \\
%DIFDELCMD < 2 & & & CDR \\
%DIFDELCMD < 3 & & & IPR \\
%DIFDELCMD < 4 & & & EAR, Pre-Campaign \\
%DIFDELCMD < 5 & & & Final report \\
%DIFDELCMD < \bottomrule
%DIFDELCMD < \end{tabular}
%DIFDELCMD < \end{table}
%DIFDELCMD < %%%
\DIFdelend %DIF > load all acronyms
\DIFaddbegin \newacronym{ait}{AIT}{Assembly, Integration and Test}
\newacronym{asap}{ASAP}{As Soon As Possible}
\newacronym{cdr}{CDR}{Critical Design Review}
\newacronym{cog}{COG}{Centre of Gravity}
\newacronym{cots}{COTS}{Commercial off-the-shelf}
\newacronym{crc}{CRC}{Cyclic Redundancy Check}
\newacronym{crp}{CRP}{Campaign Requirement Plan}
\newacronym{ctm}{CTM}{chemistry and transport models}
\newacronym{dlr}{DLR}{Deutsches Zentrum für Luft- und Raumfahrt}
\newacronym{eat}{EAT}{Experiment Acceptance Test}
\newacronym{ear}{EAR}{Experiment Acceptance Review}
\newacronym{ects}{ECTS}{European Credit Transfer System}
\newacronym{eit}{EIT}{Electrical Interface Test}
\newacronym{epm}{EPM}{Esrange Project Manager}
\newacronym{esa}{ESA}{European Space Agency}
\newacronym{esrange}{Esrange}{Esrange Space Center}
\newacronym{estec}{ESTEC}{European Space Research and Technology Centre, ESA (NL)}
\newacronym{esw}{ESW}{Experiment Selection Workshop}
\newacronym{far}{FAR}{Flight Acceptance Review}
\newacronym{fst}{FST}{Flight Simulation Test}
\newacronym{frp}{FRP}{Flight Requirement Plan}
\newacronym{frr}{FRR}{Flight Readiness Review}
\newacronym{ftir}{FTIR}{Fourier-Transform Infrared Spectroscopy}
\newacronym{gse}{GSE}{Ground Support Equipment}
\newacronym{gui}{GUI}{Graphical User Interface}
\newacronym{hk}{HK}{House Keeping}
\newacronym{hw}{H/W}{Hardware}
\newacronym{icd}{ICD}{Interface Control Document}
\newacronym{if}{I/F}{Interface}
\newacronym{ipr}{IPR}{Integration Progress Review}
\newacronym{ir}{IR}{Infrared}
\newacronym{irf}{IRF}{Institutet För Rymdfysik}
\newacronym{lo}{LO}{Lift Off}
\newacronym{lod}{LoD}{Limit of Detection}
\newacronym{lt}{LT}{Local Time}
\newacronym{los}{LOS}{Line of Sight}
\newacronym{mbps}{Mbps}{Mega Bits per Second}
\newacronym{mcu}{MCU}{Microcontroller Unit}
\newacronym{mfh}{MFH}{Mission Flight Handbook}
\newacronym{mirage}{MIRAGE}{Methane InfraRed Absorption Gas Experiment}
\newacronym{moraba}{MORABA}{Mobile Raketen Basis (DLR, EuroLaunch)}
\newacronym{mosfet}{MOSFET}{Metal-Oxide-Semiconductor Field-Effect Transistor}
\newacronym{ndir}{NDIR}{Non-Dispersive Infrared}
\newacronym{op}{OP}{Oberpfaffenhofen, DLR Center}
\newacronym{pcb}{PCB}{Printed Circuit Board}
\newacronym{pdr}{PDR}{Preliminary Design Review}
\newacronym{pid}{PID}{Proportional-Derivative-Integral}
\newacronym{psa}{PSA}{Pressure Sensitive Adhesive}
\newacronym{pst}{PST}{Payload System Test}
\newacronym{pwm}{PWM}{Pulse Width Modulation}
\newacronym{ptc}{PTC}{Positive Temperature Coefficient}
\newacronym{rbf}{RBF}{Remove Before Flight}
\newacronym{rtc}{RTC}{Real Time Clock}
\newacronym{rtm}{RTM}{radiative transfer models}
\newacronym{sed}{SED}{Student Experiment Documentation}
\newacronym{snsb}{SNSB}{Swedish National Space Board}
\newacronym{sods}{SODS}{Start Of Data Storage}
\newacronym{soe}{SOE}{Start Of Experiment}
\newacronym{stw}{STW}{Student Training Week}
\newacronym{sw}{S/W}{Software}
\newacronym{t}{T}{Time Before and After Launch Noted with + or -}
\newacronym{tbc}{TBC}{To Be Confirmed}
\newacronym{tbd}{TBD}{To Be Determined}
\newacronym{tco}{TCO}{Thermal CutOut}
\newacronym{tdlas}{TDLAS}{Tunable Diode Laser Absorption Spectroscopy}
\newacronym{wbs}{WBS}{Work Breakdown Structure}
\newacronym{zarm}{ZARM}{Zentrum für Angewandte Raumfahrttechnologie und Mikrogravitation}
\DIFaddend 


%DIF > title and changerecord
\DIFaddbegin \thispagestyle{empty}
\newpage

 \addchap{Change Record}

 \begin{tabular}{|l|l|l|l|}
    \hline
    \textbf{Version}        &   \textbf{Date}   & \textbf{Changed chapters} & \textbf{Remarks}  \\\midrule
    \hline
    0       &   2025-Dec-17      &   New Version     &   Blank Book 2025 \\\hline
    1-0       & 2026-Jan-19      &   all             &   PDR \\\hline
   \rowcolor{yellow} 1-1       & 2026-Mar-15       &   \dos{all}             &  Revision after PDR \\\hline
    2-0       &                   &                   &   CDR \\\hline
    3-0       &                   &                   &   IPR \\\hline
    4-0       &                   &                   &   EAR, Pre-Campaign \\\hline
    5-0       &                   &                   &   Final report \\ \hline
 \end{tabular}

\DIFaddend \textbf{Keywords:} \DIFaddbegin \DIFadd{REXUS or }\DIFaddend BEXUS, SED - Student Experiment Documentation\DIFdelbegin \DIFdel{, Methane, NDIR
}\DIFdelend \DIFaddbegin \DIFadd{...
}\DIFaddend 

%DIF < %%%%%%%%%%%%%%%%%%%%%%%%%%%%%%%%%%%%%%%%%%%%%%%%%%%%%%%%%%%%%%%%%%%%%%%%
%DIF <  Table of Contents
\DIFaddbegin \textit{\DIFadd{Please keep the change record as detailed as possible. It helps the
organisers to review your SED in shorter time as well as you to keep track of
your changes throughout the programme.
When updating chapters, instead of saying chapter 1, state sections chapter
1.1. and so on for technical changes (design, numbers, status…). For
linguistic changes you might write “minor changes in chapter 1/ throughout the
document” etc… }} \\
\textbf{\DIFadd{We require you to additionally highlight changes you made from the last
issue in a light colour (e.g. yellow) in addition to the regular version.}}

%DIF > %%%%%%%%%%%%%%%%%%%%%%%%%%%%%%%%%%%%%%%%%%%%%%%%%%%%%%%%%%%%%%%%%%%%%%%%%%%%%%

%DIF > table of contents
\DIFaddend \tableofcontents
\DIFaddbegin \newpage
\DIFaddend 

%DIF < %%%%%%%%%%%%%%%%%%%%%%%%%%%%%%%%%%%%%%%%%%%%%%%%%%%%%%%%%%%%%%%%%%%%%%%%
%DIF <  List of Figures
\listoffigures
\DIFaddbegin \newpage
\DIFaddend 

%DIF < %%%%%%%%%%%%%%%%%%%%%%%%%%%%%%%%%%%%%%%%%%%%%%%%%%%%%%%%%%%%%%%%%%%%%%%%
%DIF <  List of Tables
\listoftables
\DIFaddbegin \newpage
\DIFaddend 

%DIF < %%%%%%%%%%%%%%%%%%%%%%%%%%%%%%%%%%%%%%%%%%%%%%%%%%%%%%%%%%%%%%%%%%%%%%%%
%DIF < %%%%%%%%%%%%%%%%%%%%%%%% Preface %%%%%%%%%%%%%%%%%%%%%%%%%%%%%%%%%%%%%%%
\DIFdelbegin %DIFDELCMD < \chapter*{Preface}
%DIFDELCMD < \markboth{Preface}{Preface}
%DIFDELCMD < %%%
\addcontentsline{toc}{chapter}{\DIFdel{Preface}}
%DIFAUXCMD
\DIFdelend \DIFaddbegin \cleardoublepage
\newpage
\DIFaddend 

\DIFaddbegin \cleardoublepage
\newpage

\cleardoublepage
%DIF > %%%%%%%%%%%%%%%%%%%%%%%%%%%%%%%%%%%%%%%%%%%%%%%%%%%%%%%%%%%%%%%%%%%%%%%%%%%%%%

%DIF > main text
\addchap{PREFACE}
\DIFaddend The REXUS/BEXUS programme is realised under a bilateral Agency Agreement between the German Aerospace Center (DLR) and the Swedish National Space Agency (SNSA). The Swedish share of the payload has been made available to students from other European countries through a collaboration with the European Space Agency (ESA). EuroLaunch, a cooperation between the Swedish Space Corporation (SSC) and the Mobile Rocket Base (MORABA) of DLR, is responsible for the campaign management and operations of the launch vehicles. Experts from DLR, SSC, ZARM and ESA provide technical support to the student teams throughout the project. REXUS and BEXUS are launched from SSC, Esrange Space Center in northern Sweden.

\DIFdelbegin \DIFdel{\mbox{%DIFAUXCMD
\ac{mirage} }\hskip0pt%DIFAUXCMD
}\DIFdelend \DIFaddbegin \DIFadd{\mbox{%DIFAUXCMD
\gls{mirage} }\hskip0pt%DIFAUXCMD
}\DIFaddend is an experiment developed by a team of students from the Uppsala Universitet in Sweden. Its aim is to demonstrate low-cost methane sensing capabilities using stratospheric balloons. \DIFdelbegin \DIFdel{\mbox{%DIFAUXCMD
\ac{mirage} }\hskip0pt%DIFAUXCMD
}\DIFdelend \DIFaddbegin \DIFadd{\mbox{%DIFAUXCMD
\gls{mirage} }\hskip0pt%DIFAUXCMD
}\DIFaddend is scheduled to fly aboard BEXUS as part of cycle 17 of the Swedish-German programme REXUS/BEXUS.

This document contains all relevant available information regarding \DIFdelbegin \DIFdel{\mbox{%DIFAUXCMD
\ac{mirage}}\hskip0pt%DIFAUXCMD
}\DIFdelend \DIFaddbegin \DIFadd{\mbox{%DIFAUXCMD
\gls{mirage}}\hskip0pt%DIFAUXCMD
}\DIFaddend ; including the experiment's objectives and requirements, a breakdown of the experimental module's design, design verification procedures, a timeline and pre-launch campaign activities. Technical documents that pertain to the design are also included for reference, as appendix sections.\DIFaddbegin \newpage
\DIFaddend 

%DIF < %%%%%%%%%%%%%%%%%%%%%%%%%%%%%%%%%%%%%%%%%%%%%%%%%%%%%%%%%%%%%%%%%%%%%%%%
%DIF < %%%%%%%%%%%%%%%%%%%%%%%% Abstract %%%%%%%%%%%%%%%%%%%%%%%%%%%%%%%%%%%%%%
\DIFdelbegin %DIFDELCMD < \chapter*{Abstract}
%DIFDELCMD < \markboth{Abstract}{Abstract}
%DIFDELCMD < %%%
\addcontentsline{toc}{chapter}{\DIFdel{Abstract}}
%DIFAUXCMD
\DIFdelend \DIFaddbegin \addchap{ABSTRACT}
\DIFaddend 

Methane \DIFaddbegin \DIFadd{is one of the most powerfull greenhouse gases contributing to climate change with about 25 times as much warming potential as carbondioxide over the span of one hundred years. To reduce the methane content in the atmosphere, it is highly important to understand methane sources and sinks. While a lot of methane is emitted by industries and waste sites, there are also natural emissions from, for example, wetlands. This shows the diverse environments in which monitoring is necessary. For this reason, current-day large-scale monitoring is done from space, relying on atmospheric models to calculate the methane concentration profile from the averaged column density. These atmospheric models require in-situ measurement data to make them more accurate. However, not many missions have done so for methane in the stratosphere, which makes the models less reliable overall and therefore renders monitoring less reliable as well. }\\
\DIFadd{Furthermore, commonly used laser-based technologies measuring methane with high precision are often very expensive and complex. This prevents more in-situ measurements from being recorded in the stratosphere.}\\  
\DIFadd{Methane }\DIFaddend InfraRed Absorption Gas Experiment (MIRAGE) aims to \DIFaddbegin \DIFadd{demonstrate a low-cost, low-complexity alternative to laser-based technologies to }\DIFaddend measure stratospheric methane concentrations\DIFaddbegin \DIFadd{, }\DIFaddend using an off-the-shelf Non-Dispersive InfraRed (NDIR) sensor \DIFdelbegin \DIFdel{. Current measurement techniques rely on laser-based setups, requiring fine adjustment while having a high cost. While these systems provide high resolution, they prevent frequent and cost-effective measurements. MIRAGE will demonstrate that NDIR sensors can fill the niche of low-cost measurements with reasonable resolution.
}%DIFDELCMD < 

%DIFDELCMD < %%%
\DIFdel{This is done by pressurizing }\DIFdelend \DIFaddbegin \DIFadd{core. To increase the accuracy, }\DIFaddend temperature-controlled ambient air \DIFaddbegin \DIFadd{is pressurised }\DIFaddend before feeding it into the optical chamber of the NDIR sensor. With this \DIFdelbegin \DIFdel{setup it is expected to achieve an absolute concentration accuracy of $\mathrm{\Delta}c \leq 0.1$ ppm.
}\DIFdelend \DIFaddbegin \DIFadd{experiment, NDIR technology will be tested in regards to its limitations and future potential for use in atmospheric methane monitoring.
%DIF > This is done by pressurizing temperature-controlled ambient air before feeding it into the optical chamber of the NDIR sensor.
%DIF > With this setup it is expected to achieve an absolute concentration accuracy of \dos{$\mathrm{\Delta}c \leq 0.1$ ppm}.
}\DIFaddend 


%DIF < %%%%%%%%%%%%%%%%%%%%%%%%%%%%%%%%%%%%%%%%%%%%%%%%%%%%%%%%%%%%%%%%%%%%%%%%
%DIF < %%%%%%%%%%%%%%%%%%%%%%%% Content %%%%%%%%%%%%%%%%%%%%%%%%%%%%%%%%%%%%%%%
%DIF > \dos{EXPAND THIS. Expansion is not a required action, but it was commented that our abstract is very short.}\newpage

\chapter{\DIFdelbegin \DIFdel{Introduction}\DIFdelend \DIFaddbegin \DIFadd{INTRODUCTION}\DIFaddend }
\DIFdelbegin %DIFDELCMD < 

%DIFDELCMD < %%%
\DIFdelend \DIFaddbegin \label{sec:introduction}
\DIFaddend \section{Scientific/Technical Background}

Gases in the stratosphere play a pivotal role in climate modelling due to their prolonged lifetime.
Especially methane sources have been shown to play a big role in climate modelling \cite{stocker2013, saunois2017, nisbet2019, turner2019}\DIFdelbegin \DIFdel{, and the need of accurate measurements }\DIFdelend \DIFaddbegin \DIFadd{. Since remote sensing via satellites is used to monitor methane sources and sinks on a large scale they contribute a lot to climate models. However, they heavily depend on atmospheric models to compute the methane concentrations at different altitudes from the averaged column densities they can measure. Therefore, more accurate in-situ measurements of methane are needed, which }\DIFaddend was reflected on in the National Research Council Decadal Survey for Earth Science \cite{national2019} and the summary report by the Carbon Climate Workshop in 2018 \cite{sellers2018}. 

The concentration of atmospheric methane ranges from around 2 ppm \DIFaddbegin \DIFadd{down }\DIFaddend to 0.7 ppm \DIFdelbegin \DIFdel{at }\DIFdelend \DIFaddbegin \DIFadd{between }\DIFaddend 0 \DIFdelbegin \DIFdel{km }\DIFdelend and 30 km altitude\DIFdelbegin \DIFdel{, respectively}\DIFdelend .
Because of the small concentrations, accurate optically complex systems like \DIFdelbegin \DIFdel{\mbox{%DIFAUXCMD
\ac{tdlas} }\hskip0pt%DIFAUXCMD
\mbox{%DIFAUXCMD
\cite{joly2020} }\hskip0pt%DIFAUXCMD
or \mbox{%DIFAUXCMD
\ac{ftir} }\hskip0pt%DIFAUXCMD
}\DIFdelend \DIFaddbegin \DIFadd{\mbox{%DIFAUXCMD
\gls{tdlas} }\hskip0pt%DIFAUXCMD
\mbox{%DIFAUXCMD
\cite{joly2020} }\hskip0pt%DIFAUXCMD
or \mbox{%DIFAUXCMD
\gls{ftir} }\hskip0pt%DIFAUXCMD
}\DIFaddend \cite{runge2023} are widely employed.
\DIFdelbegin \DIFdel{\mbox{%DIFAUXCMD
\ac{tdlas} }\hskip0pt%DIFAUXCMD
}\DIFdelend \DIFaddbegin \DIFadd{\mbox{%DIFAUXCMD
\gls{tdlas} }\hskip0pt%DIFAUXCMD
}\DIFaddend systems send a laser beam of varying frequency close to the excitation energy of the target molecule through a gas sample and calculate the target molecule's concentration based on the measured absorption via the Beer-Lambert law.
\DIFdelbegin \DIFdel{\mbox{%DIFAUXCMD
\ac{ftir} }\hskip0pt%DIFAUXCMD
}\DIFdelend \DIFaddbegin \DIFadd{\mbox{%DIFAUXCMD
\gls{ftir} }\hskip0pt%DIFAUXCMD
}\DIFaddend works by sending a light beam containing multiple frequencies into the gas sample and recording the total absorption.
The frequencies are then changed and the process repeated.
This is often done via a Michelsen interferometer.
Both techniques \DIFaddbegin \DIFadd{have the potential to }\DIFaddend enable great accuracy \DIFdelbegin \DIFdel{but require high precision }\DIFdelend \DIFaddbegin \DIFadd{depending on their path length. However, to get longer pathlengths it is often necessary to eflect the laser beam thousands of times since measurement are usually limited in space. Therefore, high precision is required }\DIFaddend in the complex optical setups, meaning that those systems are expensive\DIFdelbegin \DIFdel{and inaccessible }\DIFdelend \DIFaddbegin \DIFadd{, inaccessible and often highly sensitive to vibrations and shocks}\DIFaddend .

During recent years, \DIFdelbegin \DIFdel{\mbox{%DIFAUXCMD
\ac{ndir} }\hskip0pt%DIFAUXCMD
}\DIFdelend \DIFaddbegin \DIFadd{\mbox{%DIFAUXCMD
\gls{ndir} }\hskip0pt%DIFAUXCMD
}\DIFaddend sensors have advanced as a cost-effective and simple alternative to laser-based devices commonly employed for methane sensing.
NDIR spectroscopy works by passing a spectrum of \DIFdelbegin \DIFdel{\mbox{%DIFAUXCMD
\ac{ir} }\hskip0pt%DIFAUXCMD
}\DIFdelend \DIFaddbegin \DIFadd{\mbox{%DIFAUXCMD
\gls{ir} }\hskip0pt%DIFAUXCMD
}\DIFaddend light through a gas sample where each gas absorbs the light at a unique wavelength.
An optical filter selects the wavelengths of the gas of interest.
The intensity of the transmitted light at the relevant wavelength can then be measured \DIFdelbegin \DIFdel{, and thus the change in intensity is known}\DIFdelend \DIFaddbegin \DIFadd{to receive the absorbance $A=\log_{10}\left(\frac{I_0}{I}\right) = c\varepsilon l$}\DIFaddend .
The concentration c of the gas can then be calculated using the Beer-Lambert law

\begin{equation}
c = \log_{10}\left( \frac{I_{0}}{I} \right) \times {(\varepsilon l)}^{- 1},
\end{equation}

where $\varepsilon$ is the molar absorption coefficient, specific to the gas species, $l$ is the path length between light source and detector and $I_{0}$ and $I$ are the initial intensity of light and the intensity at the detector.
$I_{0}$ is determined either before the measurement or simultaneously using a second detector.
A schematic is provided in Figure \ref{fig:ndir_principle}.
Due to the simple setup, the overall cost of NDIR detectors is quite small with a range of tens to a few 100s of Euros. \DIFaddbegin \DIFadd{Furthermore, using a simple setup with a broad light sources, makes the device more robust against misalignment.
}\DIFaddend 

\DIFdelbegin %DIFDELCMD < \begin{figure}[h]
%DIFDELCMD < %%%
\DIFdelendFL \DIFaddbeginFL \begin{figure}[htbp]
\DIFaddendFL \centering
\DIFdelbeginFL %DIFDELCMD < \includegraphics[width=8cm]{Graphics/Figures/NDIR schematic.png}
%DIFDELCMD < %%%
\DIFdelendFL \DIFaddbeginFL \includegraphics[width=8cm]{0_images/figures/NDIR schematic.png}
\DIFaddendFL \caption{Schematic of the non-dispersive infrared (NDIR) measurement principle.
Source: \cite{gaynullin2025}}
\label{fig:ndir_principle}
\end{figure}

\DIFdelbegin \DIFdel{Since the pressure }\DIFdelend \DIFaddbegin \DIFadd{At the same time, less reflections mean a shorter pathlength and a broader wavelength spectrum means limited applicability of the Beer-Lambert law. Another drawback of having a light source with a broader spectrum (compared to a laser) is that the absorption spectra of several molecules can overlap and make it harder to differentiate the effect of different absorbing molecules. In the case of methane, there is a large overlap with water molecules, which are usually present in a far bigger concentration in the atmosphere. This has previously limited NDIR so much that it couldn't be applied to measuring atmospheric methane. }\\
\begin{figure}[htbp]
\centering
\includegraphics[width=8cm]{0_images/figures/K96principle.png}
\caption{\DIFaddFL{Schematic of K96 sensor core, which combines three detectors for water, arbondioxide and methane.
Source: \mbox{%DIFAUXCMD
\cite{wastine2022}}\hskip0pt%DIFAUXCMD
}}
\label{fig:K96_principle}
\end{figure}
\DIFadd{However, recently a new sensor has been developed by Senseair AB \mbox{%DIFAUXCMD
\cite{wastine2022}}\hskip0pt%DIFAUXCMD
, which measures $\mathrm{CH_4}$, $\mathrm{H_2 O}$, and $\mathrm{CO_2}$ at the same time by using three different detectors in one sensor core as depicted in }\autoref{fig:K96_principle}\DIFadd{. This allows one to subtract the influence of water and carbondioxide absorption from the methane signal. Therefore, this NDIR sensor is a promising candidate for atmospheric methane monitoring. }\\  
\DIFadd{While the K96 sensor is a state of the art sensor and delivers higher resolution methane measurements than other NDIR sensors, it is not quite accurate enough to measure atmospheric methane concentrations reliably. The \mbox{%DIFAUXCMD
\gls{lod}}\hskip0pt%DIFAUXCMD
, which gives the lowest concentration at atmospheric pressures that can be detected by the K96 sensor is about 3ppm\mbox{%DIFAUXCMD
\cite{mao_improvement_2025}}\hskip0pt%DIFAUXCMD
. Therefore, it is necessary to increase the concentration. Here it is important to differentiate the terms concentration and molar fraction, which are often used interchangeably because most contexts are working under ground level conditions with a constant pressure. While concentration gives the particle number per volume $c=\frac{n}{V}$, the molar fraction means particle number of a certain gas species per total particle number including all present gas species $x_i=\frac{n_i}{n_{tot}}$. When a detector resolution is given in particles-per-million or ppm, this is usually assumed for ground level. In reality the quantity of interest for an optical sensor that is based on the Beer-Lambert law is the concentration of the gas species.}\\ 
\DIFadd{Now, since the pressure }\DIFaddend $p$ decreases exponentially with altitude, the concentration

\begin{equation}
c = \frac{n}{V}
\end{equation}

connected to the pressure by the ideal gas law

\begin{equation}
pV = nRT
\end{equation}

with the volume $V$, the ideal gas constant $R$ and the temperature $T$, decreases with about the same rate as the pressure. \DIFdelbegin \DIFdel{At the same time the }\DIFdelend \DIFaddbegin \DIFadd{If it were assumed that the composition of the air stayed constant, the molar fraction would not be affected by this drop in the particle number. However, the emitted light would now encounter way fewer particles and therefore be less absorbed. This means that a detector with low accuracy has an even lower signal if the pressure drops. Furthermore, the }\DIFaddend relative molar fraction of methane molecules compared to other molecules in the air is \DIFdelbegin \DIFdel{decreasing}\DIFdelend \DIFaddbegin \DIFadd{actually decreasing from 2\,ppm to about 1\,ppm as the atmospheric composition changes between Troposphere and Stratosphere}\DIFaddend , which means that there are \DIFdelbegin \DIFdel{very little }\DIFdelend \DIFaddbegin \DIFadd{even fewer }\DIFaddend methane molecules at the altitudes \DIFdelbegin \DIFdel{we can reach.This needs to be considered when working with a sensor offering lower resolution. }\DIFdelend \DIFaddbegin \DIFadd{eached by the BEXUS gondola.}\\
\DIFadd{Therefore it is necessary to pressurise the air to increase the concentration of methane particles. Due to the \mbox{%DIFAUXCMD
\gls{lod} }\hskip0pt%DIFAUXCMD
of the K96 of 3\,ppm it is not enough to keep the pressure at 1\,bar. Instead it is necessary to increase the pressure to at least 3\,bar to be able to measure with the K96 even at altitudes of 20\,km where the molar fraction could be already as low as 1\,ppm. Having a pressure of 3 bar in the measurement chamber would increase the "effective molar fraction" to 3\,ppm in that case, just high enough to rely on the sensor. }\\
\DIFadd{The sensor is susceptible to other factors causing drift such as temperature and humidity. These will require a thorough calibration strategy, which allows to measure the methane concentration under varying conditions. An AI-based calibration approach has been tested before on the K96 sensor\mbox{%DIFAUXCMD
\cite{gia_accurate_2025} }\hskip0pt%DIFAUXCMD
and therefore a similar approach could be employed for \mbox{%DIFAUXCMD
\gls{mirage}}\hskip0pt%DIFAUXCMD
. At the same time, the experiment should limit the temperature range of the air to be measured, which requires a thermal control system.}\\
\DIFaddend 



%DIF > %%%%%%%%%%%%%%%%%%%%%%%%%%%%%%%%%%%%%%%%%%%%%%%%%%%%%%%%%%%%%%%%%%%%%%%%%%%%%%
\section{Mission Statement}

As explained in Section 1.1 methane measurements are of great importance for climate modelling.
However, systems with which methane measurements are taken, are oftentimes expensive and complex in setup.
Therefore, \DIFdelbegin \DIFdel{our mission is }\DIFdelend \DIFaddbegin \DIFadd{\mbox{%DIFAUXCMD
\gls{mirage} }\hskip0pt%DIFAUXCMD
aims }\DIFaddend to demonstrate a low-cost, low-complexity solution for measuring methane concentrations in the upper stratosphere.
By using the novel K96 non-dispersive infrared (NDIR) sensor developed by Senseair \DIFaddbegin \DIFadd{AB }\DIFaddend (a Swedish company, specialised in state-of-the-art gas sensors) \cite{wastine2022} \DIFdelbegin \DIFdel{we will use }\DIFdelend off-the-shelf technology \DIFaddbegin \DIFadd{will be shown }\DIFaddend to provide an alternative to laser-based approaches.

%DIF > %%%%%%%%%%%%%%%%%%%%%%%%%%%%%%%%%%%%%%%%%%%%%%%%%%%%%%%%%%%%%%%%%%%%%%%%%%%%%%
\section{Experiment Objectives}

\begin{table}[H]
\centering
\caption{Experimental Objectives}
\label{tab:objectives}
\DIFdelbeginFL %DIFDELCMD < \begin{tabularx}{0.9\textwidth}{lX}
%DIFDELCMD < \toprule
%DIFDELCMD < \textbf{Priority} & \textbf{Description} \\
%DIFDELCMD < \midrule
%DIFDELCMD < Technical - Primary & Measure methane with an absolute concentration error of $\mathrm{\Delta}c \leq 0.1$ ppm in the stratosphere. \\
%DIFDELCMD < Technical - Primary & Stay within a budget of material (BoM) of 1000 €. \\
%DIFDELCMD < Scientific - Secondary & Validate recorded measurements via chemical transport models. \\
%DIFDELCMD < Scientific - Secondary & Search for fluctuations in recorded methane concentrations, possibly linked to emission regions. \\
%DIFDELCMD < \bottomrule
%DIFDELCMD < \end{tabularx}
%DIFDELCMD < %%%
\DIFdelendFL \DIFaddbeginFL \begin{tabularx}{0.9\textwidth}{lX}
\toprule
\textbf{Priority} & \textbf{Description} \\
\midrule
Technical - Primary & Measure methane, carbon dioxide and water with an NDIR sensor core up to at least 20\,km altitude.\\
Technical - Primary & Demonstrate a pressurisation system that can maintain 3 bar up to at least 20km altitude.\\
Scientific - Secondary & Validate recorded measurements via chemical transport models, satellite data and previous in-situ measurements by other groups.\\
Scientific - Secondary & Evaluate decrease of measurement uncertainty in the field by employing AI-based calibration.\\
\bottomrule
\end{tabularx}
\DIFaddendFL \end{table}

%DIF > %%%%%%%%%%%%%%%%%%%%%%%%%%%%%%%%%%%%%%%%%%%%%%%%%%%%%%%%%%%%%%%%%%%%%%%%%%%%%%
\section{Experiment Concept}

The Non-Dispersive Infra-Red (NDIR) sensor detects light absorption due to the number of absorbing particles in the optical cavity within a narrow wavelength range.
As mentioned in Section 1.1, this number is related to the concentration of the contributing gas species and their pressures and temperature.
Thus, it is possible to increase the concentration of a gas species by increasing the pressure of the gas while maintaining the same relative composition and molar fractions of the sample.

Therefore, \DIFdelbegin \DIFdel{we use }\DIFdelend an NDIR sensor fed with pressurized ambient air under temperature-controlled conditions \DIFaddbegin \DIFadd{will be used. To keep a pressure of 3\,bar up to at least 20\,km altitude is the real challenge of \mbox{%DIFAUXCMD
\gls{mirage}}\hskip0pt%DIFAUXCMD
. As explained above, ambient air needs to be pressurized to 3 bar using a combination of vaccuum pumps and a compressor, that work together. Vaccuum pumps are meant to draw in low density ambient air and pre-pressurise this air to provide air dense enough for the compressor to further pressurise it to 3 bar}\DIFaddend . \DIFdelbegin \DIFdel{The measurement resolution of the sensor is 0.1 ppm at ground level, i.e. at 1 bar . Because the uncertainty we aim for is less than 0.1 ppm, we need to pressurize the ambient air to 2 to 5 bar . Thus, we aim to use a pump to continuously pump ambient air into a pressurized chamber, in which the sensor is located. }\DIFdelend Air can leave the pressurized chamber through \DIFdelbegin \DIFdel{an electrical valve. This method of active air release can restrict the flow rate dynamically, adjusting to the pressure measured internally. }%DIFDELCMD < 

%DIFDELCMD < %%%
\DIFdel{As the detector output can drift with temperature, we need to }\DIFdelend \DIFaddbegin \DIFadd{a valve, which could be passive since the outflow rate is independent of altitude since it is limited by the speed of sound even if the pressure difference is just 2 bar as it will be on ground level. The concept and different components involved in the pressurisation system will be explained in more depth in the mechanical design section below. }\\
\DIFadd{To }\DIFaddend stabilize the temperature \DIFaddbegin \DIFadd{of the air is another complex challenge, that needs to be solved }\DIFaddend via a thermal control system\DIFdelbegin \DIFdel{. Additionally, we need information }\DIFdelend \DIFaddbegin \DIFadd{, that can handle varying ambient conditions. The air is initially very cold with a low density but due to pressurisation it can heat up substantially. At the same time, it's mass is small and it will lose heat to surrounding tubing and chamber walls. Electrical components such as pumps, that are also generating heat and active and passive heating elements need to make sure that the air is cooled and heated to the right temperatures in the right places. An extensive thermal study including more factors will analyse the thermal system needed further in the thermal design section below.}\\
\DIFadd{Additionally, information is needed }\DIFaddend about the ambient pressure and temperature, to inform the thermal control system and the flow rate of the pumps. \DIFaddbegin \DIFadd{Furthermore, this allows for calculation of the concentration of the methane particles outside of the experiment from the concentration measured via the ideal gas law.}\\
\DIFaddend As water vapor is a major factor in the uncertainty of the sensor \DIFdelbegin \DIFdel{, we also need to monitor }\DIFdelend \DIFaddbegin \DIFadd{and condensation of water could damage the sensor and other electronics, }\DIFaddend the inside and ideally outside humidity \DIFaddbegin \DIFadd{should be monitored as well}\DIFaddend .
These measurements can then be used to engage mitigation methods like \DIFdelbegin \DIFdel{stopping of }\DIFdelend \DIFaddbegin \DIFadd{closing the vent, stopping }\DIFaddend the pumps and raising \DIFdelbegin \DIFdel{of }\DIFdelend the internal temperature to \DIFdelbegin \DIFdel{prevent }\DIFdelend \DIFaddbegin \DIFadd{increase the water capacity of the air and prevent further }\DIFaddend condensation. Once the ambient humidity sensor is reporting safe levels, the \DIFaddbegin \DIFadd{pressurisation system can start working again, flushing out the more humid air. The }\DIFaddend internal temperature is to be set to normal operating conditions \DIFdelbegin \DIFdel{, and the pumps commence a full flushing cycle. This means that the pumps operate with the maximum flow rate for 2 minutes. Afterwards, the internal humidity levels are measured. If they are too high, another flushing cycle commences up to a total of 3 cycles, after which normal operations start again. }%DIFDELCMD < 

%DIFDELCMD < %%%
\DIFdel{Processing of the different sensor outputs and engaging of appropriate actions will be done via a \mbox{%DIFAUXCMD
\ac{mcu} }\hskip0pt%DIFAUXCMD
like an Arduino Nano.  
}%DIFDELCMD < 

%DIFDELCMD < %%%
\DIFdelend \DIFaddbegin \DIFadd{again as well. }\\
\DIFadd{Many sensors will be required to monitor those different ambient conditions and internal temperature, pressure and humidity. Depending on the location within the setup they need to be able to work under different conditions and fulfill different accuracy and range requirements. An analysis of their selection and the state of the art is done in the electrical design section below.}\\
\DIFadd{A \mbox{%DIFAUXCMD
\gls{mcu} }\hskip0pt%DIFAUXCMD
will be integrated into a PCB designed to handle all of the different processes that the experiment has to fulfill to ensure reliable measurements of the atmospheric methane concentrations up to an altitude of at least 20\,km.  
}\DIFaddend A block diagram showcasing the different functional/physical experiment blocks is provided below in Figure \ref{fig:block_diagram}.

\DIFdelbegin %DIFDELCMD < \begin{figure}[h]
%DIFDELCMD < %%%
\DIFdelendFL \DIFaddbeginFL \begin{figure}[t]
\DIFaddendFL \centering
\DIFdelbeginFL %DIFDELCMD < \includegraphics[width=\textwidth]{Graphics/Figures/Components block diagram.png}
%DIFDELCMD < %%%
\DIFdelendFL \DIFaddbeginFL \includegraphics[width=\textwidth]{0_images/figures/Components block diagram.png}
\DIFaddendFL \caption{Block diagram with different physical or functional components of the experiment.}
\label{fig:block_diagram}
\end{figure}

%DIF > %%%%%%%%%%%%%%%%%%%%%%%%%%%%%%%%%%%%%%%%%%%%%%%%%%%%%%%%%%%%%%%%%%%%%%%%%%%%%%
\section{Team Details}

The team of \DIFdelbegin \DIFdel{\mbox{%DIFAUXCMD
\ac{mirage} }\hskip0pt%DIFAUXCMD
}\DIFdelend \DIFaddbegin \DIFadd{\mbox{%DIFAUXCMD
\gls{mirage} }\hskip0pt%DIFAUXCMD
}\DIFaddend is still in the process of recruiting team members, especially for the bachelor thesis projects.
All team members are studying at Uppsala University.

%DIF > %%%%%%%%%%%%%%%%%%%%%%%%%%%%%%%%%%%%%%%%%%%%%%%%%%%%%%%%%%%%%%%%%%%%%%%%%%%%%%
\subsection{Contact \DIFdelbegin \DIFdel{Point}\DIFdelend \DIFaddbegin \DIFadd{point}\DIFaddend }

Team leader phone number: +46 76 456 08 02\\
Email: abdullah.alsakka.5356@student.uu.se\\
Team leader address: Studentvägen 18, Uppsala SE75234

%DIF > %%%%%%%%%%%%%%%%%%%%%%%%%%%%%%%%%%%%%%%%%%%%%%%%%%%%%%%%%%%%%%%%%%%%%%%%%%%%%%
\subsection{Team Members}

\begin{itemize}
    \item \textbf{Name:} Abdullah Alsakka \\
    \textbf{Education (Year):} Masters in Astronomy and Space Physics (2\textsuperscript{nd} year) \\
    \textbf{Role:} Team Leader and Science team

    \item \textbf{Name:} Jonathan Remmert \\
    \textbf{Education (Year):} Masters in Astronomy and Space Physics (2\textsuperscript{nd} year) \\
    \textbf{Role:} Science team

    \item \textbf{Name:} Oliver Bolin \\
    \textbf{Education (Year):} Electrical Engineering (2\textsuperscript{nd} year/ 5 year program) \\
    \textbf{Role:} Electrical Team

    \item \textbf{Name:} Max Steen \\
    \textbf{Education (Year):} Electrical Engineering (2\textsuperscript{nd} year/ 5 year program) \\
    \textbf{Role:} Electrical Team

    \item \textbf{Name:} Anna Tschesche \\
    \textbf{Education (Year):} Master in Material Science (2\textsuperscript{nd} year) \\
    \textbf{Role:} Science Team

    \item \textbf{Name:} Olle Hallberg \\
    \textbf{Education (Year):} Engineering Physics (3\textsuperscript{rd} year/ 5 year program) \\
    \textbf{Role:} Science Team

    \item \textbf{Name:} Molly Sundell \\
    \textbf{Education (Year):} Material Engineering (3\textsuperscript{rd} year/ 5 year program) \\
    \textbf{Role:} Outreach Team

    \item \textbf{Name:} Antoni Marciniewski \\
    \textbf{Education (Year):} Mechanical Engineering (1\textsuperscript{st} year) \\
    \textbf{Role:} Mechanical Team

    \item \textbf{Name:} Leo Adolphus Wattin \\
    \textbf{Education (Year):} Mechanical Engineering (2\textsuperscript{nd} year) \\
    \textbf{Role:} Mechanical Team
\end{itemize}

\DIFdelbegin \chapter{\DIFdel{Experiment Requirements and Constraints}}
%DIFAUXCMD
\addtocounter{chapter}{-1}%DIFAUXCMD
\DIFdelend %DIF > %%%%%%%%%%%%%%%%%%%%%%%%%%%%%%%%%%%%%%%%%%%%%%%%%%%%%%%%%%%%%%%%%%%%%%%%%%%%%%
\DIFaddbegin \clearpage
\newpage
\DIFaddend 

\DIFaddbegin \chapter{\DIFadd{EXPERIMENT REQUIREMENTS AND CONSTRAINTS}}
\label{sec:requirements}
\DIFaddend \section{Functional Requirements}

\begin{longtable}{|p{2cm}|p{12cm}|}
\caption{Functional Requirements} \\
\DIFdelbegin %DIFDELCMD < \toprule
%DIFDELCMD < %%%
\DIFdelend \DIFaddbegin \hline 
\DIFaddend \textbf{ID} & \textbf{Requirement text} \\
\DIFaddbegin \hline 
\DIFaddend \endfirsthead
\DIFdelbegin %DIFDELCMD < \multicolumn{2}{c}{%
{%DIFAUXCMD
%DIFDELCMD < \tablename%%%
\DIFdel{\ }%DIFDELCMD < \thetable%%%
\DIFdel{\ -- }\textit{\DIFdel{Continued from previous page}}%DIFAUXCMD
} %DIFAUXCMD
\DIFdelend \DIFaddbegin \caption[]{\DIFadd{Functional Requirements -- }\textit{\DIFadd{Continued from previous page}}} \DIFaddend \\
\DIFdelbegin %DIFDELCMD < \toprule
%DIFDELCMD < %%%
\DIFdelend \DIFaddbegin \hline
\DIFaddend \textbf{ID} & \textbf{Requirement text} \\
\DIFaddbegin \hline
\DIFaddend \endhead
\DIFdelbegin %DIFDELCMD < \bottomrule
%DIFDELCMD < %%%
\DIFdelend \DIFaddbegin \hline
\DIFaddend \multicolumn{2}{r}{\textit{Continued on next page}} \\
\endfoot
\DIFdelbegin %DIFDELCMD < \bottomrule
%DIFDELCMD < \endlastfoot
%DIFDELCMD < %%%
\DIFdelend \hline
\DIFaddbegin \endlastfoot
\DIFaddend F.1. & The experiment shall measure the atmospheric \DIFdelbegin \DIFdel{methane, CO2 and H2O }\DIFdelend \DIFaddbegin \DIFadd{CH\textsubscript{4}, CO\textsubscript{2} and H\textsubscript{2}O }\DIFaddend concentrations of ambient air simultaneously with an off-the-shelf \DIFdelbegin \DIFdel{\mbox{%DIFAUXCMD
\ac{ndir} }\hskip0pt%DIFAUXCMD
}\DIFdelend \DIFaddbegin \DIFadd{\mbox{%DIFAUXCMD
\gls{ndir} }\hskip0pt%DIFAUXCMD
}\DIFaddend sensor. \\
\hline
F.2. & The measurement chamber shall be pressurised with ambient air. \\
\hline
F.3. & The air in the measurement chamber shall be exchanged\DIFdelbegin \DIFdel{continuously}\DIFdelend . \\
\hline
F.4. & The experiment shall measure \DIFaddbegin \DIFadd{the }\DIFaddend pressure inside of the measurement chamber. \\
\hline
F.5. & The experiment shall measure the temperature inside of the measurement chamber. \\
\hline
F.6\DIFaddbegin \DIFadd{. }\DIFaddend & The experiment shall stabilize the temperature inside \DIFaddbegin \DIFadd{of }\DIFaddend the measurement chamber. \\
\hline
F.7\DIFaddbegin \DIFadd{. }\DIFaddend & The experiment shall measure the humidity inside of the measurement chamber. \\
\hline
F.8\DIFaddbegin \DIFadd{. }\DIFaddend & The experiment shall prevent the condensation of water inside \DIFaddbegin \DIFadd{of }\DIFaddend the measurement chamber. \\
\hline
%DIF > F.9. & The experiment should change between measurement times based on the flight phase of the balloon. \\
%DIF > \hline
F.9\DIFaddbegin \DIFadd{. }\DIFaddend & \DIFdelbegin \DIFdel{The experiment should change between measurement times based on the flight phase of the balloon. }%DIFDELCMD < \\
%DIFDELCMD < \hline
%DIFDELCMD < %%%
\DIFdel{F.10 }%DIFDELCMD < & %%%
\DIFdelend The experiment should measure ambient temperature outside. \\
\hline
F\DIFdelbegin \DIFdel{.11 }\DIFdelend \DIFaddbegin \DIFadd{.10. }\DIFaddend & The experiment should measure ambient pressure outside. \\
\hline
F\DIFdelbegin \DIFdel{.12 }\DIFdelend \DIFaddbegin \DIFadd{.11. }\DIFaddend & The experiment should measure ambient humidity outside. \\
\hline
F\DIFdelbegin \DIFdel{.13 }\DIFdelend \DIFaddbegin \DIFadd{.12. }\DIFaddend & The experiment shall send data via E-Link. \\
\hline
F\DIFdelbegin \DIFdel{.14 }\DIFdelend \DIFaddbegin \DIFadd{.13. }\DIFaddend & The experiment \DIFdelbegin \DIFdel{should }\DIFdelend \DIFaddbegin \DIFadd{shall }\DIFaddend receive data via E-Link. \\
\DIFdelbegin %DIFDELCMD < \hline
%DIFDELCMD < %%%
\DIFdelend \end{longtable}

%DIF > %%%%%%%%%%%%%%%%%%%%%%%%%%%%%%%%%%%%%%%%%%%%%%%%%%%%%%%%%%%%%%%%%%%%%%%%%%%%%%
\section{Performance Requirements}

\begin{longtable}{|p{2cm}|p{12cm}|}
\caption{Performance Requirements}
\DIFaddbegin \label{tab:performance_req}\DIFaddend \\
\DIFdelbegin %DIFDELCMD < \toprule
%DIFDELCMD < %%%
\DIFdelend \DIFaddbegin \hline
\DIFaddend \textbf{ID} & \textbf{Requirement text} \\
\DIFaddbegin \hline
\DIFaddend \endfirsthead
\DIFdelbegin %DIFDELCMD < \multicolumn{2}{c}{%
{%DIFAUXCMD
%DIFDELCMD < \tablename%%%
\DIFdel{\ }%DIFDELCMD < \thetable%%%
\DIFdel{\ -- }\textit{\DIFdel{Continued from previous page}}%DIFAUXCMD
} %DIFAUXCMD
\DIFdelend \DIFaddbegin \caption[]{\DIFadd{Performance Requirements -- }\textit{\DIFadd{Continued from previous page}}} \DIFaddend \\
\DIFdelbegin %DIFDELCMD < \toprule
%DIFDELCMD < %%%
\DIFdelend \DIFaddbegin \hline
\DIFaddend \textbf{ID} & \textbf{Requirement text} \\
\DIFaddbegin \hline
\DIFaddend \endhead
\DIFdelbegin %DIFDELCMD < \bottomrule
%DIFDELCMD < %%%
\DIFdelend \DIFaddbegin \hline
\DIFaddend \multicolumn{2}{r}{\textit{Continued on next page}} \\
\endfoot
\DIFdelbegin %DIFDELCMD < \bottomrule
%DIFDELCMD < \endlastfoot
%DIFDELCMD < %%%
\DIFdelend \hline
\DIFaddbegin \endlastfoot
\DIFaddend P.1.1 & The \DIFdelbegin \DIFdel{\mbox{%DIFAUXCMD
\ac{ndir} }\hskip0pt%DIFAUXCMD
}\DIFdelend \DIFaddbegin \DIFadd{\mbox{%DIFAUXCMD
\gls{ndir} }\hskip0pt%DIFAUXCMD
}\DIFaddend sensor shall measure the concentration of methane with a range of 0.4 to 3.0 ppm. \\
\hline
P.1.2 \DIFdelbegin \DIFdel{. }\DIFdelend & The \DIFdelbegin \DIFdel{\mbox{%DIFAUXCMD
\ac{ndir} }\hskip0pt%DIFAUXCMD
}\DIFdelend \DIFaddbegin \DIFadd{\mbox{%DIFAUXCMD
\gls{ndir} }\hskip0pt%DIFAUXCMD
}\DIFaddend sensor shall measure the concentration of methane with an uncertainty of \DIFdelbegin \DIFdel{0.1 }\DIFdelend \DIFaddbegin \DIFadd{0.5 }\DIFaddend ppm or less. \\
\hline
P.2.1 & The \DIFdelbegin \DIFdel{\mbox{%DIFAUXCMD
\ac{ndir} }\hskip0pt%DIFAUXCMD
}\DIFdelend \DIFaddbegin \DIFadd{\mbox{%DIFAUXCMD
\gls{ndir} }\hskip0pt%DIFAUXCMD
}\DIFaddend sensor shall measure the concentration of \DIFdelbegin \DIFdel{CO\textsubscript{2} }\DIFdelend \DIFaddbegin \DIFadd{carbondioxide }\DIFaddend with a range of 200 to 500 ppm. \\
\hline
P.2.2 & The \DIFdelbegin \DIFdel{\mbox{%DIFAUXCMD
\ac{ndir} }\hskip0pt%DIFAUXCMD
}\DIFdelend \DIFaddbegin \DIFadd{\mbox{%DIFAUXCMD
\gls{ndir} }\hskip0pt%DIFAUXCMD
}\DIFaddend sensor shall measure the concentration of \DIFdelbegin \DIFdel{CO\textsubscript{2} }\DIFdelend \DIFaddbegin \DIFadd{carbondioxide }\DIFaddend with an uncertainty of \DIFdelbegin \DIFdel{1ppm }\DIFdelend \DIFaddbegin \DIFadd{5 ppm }\DIFaddend or less. \\
\hline
P.3.1 & The \DIFdelbegin \DIFdel{\mbox{%DIFAUXCMD
\ac{ndir} }\hskip0pt%DIFAUXCMD
}\DIFdelend \DIFaddbegin \DIFadd{\mbox{%DIFAUXCMD
\gls{ndir} }\hskip0pt%DIFAUXCMD
}\DIFaddend sensor shall measure the concentration of \DIFdelbegin \DIFdel{H\textsubscript{2}O }\DIFdelend \DIFaddbegin \DIFadd{water }\DIFaddend with a range of 1 to 30,\DIFdelbegin \DIFdel{000ppm}\DIFdelend \DIFaddbegin \DIFadd{000 ppm}\DIFaddend . \\
\hline
P.3.2 & The \DIFdelbegin \DIFdel{\mbox{%DIFAUXCMD
\ac{ndir} }\hskip0pt%DIFAUXCMD
}\DIFdelend \DIFaddbegin \DIFadd{\mbox{%DIFAUXCMD
\gls{ndir} }\hskip0pt%DIFAUXCMD
}\DIFaddend sensor shall measure the concentration of \DIFdelbegin \DIFdel{H\textsubscript{2}O }\DIFdelend \DIFaddbegin \DIFadd{water }\DIFaddend with an uncertainty of \DIFdelbegin \DIFdel{1ppm or lessat all altitudes}\DIFdelend \DIFaddbegin \DIFadd{5 ppm or less}\DIFaddend . \\
\hline
%DIF > P.4.1 & The \gls{ndir} sensor shall have a measurement time for the CH\textsubscript{4}, CO\textsubscript{2} and H\textsubscript{2}O concentrations of 20 s during ascent and descent. \\
%DIF > \hline
%DIF > P.4.2 & The \gls{ndir} sensor shall have a measurement time for the CH\textsubscript{4}, CO\textsubscript{2} and H\textsubscript{2}O concentrations of 50 s during the floating phase. \\
%DIF > \hline
P.4\DIFaddbegin \DIFadd{.1 }\DIFaddend & The pressure inside \DIFaddbegin \DIFadd{of }\DIFaddend the measurement chamber \DIFdelbegin \DIFdel{during methane measurements shall be at least 2 barduring measurements}\DIFdelend \DIFaddbegin \DIFadd{shall stay within a range of 2.5 and 3.2 bar}\DIFaddend . \\
\hline
P\DIFdelbegin \DIFdel{.5.1 }\DIFdelend \DIFaddbegin \DIFadd{.4.2 }\DIFaddend & The \DIFdelbegin \DIFdel{air }\DIFdelend \DIFaddbegin \DIFadd{pressure }\DIFaddend inside of the measurement chamber \DIFdelbegin \DIFdel{shall be exchanged continuously}\DIFdelend \DIFaddbegin \DIFadd{should stay within a range of 2.8 and 3 bar}\DIFaddend . \\
\hline
P.5 \DIFdelbegin \DIFdel{.2 }\DIFdelend & The \DIFdelbegin \DIFdel{air should be sufficiently (97\%) exchanged within at least one minute}\DIFdelend \DIFaddbegin \DIFadd{pressurisation system shall have an air throughput of at least 1 l/min}\DIFaddend .\\
\hline
P.6.1 & The experiment shall measure the pressure inside \DIFaddbegin \DIFadd{of }\DIFaddend the measurement chamber in a range of 0.3 to \DIFdelbegin \DIFdel{5.1 }\DIFdelend \DIFaddbegin \DIFadd{4 }\DIFaddend bar. \\
\hline
P.6.2 & The experiment shall measure the pressure inside \DIFaddbegin \DIFadd{of }\DIFaddend the measurement chamber with a maximum uncertainty of 0.02 \DIFdelbegin \DIFdel{m}\DIFdelend bar . \\
\hline
P.7 \DIFdelbegin \DIFdel{.1 }\DIFdelend & The \DIFdelbegin \DIFdel{experiment shall measure the temperature inside }\DIFdelend \DIFaddbegin \DIFadd{temperature of the air inside of }\DIFaddend the measurement chamber \DIFdelbegin \DIFdel{in }\DIFdelend \DIFaddbegin \DIFadd{shall stay within }\DIFaddend a range of \DIFdelbegin \DIFdel{-40 to +60 }\DIFdelend \DIFaddbegin \DIFadd{15 and 40}\DIFaddend \degree C \DIFaddbegin \DIFadd{while active}\DIFaddend . \\
\hline
P\DIFdelbegin \DIFdel{.7.2 }\DIFdelend \DIFaddbegin \DIFadd{.8 }\DIFaddend & The \DIFdelbegin \DIFdel{experiment shall measure the temperature inside }\DIFdelend \DIFaddbegin \DIFadd{temperature of the air inside of }\DIFaddend the measurement chamber \DIFdelbegin \DIFdel{with a maximum uncertainty of 1.5}\DIFdelend \DIFaddbegin \DIFadd{should stay within a range of 19 and 24}\DIFaddend \degree C \DIFaddbegin \DIFadd{while active}\DIFaddend . \\
\hline
P\DIFdelbegin \DIFdel{.8}\DIFdelend \DIFaddbegin \DIFadd{.9}\DIFaddend .1 & The experiment shall measure the \DIFdelbegin \DIFdel{humidity inside }\DIFdelend \DIFaddbegin \DIFadd{temperature inside of }\DIFaddend the measurement chamber in a range of \DIFdelbegin \DIFdel{0 to 100\% r. H. }\DIFdelend \DIFaddbegin \DIFadd{-40 to +85 }\degree \DIFadd{C. }\DIFaddend \\
\hline
P\DIFdelbegin \DIFdel{.8}\DIFdelend \DIFaddbegin \DIFadd{.9}\DIFaddend .2 & The experiment shall measure the \DIFdelbegin \DIFdel{humidity inside }\DIFdelend \DIFaddbegin \DIFadd{temperature inside of }\DIFaddend the measurement chamber with a maximum uncertainty of \DIFdelbegin \DIFdel{3\% r. H. }\DIFdelend \DIFaddbegin \DIFadd{1.5}\degree \DIFadd{C. }\DIFaddend \\
\hline
P\DIFdelbegin \DIFdel{.9 }\DIFdelend \DIFaddbegin \DIFadd{.10 }\DIFaddend & The \DIFdelbegin \DIFdel{temperature of the air inside }\DIFdelend \DIFaddbegin \DIFadd{experiment shall keep the humidity inside of }\DIFaddend the measurement chamber \DIFdelbegin \DIFdel{shall stay within a range of 0 and 40}%DIFDELCMD < \degree %%%
\DIFdel{C.}\DIFdelend \DIFaddbegin \DIFadd{below 90\% r.H. }\DIFaddend \\
\hline
P\DIFdelbegin \DIFdel{.10 }\DIFdelend \DIFaddbegin \DIFadd{.11.1 }\DIFaddend & The experiment \DIFdelbegin \DIFdel{should keep the temperature of the air inside the measurement chamber within }\DIFdelend \DIFaddbegin \DIFadd{shall measure the humidity inside of the measurement chamber in }\DIFaddend a range of \DIFdelbegin \DIFdel{19 and 23}%DIFDELCMD < \degree %%%
\DIFdel{C.}\DIFdelend \DIFaddbegin \DIFadd{0 to 100\% r.H. }\DIFaddend \\
\hline
P.11\DIFaddbegin \DIFadd{.2 }\DIFaddend & The experiment shall \DIFdelbegin \DIFdel{keep }\DIFdelend \DIFaddbegin \DIFadd{measure }\DIFaddend the humidity inside of the measurement chamber \DIFdelbegin \DIFdel{below 90}\DIFdelend \DIFaddbegin \DIFadd{with a maximum uncertainty of 3}\DIFaddend \% r.H. \\
\hline
P.12 & The experiment should measure the temperature, pressure and humidity inside of the measurement chamber at a rate \DIFdelbegin \DIFdel{of 10Hz}\DIFdelend \DIFaddbegin \DIFadd{over 10 Hz}\DIFaddend . \\
\hline
P.13.1 & The \DIFdelbegin \DIFdel{\mbox{%DIFAUXCMD
\ac{ndir} }\hskip0pt%DIFAUXCMD
sensor shall have a measurement time for the CH4, CO2 and H2O concentrations of 20s during ascent and descent. }%DIFDELCMD < \\
%DIFDELCMD < \hline
%DIFDELCMD < %%%
\DIFdel{P.13.2 }%DIFDELCMD < & %%%
\DIFdel{The \mbox{%DIFAUXCMD
\ac{ndir} }\hskip0pt%DIFAUXCMD
sensor shall have a measurement time for the CH4, CO2 and H2O concentrations of 50s during the floating phase. }%DIFDELCMD < \\
%DIFDELCMD < \hline
%DIFDELCMD < %%%
\DIFdel{P.14.1 }%DIFDELCMD < & %%%
\DIFdel{The }\DIFdelend experiment should measure the ambient pressure outside in a range of 20 to 1100 mbar. \\
\hline
P\DIFdelbegin \DIFdel{.14}\DIFdelend \DIFaddbegin \DIFadd{.13}\DIFaddend .2 & The experiment should measure the ambient pressure outside with a maximum uncertainty of \DIFdelbegin \DIFdel{0.05 }\DIFdelend \DIFaddbegin \DIFadd{20 }\DIFaddend mbar. \\
\hline
P\DIFdelbegin \DIFdel{.15}\DIFdelend \DIFaddbegin \DIFadd{.14}\DIFaddend .1 & The experiment should measure the ambient temperature outside in a range of \DIFdelbegin \DIFdel{-100 }\DIFdelend \DIFaddbegin \DIFadd{-70 }\DIFaddend to +60 \degree C. \\
\hline
P\DIFdelbegin \DIFdel{.15}\DIFdelend \DIFaddbegin \DIFadd{.14}\DIFaddend .2 & The experiment should measure the ambient temperature outside with a maximum uncertainty of 1.5\degree C. \\
\hline
P\DIFdelbegin \DIFdel{.16}\DIFdelend \DIFaddbegin \DIFadd{.15}\DIFaddend .1 & The experiment should measure the ambient humidity outside in a range of 0 to 100\DIFaddbegin \DIFadd{\% }\DIFaddend r.H. \\
\hline
P\DIFdelbegin \DIFdel{.16}\DIFdelend \DIFaddbegin \DIFadd{.15}\DIFaddend .2 & The experiment should measure the \DIFdelbegin \DIFdel{humidity outside the sensor box }\DIFdelend \DIFaddbegin \DIFadd{ambient humidity outside }\DIFaddend with a maximum uncertainty of 3\% r.H. \\
\hline
P\DIFdelbegin \DIFdel{.17 }\DIFdelend \DIFaddbegin \DIFadd{.16 }\DIFaddend & The experiment should measure the ambient temperature, pressure and humidity outside at a rate \DIFdelbegin \DIFdel{of 10Hz}\DIFdelend \DIFaddbegin \DIFadd{over 10 Hz}\DIFaddend . \\
\hline
P\DIFdelbegin \DIFdel{.18 }\DIFdelend \DIFaddbegin \DIFadd{.17 }\DIFaddend & The experiment should receive data via E-Link with an uplink rate \DIFdelbegin \DIFdel{of 10kbps}\DIFdelend \DIFaddbegin \DIFadd{under 10 kbps}\DIFaddend . \\
\hline
P\DIFdelbegin \DIFdel{.19 }\DIFdelend \DIFaddbegin \DIFadd{.18 }\DIFaddend & The experiment shall send the obtained data with the E-Link with a downlink rate \DIFdelbegin \DIFdel{of }\DIFdelend \DIFaddbegin \DIFadd{under }\DIFaddend 100 kbps. \\
\DIFdelbegin %DIFDELCMD < \hline
%DIFDELCMD < %%%
\DIFdelend \end{longtable}

%DIF > %%%%%%%%%%%%%%%%%%%%%%%%%%%%%%%%%%%%%%%%%%%%%%%%%%%%%%%%%%%%%%%%%%%%%%%%%%%%%%
\section{Design Requirements}

\begin{longtable}{|p{2cm}|p{12cm}|}
\caption{Design Requirements} \\
\DIFdelbegin %DIFDELCMD < \toprule
%DIFDELCMD < %%%
\DIFdelend \DIFaddbegin \hline
\DIFaddend \textbf{ID} & \textbf{Requirement text} \\
\DIFaddbegin \hline
\DIFaddend \endfirsthead
\DIFdelbegin %DIFDELCMD < \multicolumn{2}{c}{%
{%DIFAUXCMD
%DIFDELCMD < \tablename%%%
\DIFdel{\ }%DIFDELCMD < \thetable%%%
\DIFdel{\ -- }\textit{\DIFdel{Continued from previous page}}%DIFAUXCMD
} %DIFAUXCMD
\DIFdelend \DIFaddbegin \caption[]{\DIFadd{Design Requirements -- }\textit{\DIFadd{Continued from previous page}}} \DIFaddend \\
\DIFdelbegin %DIFDELCMD < \toprule
%DIFDELCMD < %%%
\DIFdelend \DIFaddbegin \label{tab:design_req}\\
\hline
\DIFaddend \textbf{ID} & \textbf{Requirement text} \\
\DIFaddbegin \hline
\DIFaddend \endhead
\DIFdelbegin %DIFDELCMD < \bottomrule
%DIFDELCMD < %%%
\DIFdelend \DIFaddbegin \hline
\DIFaddend \multicolumn{2}{r}{\textit{Continued on next page}} \\
\endfoot
\DIFdelbegin %DIFDELCMD < \bottomrule
%DIFDELCMD < \endlastfoot
%DIFDELCMD < %%%
\DIFdelend \hline
\DIFaddbegin \endlastfoot
\DIFaddend D.1\DIFaddbegin \DIFadd{.1 }\DIFaddend & The experiment shall operate in the temperature profile of the BEXUS vehicle flight and launch. \\
\hline
D\DIFaddbegin \DIFadd{.1}\DIFaddend .2 \DIFaddbegin & \DIFadd{The temperature in the electronic compartment shall stay within a range of -25 and 85°C.}\\
\hline
\DIFadd{D}\DIFaddend .1\DIFaddbegin \DIFadd{.3 }\DIFaddend & The \DIFaddbegin \DIFadd{temperature in the electronic compartment should stay within a range of -15 and 60°C.}\\
\hline
\DIFadd{D.2.1 }& \DIFadd{The }\DIFaddend experiment shall operate in the pressure profile of the BEXUS vehicle flight and launch. \\
\hline
D.2.2 & The measurement chamber shall withstand a pressure difference of up to \DIFdelbegin \DIFdel{5.5 }\DIFdelend \DIFaddbegin \DIFadd{4 }\DIFaddend bar. \\
\hline
D.3 & The experiment shall operate in the vibration profile of the BEXUS vehicle flight and launch. \\
\hline
D.4 & The experiment \DIFaddbegin \DIFadd{shall operate in the shock profile of the BEXUS vehicle flight and launch. }\\
\hline
\DIFadd{D.5 }& \DIFadd{The experiment }\DIFaddend power connector shall be an Amphenol PT02E8-4P, Code A.\\
\hline
D\DIFdelbegin \DIFdel{.5 }\DIFdelend \DIFaddbegin \DIFadd{.6 }\DIFaddend & The experiment \DIFdelbegin \DIFdel{shall }\DIFdelend \DIFaddbegin \DIFadd{data connector shall be an Amphenol RJF21B connector. }\\
\hline
\DIFadd{D.7 }& \DIFadd{The experiment shall }\DIFaddend not disturb or harm the launch vehicle. \\
\hline
D\DIFdelbegin \DIFdel{.6 }\DIFdelend \DIFaddbegin \DIFadd{.8 }\DIFaddend & \DIFdelbegin \DIFdel{The \mbox{%DIFAUXCMD
\ac{ndir} }\hskip0pt%DIFAUXCMD
sensor shall be separate from the other electronics in the measurement chamber , which will be flushed and pressurised before each measurement}\DIFdelend \DIFaddbegin \DIFadd{Only the measurement chamber containing the \mbox{%DIFAUXCMD
\gls{ndir} }\hskip0pt%DIFAUXCMD
sensor core and an additional pressure sensor shall be pressurised}\DIFaddend .\\
\hline
D\DIFdelbegin \DIFdel{.7 }\DIFdelend \DIFaddbegin \DIFadd{.9 }\DIFaddend & The experiment shall consume less than 150 Wh during the flight. \\
\hline
D\DIFdelbegin \DIFdel{.8 }\DIFdelend \DIFaddbegin \DIFadd{.10 }\DIFaddend & The experiment shall \DIFaddbegin \DIFadd{operate under 28.8\,V. }\\
\hline
\DIFadd{D.11.1 }& \DIFadd{The experiment shall }\DIFaddend have no currents exceeding \DIFaddbegin \DIFadd{2.4 A. }\\
\hline
\DIFadd{D.11.2 }& \DIFadd{The experiment should have no currents exceeding }\DIFaddend 1.4 A\DIFdelbegin \DIFdel{at any time during the flight}\DIFdelend . \\
\hline
D\DIFdelbegin \DIFdel{.9 }\DIFdelend \DIFaddbegin \DIFadd{.12 }\DIFaddend & The experiment should stay within a size limit of \DIFdelbegin \DIFdel{12x12x5 }\DIFdelend \DIFaddbegin \DIFadd{15x20x20 }\DIFaddend cm\textsuperscript{3}. \\
\hline
D\DIFdelbegin \DIFdel{.10 }\DIFdelend \DIFaddbegin \DIFadd{.13 }\DIFaddend & The experiment \DIFdelbegin \DIFdel{should not exceed a mass of 2 }\DIFdelend \DIFaddbegin \DIFadd{shall stay within a weight limit of 5 }\DIFaddend kg. \\
\hline
D\DIFdelbegin \DIFdel{.11 }\DIFdelend \DIFaddbegin \DIFadd{.14 }\DIFaddend & The air sample to be measured shall not be changed in its composition by the experiment, the vehicle or other experiments. \\
\hline
D\DIFdelbegin \DIFdel{.12 }\DIFdelend \DIFaddbegin \DIFadd{.15 }\DIFaddend & The experiment shall have access to fresh air from outside. \\
\hline
D\DIFdelbegin \DIFdel{.13 }\DIFdelend \DIFaddbegin \DIFadd{.16 }\DIFaddend & The data storage shall be \DIFaddbegin \DIFadd{able to store at least 2GB. }\\
\hline
\DIFadd{D.17 }& \DIFadd{The data storage shall be }\DIFaddend protected from water due to condensation or during landing. \\
\hline
D\DIFdelbegin \DIFdel{.14 }\DIFdelend \DIFaddbegin \DIFadd{.18 }\DIFaddend & The \DIFdelbegin \DIFdel{\mbox{%DIFAUXCMD
\ac{ndir} }\hskip0pt%DIFAUXCMD
}\DIFdelend \DIFaddbegin \DIFadd{\mbox{%DIFAUXCMD
\gls{ndir} }\hskip0pt%DIFAUXCMD
}\DIFaddend sensor should still function after recovery post flight. \\
\hline
D\DIFdelbegin \DIFdel{.15 }\DIFdelend \DIFaddbegin \DIFadd{.19 }\DIFaddend & The budget for all components except the K96 sensor should stay below \DIFdelbegin \DIFdel{1000}\DIFdelend \DIFaddbegin \DIFadd{2000}\DIFaddend \texteuro. \\
\DIFdelbegin %DIFDELCMD < \hline
%DIFDELCMD < %%%
\DIFdelend \end{longtable}

%DIF > %%%%%%%%%%%%%%%%%%%%%%%%%%%%%%%%%%%%%%%%%%%%%%%%%%%%%%%%%%%%%%%%%%%%%%%%%%%%%%
\section{Operational Requirements}

\begin{longtable}{|p{2cm}|p{12cm}|}
\caption{Operational Requirements} \\
\DIFdelbegin %DIFDELCMD < \toprule
%DIFDELCMD < %%%
\DIFdelend \DIFaddbegin \hline
\DIFaddend \textbf{ID} & \textbf{Requirement text} \\
\DIFaddbegin \hline
\DIFaddend \endfirsthead
\DIFdelbegin %DIFDELCMD < \multicolumn{2}{c}{%
{%DIFAUXCMD
%DIFDELCMD < \tablename%%%
\DIFdel{\ }%DIFDELCMD < \thetable%%%
\DIFdel{\ -- }\textit{\DIFdel{Continued from previous page}}%DIFAUXCMD
} %DIFAUXCMD
\DIFdelend \DIFaddbegin \caption[]{\DIFadd{Operational Requirements -- }\textit{\DIFadd{Continued from previous page}}} \DIFaddend \\
\DIFdelbegin %DIFDELCMD < \toprule
%DIFDELCMD < %%%
\DIFdelend \DIFaddbegin \hline
\DIFaddend \textbf{ID} & \textbf{Requirement text} \\
\DIFaddbegin \hline
\DIFaddend \endhead
\DIFdelbegin %DIFDELCMD < \bottomrule
%DIFDELCMD < %%%
\DIFdelend \DIFaddbegin \hline
\DIFaddend \multicolumn{2}{r}{\textit{Continued on next page}} \\
\endfoot
\DIFdelbegin %DIFDELCMD < \bottomrule
%DIFDELCMD < \endlastfoot
%DIFDELCMD < %%%
\DIFdelend \hline
\DIFaddbegin \endlastfoot
\DIFaddend O.1 & The experiment shall be able to conduct measurements autonomously in case connection with the ground segment is lost. \\
\hline
O.2 & The experiment shall accept a request for radio silence at any time while on the launch pad. \\
\hline
O.3 & The experiment \DIFdelbegin \DIFdel{should switch automatically to a longer or shorter measurement time depending on the change in altitude. }%DIFDELCMD < \\
%DIFDELCMD < \hline
%DIFDELCMD < %%%
\DIFdel{O.4 }%DIFDELCMD < & %%%
\DIFdel{The experiment shall respond to manual changes of measurement time from the ground. }%DIFDELCMD < \\
%DIFDELCMD < \hline
%DIFDELCMD < %%%
\DIFdel{O.5 }%DIFDELCMD < & %%%
\DIFdel{The experiment shall }\DIFdelend \DIFaddbegin \DIFadd{shall }\DIFaddend store all data on board and continuously send \DIFdelbegin \DIFdel{condensed data sets }\DIFdelend \DIFaddbegin \DIFadd{data }\DIFaddend to ground via E-Link. \\
\hline
O\DIFdelbegin \DIFdel{.6 }\DIFdelend \DIFaddbegin \DIFadd{.4 }\DIFaddend & The experiment shall automatically register and respond to high r.H. inside the chamber in order to lower the humidity. \\
\DIFdelbegin %DIFDELCMD < \hline
%DIFDELCMD < %%%
\DIFdelend \end{longtable}

%DIF > %%%%%%%%%%%%%%%%%%%%%%%%%%%%%%%%%%%%%%%%%%%%%%%%%%%%%%%%%%%%%%%%%%%%%%%%%%%%%%
\section{Constraints}

\begin{longtable}{|p{2cm}|p{12cm}|}
\caption{Constraints} \\
\DIFdelbegin %DIFDELCMD < \toprule
%DIFDELCMD < %%%
\DIFdelend \DIFaddbegin \hline
\DIFaddend \textbf{ID} & \textbf{Requirement text} \\
\DIFaddbegin \hline
\DIFaddend \endfirsthead
\DIFdelbegin %DIFDELCMD < \multicolumn{2}{c}{%
{%DIFAUXCMD
%DIFDELCMD < \tablename%%%
\DIFdel{\ }%DIFDELCMD < \thetable%%%
\DIFdel{\ -- }\textit{\DIFdel{Continued from previous page}}%DIFAUXCMD
} %DIFAUXCMD
\DIFdelend \DIFaddbegin \caption[]{\DIFadd{Constraints -- }\textit{\DIFadd{Continued from previous page}}} \DIFaddend \\
\DIFdelbegin %DIFDELCMD < \toprule
%DIFDELCMD < %%%
\DIFdelend \DIFaddbegin \hline
\DIFaddend \textbf{ID} & \textbf{Requirement text} \\
\DIFaddbegin \hline
\DIFaddend \endhead
\DIFdelbegin %DIFDELCMD < \bottomrule
%DIFDELCMD < %%%
\DIFdelend \DIFaddbegin \hline
\DIFaddend \multicolumn{2}{r}{\textit{Continued on next page}} \\
\endfoot
\DIFdelbegin %DIFDELCMD < \bottomrule
%DIFDELCMD < \endlastfoot
%DIFDELCMD < %%%
\DIFdelend \hline
\DIFaddbegin \endlastfoot
\DIFaddend C.1 & The thermal chamber at the \DIFdelbegin \DIFdel{\mbox{%DIFAUXCMD
\ac{irf} }\hskip0pt%DIFAUXCMD
}\DIFdelend \DIFaddbegin \DIFadd{\mbox{%DIFAUXCMD
\gls{irf} }\hskip0pt%DIFAUXCMD
}\DIFaddend goes only down to -40\degree C. \\
\DIFdelbegin %DIFDELCMD < \hline
%DIFDELCMD < %%%
\DIFdelend \end{longtable}\DIFaddbegin \newpage
\DIFaddend 

\chapter{\DIFdelbegin \DIFdel{Project Planning}\DIFdelend \DIFaddbegin \DIFadd{PROJECT PLANNING}\DIFaddend }
\DIFaddbegin \label{sec:project_planning}
\DIFaddend 

\DIFdelbegin %DIFDELCMD < \begin{figure}
%DIFDELCMD <   \centering
%DIFDELCMD <   \includegraphics[width = \textwidth]{Graphics/Figures/wbs_manegerial.png}
%DIFDELCMD <   %%%
%DIFDELCMD < \caption{%
{%DIFAUXCMD
\DIFdelFL{Work breakdown structure for the manegerial tasks}}
  %DIFAUXCMD
%DIFDELCMD < \label{fig:wbs1}
%DIFDELCMD < \end{figure}
%DIFDELCMD < \begin{figure}
%DIFDELCMD <   \centering
%DIFDELCMD <   \includegraphics[width = \textwidth]{Graphics/Figures/wbs_science.png}
%DIFDELCMD <   %%%
%DIFDELCMD < \caption{%
{%DIFAUXCMD
\DIFdelFL{Work breakdown structure for the science tasks}}
  %DIFAUXCMD
%DIFDELCMD < \label{fig:wbs2}
%DIFDELCMD < \end{figure}
%DIFDELCMD < \begin{figure}
%DIFDELCMD <   \centering
%DIFDELCMD <   \includegraphics[width = \textwidth]{Graphics/Figures/wbs_electrical.png}
%DIFDELCMD <   %%%
%DIFDELCMD < \caption{%
{%DIFAUXCMD
\DIFdelFL{Work breakdown structure for the electrical tasks}}
  %DIFAUXCMD
%DIFDELCMD < \label{fig:wbs3}
%DIFDELCMD < \end{figure}
%DIFDELCMD < \begin{figure}
%DIFDELCMD <   \centering
%DIFDELCMD <   \includegraphics[width = \textwidth]{Graphics/Figures/wbs_mechanical.png}
%DIFDELCMD <   %%%
%DIFDELCMD < \caption{%
{%DIFAUXCMD
\DIFdelFL{Work breakdown structure for the mechanical tasks}}
  %DIFAUXCMD
%DIFDELCMD < \label{fig:wbs4}
%DIFDELCMD < \end{figure}
%DIFDELCMD < \begin{figure}
%DIFDELCMD <   \centering
%DIFDELCMD <   \includegraphics[width = \textwidth]{Graphics/Figures/wbs_software.png}
%DIFDELCMD <   %%%
%DIFDELCMD < \caption{%
{%DIFAUXCMD
\DIFdelFL{Work breakdown structure for the software tasks}}
  %DIFAUXCMD
%DIFDELCMD < \label{fig:wbs5}
%DIFDELCMD < \end{figure}
%DIFDELCMD < \begin{figure}
%DIFDELCMD <   \centering
%DIFDELCMD <   \includegraphics[width = \textwidth]{Graphics/Figures/wbs_outreach.png}
%DIFDELCMD <   %%%
%DIFDELCMD < \caption{%
{%DIFAUXCMD
\DIFdelFL{Work breakdown structure for the outreach tasks}}
  %DIFAUXCMD
%DIFDELCMD < \label{fig:wbs6}
%DIFDELCMD < \end{figure}
%DIFDELCMD < 

%DIFDELCMD < %%%
\DIFdelend \section{Work Breakdown Structure (WBS)}

\DIFdelbegin \DIFdel{In figures \ref{fig:wbs1}, \ref{fig:wbs2}, \ref{fig:wbs3}, \ref{fig:wbs4}, \ref{fig:wbs5} and \ref{fig:wbs6} the \mbox{%DIFAUXCMD
\ac{wbs} }\hskip0pt%DIFAUXCMD
of the project is presented with work packages seperated under each responsible subteams }\DIFdelend \DIFaddbegin \DIFadd{The wbs is shown in \fref{fig:wbs}. It is divided into subteams with a "general" title for all system-wide tasks}\DIFaddend .

\DIFaddbegin \begin{center}
\includegraphics[width=0.8\textwidth]{0_images/figures/wbs_general.png}

\includegraphics[width=0.8\textwidth]{0_images/figures/wbs_manegerial.png}

\includegraphics[width=0.8\textwidth]{0_images/figures/wbs_science.png}

\includegraphics[width=0.8\textwidth]{0_images/figures/wbs_mechanical.png}

\includegraphics[width=0.8\textwidth]{0_images/figures/wbs_electrical.png}

\includegraphics[width=0.8\textwidth]{0_images/figures/wbs_software.png}

\includegraphics[width=0.8\textwidth]{0_images/figures/wbs_outreach.png}

\captionof{figure}{WBS for all subsystems and general tasks}
\label{fig:wbs}
\end{center}

%DIF >  %%%%%%%%%%%%%%%%%%%%%%%%%%%%%%%%%%%%%%%%%%%%%%%%%%%%%%%%%%%%%%%%%%%%%%%%%%%%%%%
\DIFaddend \section{Schedule}

\DIFdelbegin %DIFDELCMD < \begin{figure}
%DIFDELCMD <   \centering
%DIFDELCMD <   \includegraphics[width=\linewidth]{Graphics/Figures/gantt_chart.jpg}
%DIFDELCMD <   %%%
%DIFDELCMD < \caption{%
{%DIFAUXCMD
\DIFdelFL{Gannt Chart}}
  %DIFAUXCMD
%DIFDELCMD < \label{fig:gannt}
%DIFDELCMD <  \end{figure}
%DIFDELCMD < %%%
\DIFdelend \DIFaddbegin \DIFadd{A Gantt chart has been made to visualize the work to be done and the duration of specific tasks. It was divided into the different subteams in addition to a system wide \mbox{%DIFAUXCMD
\gls{ait}}\hskip0pt%DIFAUXCMD
. The Gantt chart is shown in \fref{fig:gantt}, with the critical path highlighted in red.
}\DIFaddend 

\DIFdelbegin \DIFdel{The Gannt chart in \fref{fig:gannt} shows the general schedule of the project}\DIFdelend \DIFaddbegin \DIFadd{The critical path is Define Scientific Objectives - Define Requirements - Order Materials - Build Pressurized Chamber. This demonstrates the importance of the pressurized chamber on the whole experiment and thus the project planning.
}

\DIFadd{\mbox{%DIFAUXCMD
\gls{mirage} }\hskip0pt%DIFAUXCMD
team members have declared their availability per week throughout the project duration, so that certain people were assigned different tasks.
}

\DIFadd{OnlineGantt was used to manage scheduling and assigning the tasks, whereas Discord is used for communications within the team}\DIFaddend .
\DIFdelbegin \DIFdel{Due to several dates being uncertain yet, it was approximated accross a longer period.
The color coding represents the responsible subteam.
}\DIFdelend 

\DIFaddbegin \begin{figure*}[htbp]
    \centering
    \includegraphics[height=22cm]{0_images/figures/gantt.png}
    \caption{\DIFaddFL{Gantt chart for the MIRAGE experiment. Critical path is hightlighted in red.}}
    \label{fig:gantt}
\end{figure*}

%DIF >  %%%%%%%%%%%%%%%%%%%%%%%%%%%%%%%%%%%%%%%%%%%%%%%%%%%%%%%%%%%%%%%%%%%%%%%%%%%%%%%
\DIFaddend \section{Resources}

\DIFdelbegin \DIFdel{\mbox{%DIFAUXCMD
\ac{mirage} }\hskip0pt%DIFAUXCMD
}\DIFdelend \DIFaddbegin \DIFadd{\mbox{%DIFAUXCMD
\gls{mirage} }\hskip0pt%DIFAUXCMD
}\DIFaddend is a project conducted by the RiseUpp association, which is under the umbrella of Astronomisk Ungdom. The RiseUpp association has a dedicated outreach team working on finding sponsors. Currently \DIFdelbegin \DIFdel{\mbox{%DIFAUXCMD
\ac{mirage} }\hskip0pt%DIFAUXCMD
}\DIFdelend \DIFaddbegin \DIFadd{\mbox{%DIFAUXCMD
\gls{mirage} }\hskip0pt%DIFAUXCMD
}\DIFaddend is the only project RiseUpp is working on, therefore a majority of this budget will go to \DIFdelbegin \DIFdel{\mbox{%DIFAUXCMD
\ac{mirage}}\hskip0pt%DIFAUXCMD
}\DIFdelend \DIFaddbegin \DIFadd{\mbox{%DIFAUXCMD
\gls{mirage}}\hskip0pt%DIFAUXCMD
}\DIFaddend .

As part of Astronomisk Ungdom, RiseUpp get annually 2000 SEK and 20 SEK per member. RiseUpp aims to be an association at Uppsala university that gathers all students interested in space, hosts events such as lunch lectures, and participate in programmes like Rexus Bexus.

The \DIFdelbegin \DIFdel{\mbox{%DIFAUXCMD
\ac{mirage} }\hskip0pt%DIFAUXCMD
}\DIFdelend \DIFaddbegin \DIFadd{\mbox{%DIFAUXCMD
\gls{mirage} }\hskip0pt%DIFAUXCMD
}\DIFaddend project will consist mostly of students working on their thesis, therefore Uppsala university \DIFdelbegin \DIFdel{provides }\DIFdelend \DIFaddbegin \DIFadd{covers costs up to }\DIFaddend 1000 SEK per student. In addition, through \DIFdelbegin \DIFdel{our }\DIFdelend \DIFaddbegin \DIFadd{the team's }\DIFaddend supervisor, Uwe Zimmermann, \DIFdelbegin \DIFdel{we have }\DIFdelend access to a laboratory with the necessary tools and a vacuum chamber for testing \DIFaddbegin \DIFadd{is given}\DIFaddend . Uppsala university also hosts the Swedish Institute for Space Physics (Institutet för Rymdfysik), which has a thermal chamber that \DIFdelbegin \DIFdel{we will be using }\DIFdelend \DIFaddbegin \DIFadd{will be used }\DIFaddend for testing.

RiseUpp has an agreement with the Swedish company \DIFdelbegin \DIFdel{, SenseAir }\DIFdelend \DIFaddbegin \DIFadd{SenseAir AB}\DIFaddend , in which SenseAir \DIFdelbegin \DIFdel{has provided }\DIFdelend \DIFaddbegin \DIFadd{AB has provided a brand new K96 sensor, a used K96 sensor for testing, \mbox{%DIFAUXCMD
\gls{ir} }\hskip0pt%DIFAUXCMD
lamps (which are a critical part of }\DIFaddend the K96 sensor\DIFaddbegin \DIFadd{) }\DIFaddend and will provide extra gas sensors if needed.

\DIFaddbegin \DIFadd{In addition, the sensor calibrations will be done in collaboration with the Norunda weather stations; who will provide gas samples and a Picarro gas analyzer for reference.
}

\DIFadd{The gas pumps and the compressor are purchased and received, which is another crucial component to the mission.
}

%DIF >  %%%%%%%%%%%%%%%%%%%%%%%%%%%%%%%%%%%%%%%%%%%%%%%%%%%%%%%%%%%%%%%%%%%%%%%%%%%%%%%
\DIFaddend \subsection{Workforce}

The workforce consists of voluntary and thesis students. Voluntary students spend around 10 hours a week and will be the main workforce until March, they do this voluntarily and will be the core team members. Thesis students \DIFdelbegin \DIFdel{will be }\DIFdelend \DIFaddbegin \DIFadd{are }\DIFaddend recruited to begin working on the project from March, they will work in project groups under the supervision of \DIFdelbegin \DIFdel{our }\DIFdelend \DIFaddbegin \DIFadd{the team's }\DIFaddend supervisor, Associate Professor Uwe Zimmermann, and will get credited for the work. The thesis \DIFaddbegin \DIFadd{project }\DIFaddend will be 15 credits, and the students are expected to work full-time on the project \DIFaddbegin \DIFadd{from 23rd of March until the 5th of June}\DIFaddend .

\DIFaddbegin \DIFadd{More information about the thesis projects is shown in \tref{tab:thesis}. Thesis students will not directly write into the \mbox{%DIFAUXCMD
\gls{sed} }\hskip0pt%DIFAUXCMD
but they will write their own university reports.
}

\DIFadd{One task for the voluntary students is to mentor and follow up with the thesis students. The voluntary students cover all subteams and will act as the interface between the thesis projects and the \mbox{%DIFAUXCMD
\gls{mirage} }\hskip0pt%DIFAUXCMD
project.
}

\DIFadd{There are weekly meetings for the whole team in which the core team members will follow up closely with the thesis projects and update on the progress, in addition to an extra subteam meeting each week to work on given tasks and plan the next steps. Thesis students will have the chance, and be encouraged, to continue working on the project voluntarily.
}

\begin{table}[htbp]
\centering
\caption{\DIFaddFL{Thesis Projects Overview}}
\label{tab:thesis}
\begin{tabular}{|l|L{3.5cm}|L{5.5cm}|}
\hline
\textbf{Project Number} & \textbf{Project Name} & \textbf{Project Description} \\ \hline

1 & Finite Element Analysis (FEA) and Computational Fluid Dynamics (CFD) 
& Finite Element Analysis (FEA) of temperature and stress \& Computational Fluid Dynamics (CFD) to analyze airflow. 
\\ \hline

2 & PCB Design and Embedded Programming 
& Printed circuit board (PCB) design, component selection, and embedded programming. 
\\ \hline

3 & Pressurized Gas Sensing Compartment 
& Design and assembly of a pressurized gas sensing experiment compartment. 
 \\ \hline

4 & Data Analysis Pipeline 
& Data analysis pipeline programming. 
 \\ \hline

5 & NDIR Sensor Testing and Calibration 
& Sensor testing and calibration of an NDIR gas sensor for stratospheric pressure and temperature conditions. 
 \\ \hline

6 & Thermal System Design 
& Thermal system design. 
 \\ \hline

7 & ECSS-Based Material Selection 
& Material selections in accordance with ECSS guidelines for a leakage-free and heat-resistant compartment. 
 \\ \hline

8 & Power Distribution Board Design 
& Power distribution board design. 
\\ \hline

\end{tabular}
\end{table}

\DIFaddend In addition to the \DIFdelbegin \DIFdel{\mbox{%DIFAUXCMD
\ac{mirage} }\hskip0pt%DIFAUXCMD
}\DIFdelend \DIFaddbegin \DIFadd{\mbox{%DIFAUXCMD
\gls{mirage} }\hskip0pt%DIFAUXCMD
}\DIFaddend workforce, there is the RiseUpp board which will ensure the consistency of the team and assemble a team for the next Rexus Bexus cycles.
\DIFaddbegin \clearpage
\DIFaddend 

%DIF >  %%%%%%%%%%%%%%%%%%%%%%%%%%%%%%%%%%%%%%%%%%%%%%%%%%%%%%%%%%%%%%%%%%%%%%%%%%%%%%%
\subsection{Budget}

The estimated total budget of the experiment \DIFdelbegin \textit{\DIFdel{parts}} %DIFAUXCMD
\DIFdel{only with a 10}\DIFdelend \DIFaddbegin \emph{\DIFadd{parts}} \DIFadd{without the K96 sensor is 6089 SEK, and with a 15}\DIFaddend \% buffer is \DIFdelbegin \DIFdel{9500 }\DIFdelend \DIFaddbegin \DIFadd{7000 }\DIFaddend SEK ($\approx$ \DIFdelbegin \DIFdel{900 }\DIFdelend \DIFaddbegin \DIFadd{650 }\DIFaddend EUR). \DIFdelbegin \DIFdel{However, the }\DIFdelend \DIFaddbegin \DIFadd{The }\DIFaddend K96 sensor, which \DIFdelbegin \DIFdel{we got }\DIFdelend \DIFaddbegin \DIFadd{was }\DIFaddend sponsored, is not commercially available and was estimated to be 5000 SEK\DIFaddbegin \DIFadd{; making the total cost of the experiment around 1200 EUR}\DIFaddend . Additional managerial costs like outreach and travel expenses lead to a sum of approximately \DIFdelbegin \DIFdel{90000 }\DIFdelend \DIFaddbegin \DIFadd{85000 }\DIFaddend SEK ($\approx$ \DIFdelbegin \DIFdel{8500 }\DIFdelend \DIFaddbegin \DIFadd{8000 }\DIFaddend EUR).

\DIFdelbegin %DIFDELCMD < \begin{longtable}{p{3cm}p{1.7cm}p{2cm}p{1.5cm}p{5cm}}
%DIFDELCMD < %%%
\DIFdelend \DIFaddbegin \DIFadd{For the hardware, \mbox{%DIFAUXCMD
\gls{mirage} }\hskip0pt%DIFAUXCMD
was sponsored by SenseAir AB and received 2 K96 sensors; one for testing and one for use. Additionally, the pumps and the compressor were ordered.
}

\DIFadd{All of the hardware costs can be covered by the money given through Uppsala university for the thesis projects, however \mbox{%DIFAUXCMD
\gls{mirage} }\hskip0pt%DIFAUXCMD
aims to find external sponsors for additional outreach budget. Additional material such as cables, screws, RTCs and small microcontrollers are available for the team in the designated lab at Uppsala university and is provided; they were not added to the total cost.
}

\begin{longtable}{p{3cm}p{1.7cm}p{1.5cm}p{3cm}p{5cm}}
\DIFaddend \caption{Budget Breakdown}
\label{tab:budget}\\
\DIFaddbegin 

\DIFaddend \toprule
\textbf{Cost Name} & \textbf{Quantity} &
\DIFdelbegin \textbf{\DIFdel{Category}} %DIFAUXCMD
%DIFDELCMD < &
%DIFDELCMD < %%%
\DIFdelend \textbf{Unit Price (SEK)} &  \DIFaddbegin \textbf{\DIFadd{Status}} & \DIFaddend \textbf{Description} \\
\midrule
\endfirsthead
\DIFaddbegin 

\DIFaddend \multicolumn{5}{c}{\tablename\ \thetable\ -- \textit{Continued from previous page}}\\
\toprule
\textbf{Cost Name} & \textbf{Quantity} &
\DIFdelbegin \textbf{\DIFdel{Category}} %DIFAUXCMD
%DIFDELCMD < &
%DIFDELCMD < %%%
\DIFdelend \textbf{Unit Price (SEK)} &  \DIFaddbegin \textbf{\DIFadd{Status}} & \DIFaddend \textbf{Description} \\
\midrule
\endhead
\DIFaddbegin 

\DIFaddend \midrule
\multicolumn{5}{r}{\textit{Continued on next page}}\\
\endfoot
\DIFaddbegin 

\DIFaddend \bottomrule
\endlastfoot

\DIFdelbegin \DIFdel{Aluminum Block }\DIFdelend \DIFaddbegin \DIFadd{Enclosure }\DIFaddend & 1 & \DIFdelbegin \DIFdel{Mechanical }\DIFdelend \DIFaddbegin \DIFadd{720 }\DIFaddend & \DIFdelbegin \DIFdel{200 }\DIFdelend \DIFaddbegin \DIFadd{Not ordered }\DIFaddend &  \DIFdelbegin \DIFdel{For enclosure }\DIFdelend \\
\DIFdelbegin \DIFdel{Plexy Glass }%DIFDELCMD < & %%%
\DIFdel{1 }%DIFDELCMD < & %%%
\DIFdel{Mechanical }%DIFDELCMD < & %%%
\DIFdel{90 }%DIFDELCMD < & %%%
\DIFdel{For lid }%DIFDELCMD < \\
%DIFDELCMD < %%%
\DIFdel{Screws }%DIFDELCMD < & %%%
\DIFdel{8 }%DIFDELCMD < & %%%
\DIFdel{Mechanical }%DIFDELCMD < & %%%
\DIFdel{1 }%DIFDELCMD < & \\
%DIFDELCMD < %%%
\DIFdel{High Precision Valve }%DIFDELCMD < & %%%
\DIFdel{1 }%DIFDELCMD < & %%%
\DIFdel{Mechanical }%DIFDELCMD < & %%%
\DIFdel{1000 }%DIFDELCMD < & \\
%DIFDELCMD < %%%
\DIFdel{Cables }%DIFDELCMD < & %%%
\DIFdel{20 }%DIFDELCMD < & %%%
\DIFdel{Electrical }%DIFDELCMD < & %%%
\DIFdel{0.2 }%DIFDELCMD < & \\
%DIFDELCMD < %%%
\DIFdelend MS5607 Pressure Sensor & \DIFdelbegin \DIFdel{1 }\DIFdelend \DIFaddbegin \DIFadd{2 }\DIFaddend & \DIFdelbegin \DIFdel{Electrical }%DIFDELCMD < & %%%
\DIFdelend 81.75 & \DIFaddbegin \DIFadd{Not ordered }& \DIFaddend \\
\DIFaddbegin \DIFadd{SHT45 }\DIFaddend Temperature Sensor & \DIFdelbegin \DIFdel{1 }\DIFdelend \DIFaddbegin \DIFadd{3 }\DIFaddend & \DIFdelbegin \DIFdel{Electrical }\DIFdelend \DIFaddbegin \DIFadd{51.88 }\DIFaddend & \DIFdelbegin \DIFdel{93.74 }\DIFdelend \DIFaddbegin \DIFadd{Not ordered }\DIFaddend & \\
\DIFdelbegin \DIFdel{SHT3x-DIS Humidity Sensor }\DIFdelend \DIFaddbegin \DIFadd{Voltage Regulators }\DIFaddend & \DIFdelbegin \DIFdel{1 }\DIFdelend \DIFaddbegin \DIFadd{3 }\DIFaddend & \DIFdelbegin \DIFdel{Electrical }\DIFdelend \DIFaddbegin \DIFadd{65.29 }\DIFaddend & \DIFdelbegin \DIFdel{103.01 }\DIFdelend \DIFaddbegin \DIFadd{Not ordered }\DIFaddend & \\
\DIFdelbegin \DIFdel{Arduino Ethernet Shield }\DIFdelend \DIFaddbegin \DIFadd{SD Card }\DIFaddend & 1 & \DIFdelbegin \DIFdel{Electrical }\DIFdelend \DIFaddbegin \DIFadd{50.04 }\DIFaddend & \DIFdelbegin \DIFdel{369 }\DIFdelend \DIFaddbegin \DIFadd{Not ordered }\DIFaddend & \DIFaddbegin \DIFadd{16 GB }\DIFaddend \\
\DIFdelbegin \DIFdel{Arduino Mega }\DIFdelend \DIFaddbegin \DIFadd{\mbox{%DIFAUXCMD
\gls{pcb} }\hskip0pt%DIFAUXCMD
}\DIFaddend & 1 & \DIFdelbegin \DIFdel{Electrical }\DIFdelend \DIFaddbegin \DIFadd{200 }\DIFaddend & \DIFdelbegin \DIFdel{579 }\DIFdelend \DIFaddbegin \DIFadd{Not ordered }\DIFaddend & \\
\DIFdelbegin \DIFdel{Voltage Regulators }\DIFdelend \DIFaddbegin \DIFadd{Vacuum Pump }\DIFaddend & \DIFdelbegin \DIFdel{3 }\DIFdelend \DIFaddbegin \DIFadd{2 }\DIFaddend & \DIFdelbegin \DIFdel{Electrical }\DIFdelend \DIFaddbegin \DIFadd{1353.91 }\DIFaddend & \DIFdelbegin \DIFdel{60 }\DIFdelend \DIFaddbegin \DIFadd{Ordered }\DIFaddend & \\
\DIFdelbegin \DIFdel{SD Card }\DIFdelend \DIFaddbegin \DIFadd{Compressor }\DIFaddend & 1 & \DIFdelbegin \DIFdel{Electrical }\DIFdelend \DIFaddbegin \DIFadd{1169.77 }\DIFaddend & \DIFdelbegin \DIFdel{119 }\DIFdelend \DIFaddbegin \DIFadd{Ordered }\DIFaddend & \DIFdelbegin \DIFdel{32 GB }\DIFdelend \\
\DIFdelbegin \DIFdel{\mbox{%DIFAUXCMD
\ac{pcb} }\hskip0pt%DIFAUXCMD
}\DIFdelend \DIFaddbegin \DIFadd{Pump DevBoard }\DIFaddend & 1 & \DIFdelbegin \DIFdel{Electrical }\DIFdelend \DIFaddbegin \DIFadd{53.96 }\DIFaddend & \DIFdelbegin \DIFdel{200 }\DIFdelend \DIFaddbegin \DIFadd{Ordered }\DIFaddend & \\
\DIFdelbegin \DIFdel{Pump }\DIFdelend \DIFaddbegin \DIFadd{Power connector }\DIFaddend & 1 & \DIFdelbegin \DIFdel{Electrical }\DIFdelend \DIFaddbegin \DIFadd{80.55 }\DIFaddend & \DIFdelbegin \DIFdel{500 }\DIFdelend \DIFaddbegin \DIFadd{Not ordered }\DIFaddend & \\
\DIFdelbegin \DIFdel{Thermal System }\DIFdelend \DIFaddbegin \DIFadd{Amphenol Data Connector (RJF21B) }\DIFaddend & 1 & \DIFdelbegin \DIFdel{Electrical }\DIFdelend \DIFaddbegin \DIFadd{571.25 }\DIFaddend & \DIFdelbegin \DIFdel{271.3 }\DIFdelend \DIFaddbegin \DIFadd{Not ordered }\DIFaddend & \\
\DIFdelbegin \DIFdel{Capacitor }\DIFdelend \DIFaddbegin \DIFadd{10 W heating pad }\DIFaddend & \DIFdelbegin \DIFdel{3 }\DIFdelend \DIFaddbegin \DIFadd{1 }\DIFaddend & \DIFdelbegin \DIFdel{Electrical }\DIFdelend \DIFaddbegin \DIFadd{20.63 }\DIFaddend & \DIFdelbegin \DIFdel{1 }\DIFdelend \DIFaddbegin \DIFadd{Not ordered }\DIFaddend & \\
K96 & \DIFdelbegin \DIFdel{1 }\DIFdelend \DIFaddbegin \DIFadd{2 }\DIFaddend & \DIFdelbegin \DIFdel{Electrical }%DIFDELCMD < & %%%
\DIFdelend 5000 & \DIFaddbegin \DIFadd{Available }& \DIFaddend Non-commercial, estimate of value\\
\DIFaddbegin \hline
\DIFadd{15\% Margin for Hardware }& & \DIFadd{913 }& & \\
\DIFadd{Total for Hardware }& & \DIFadd{12000 }& & \\
\hline
\DIFaddend Travel Expenses & 5 & \DIFdelbegin \DIFdel{Managerial }%DIFDELCMD < & %%%
\DIFdelend 3000 & \DIFaddbegin \DIFadd{Not ordered }& \DIFaddend For the launch campaign, counting only additional nonsponsored students \\
Accommodation & 5 & \DIFdelbegin \DIFdel{Managerial }%DIFDELCMD < & %%%
\DIFdelend 10000 & \DIFaddbegin \DIFadd{Not ordered }& \DIFaddend For the launch campaign, counting only additional nonsponsored students \\
Team Building & 15 & \DIFdelbegin \DIFdel{Managerial }%DIFDELCMD < & %%%
\DIFdelend 500 & \DIFaddbegin \DIFadd{Not ordered }& \DIFaddend \\
Team Jersey & 10 & \DIFdelbegin \DIFdel{Managerial }%DIFDELCMD < & %%%
\DIFdelend 100 & \DIFaddbegin \DIFadd{Not ordered }& \DIFaddend \\
\DIFaddbegin \hline
\DIFadd{15\% Margin for Manegerial Costs }& & \DIFadd{11025 }& & \\
\DIFadd{Total for Manegerial Costs }& & \DIFadd{84500 }& & \\
\DIFaddend \end{longtable}


%
%\begin{table}[H]
%\centering
%\caption{Budget Breakdown}
%\label{tab:budget}
%\begin{tabularx}{\textwidth}{p{3cm}p{1.7cm}p{2cm}p{1cm}X}
%\toprule
%\textbf{Cost Name} & \textbf{Quantity} & \textbf{Category} & \textbf{Unit Price (SEK)} & \textbf{Description} \\
%\midrule
%Aluminum Block & 1 & Mechanical & 200 & For enclosure \\
%Plexy Glass & 1 & Mechanical & 90 & For lid \\
%Screws & 8 & Mechanical & 1 & \\
%High Precision Valve & 1 & Mechanical & 1000 & \\
%Cables & 20 & Electrical & 0.2 & \\
%MS5607 Pressure Sensor & 1 & Electrical & 81.75 & \\
%Temperature Sensor & 1 & Electrical & 93.74 & \\
%SHT3x-DIS Humidity Sensor & 1 & Electrical & 103.01 & \\
%Arduino Ethernet Shield & 1 & Electrical & 369 & \\
%Arduino Mega & 1 & Electrical & 579 & \\
%Voltage Regulators & 3 & Electrical & 60 & \\
%SD Card & 1 & Electrical & 119 & 32 GB \\
%DIF < \ac{pcb} & 1 & Electrical & 200 & \\
%DIF > \gls{pcb} & 1 & Electrical & 200 & \\
%Pump & 1 & Electrical & 500 & \\
%Thermal System & 1 & Electrical & 271.3 & \\
%Capacitor & 3 & Electrical & 1 & \\
%K96 & 1 & Electrical & 5000 & Non-commercial, estimate of value \\
%Travel Expenses & 5 & Managerial & 3000 & For the launch campaign, counting only additional nonsponsored students \\
%Accomodation & 5 & Managerial & 10000 & For the launch campaign, counting only additional nonsponsored students \\
%Team Building & 15 & Managerial & 500 & \\
%Team Jersey & 10 & Managerial & 100 & \\
%\bottomrule
%\end{tabularx}
%\end{table}

%DIF >  %%%%%%%%%%%%%%%%%%%%%%%%%%%%%%%%%%%%%%%%%%%%%%%%%%%%%%%%%%%%%%%%%%%%%%%%%%%%%%%
\subsection{External Support}

\DIFdelbegin \DIFdel{\mbox{%DIFAUXCMD
\ac{mirage} }\hskip0pt%DIFAUXCMD
}\DIFdelend \DIFaddbegin \DIFadd{\mbox{%DIFAUXCMD
\gls{mirage} }\hskip0pt%DIFAUXCMD
}\DIFaddend is endorsed and adviced by the Uppsala university department of meteorology. In addition, \DIFdelbegin \DIFdel{\mbox{%DIFAUXCMD
\ac{mirage} }\hskip0pt%DIFAUXCMD
}\DIFdelend \DIFaddbegin \DIFadd{\mbox{%DIFAUXCMD
\gls{mirage} }\hskip0pt%DIFAUXCMD
}\DIFaddend is in the process of becoming an associated association to the student union for science and engineering at Uppsala university (UTN). This will open up the doors to apply for research grants provided by the UTN. Astronomisk Ungdom has a research grant pool of 100000 SEK which is available for associations under their umbrella, which will be applied to.

\DIFdelbegin \DIFdel{We are }\DIFdelend \DIFaddbegin \DIFadd{\mbox{%DIFAUXCMD
\gls{mirage} }\hskip0pt%DIFAUXCMD
is }\DIFaddend in talks with Uppsala University Innovation and the development office at the university about financial and outreach support. They will connect \DIFdelbegin \DIFdel{us }\DIFdelend \DIFaddbegin \DIFadd{\mbox{%DIFAUXCMD
\gls{mirage} }\hskip0pt%DIFAUXCMD
}\DIFaddend with companies that work in the field and help \DIFdelbegin \DIFdel{us }\DIFdelend with social events organized by RiseUpp.

\DIFdelbegin \section{\DIFdel{Outreach Approach}}
%DIFAUXCMD
\addtocounter{section}{-1}%DIFAUXCMD
\DIFdelend %DIF >  %%%%%%%%%%%%%%%%%%%%%%%%%%%%%%%%%%%%%%%%%%%%%%%%%%%%%%%%%%%%%%%%%%%%%%%%%%%%%%%

\DIFdelbegin \DIFdel{The \mbox{%DIFAUXCMD
\ac{mirage} }\hskip0pt%DIFAUXCMD
RexusBexus team will contact local news agencies, such as Uppsala Nya Tidning, Techna (student university magazine) and appear on Student radio (local independent radiostation).
The team also aims to apear on Uppsala University website. 
The RiseUpp association is under the Astronomisk Ungdom association and a website will be ran under their domain. A project documentary movie will follow the Rexus }\DIFdelend \DIFaddbegin \section{\DIFadd{Outreach Approach}}
\subsection{\DIFadd{Goal of the outreach}}
\DIFadd{The aim is to spread the word of the Rexus}\DIFaddend /Bexus \DIFdelbegin \DIFdel{team as they work on the project and will document the flight. }\DIFdelend \DIFaddbegin \DIFadd{programme and MIRAGE to as large an audiance as possible. Especially inspiring STEM majors at Uppsala university to engage in MIRAGE or similar Rexus/Bexus projects and experiments in the future.
}\vspace{3mm}\\
\DIFadd{The MIRAGE team has advertised themselves by posting continuous updates on social media about the experiment, team meetings, team building activities and humorous content. They have also written an informative letter about the Rexus Bexus programme, the RiseUpp association and MIRAGE that has been announced on several engineering program's studium pages for recruitment of bachelor thesis students.
}\vspace{3mm}\\
\DIFadd{Other media sources that the team has advertised themselves via, are newspapers. An article was written in the students university´s magazine Techna.
}\vspace{3mm}\\
\DIFadd{To cover equipment costs for the experiment and travel expenses the team has contacted local and national STEM-companies, authorities and foundations to ask for sponsorships and collaborations. One foundation that has been contacted was Sparbanksstiftelsen Upland, Swedbank and the team have an ongoing collaboration with RedBull.
}\DIFaddend 

\begin{table}[H]
\centering
\caption{Team social media links}
\label{tab:social_media}
\DIFdelbeginFL %DIFDELCMD < \begin{tabular}{llll}
%DIFDELCMD < \toprule
%DIFDELCMD < \textbf{Team} & \textbf{Instagram} & \textbf{LinkedIn} & \textbf{Website} \\
%DIFDELCMD < \midrule
%DIFDELCMD < Mirage & \href{https://www.instagram.com/riseupp_uu/?igsh=MXhsM3lsZmMyMjh4ZA\%3D\%3D#}{\text{Riseupp\_uu}} & \href{https://www.linkedin.com/company/riseupp-uu/posts/?feedView=all}{riseupp-uu}& Coming soon   \\
%DIFDELCMD < \bottomrule
%DIFDELCMD < \end{tabular}
%DIFDELCMD < %%%
\DIFdelendFL \DIFaddbeginFL \begin{tabular}{|c|c|c|c|}\hline
\textbf{Team} & \textbf{Instagram} & \textbf{LinkedIn} & \textbf{Website} \\\hline
Mirage & \href{https://www.instagram.com/riseupp_uu/?igsh=MXhsM3lsZmMyMjh4ZA\%3D\%3D#}{\text{riseupp\_uu}} & \href{https://www.linkedin.com/company/riseupp-uu/posts/?feedView=all}{riseupp-uu}&  https://www.spacestudents.se/riseupp/   \\\hline
\end{tabular}
\DIFaddendFL \end{table}

\DIFaddbegin \subsection{\DIFadd{Future outreach plans}}
\DIFadd{In the future, the team plans to arrange lunch lectures with SSC, Christer Fuglesang, Marcus Wandt. These events can be sponsored by Uppsala kommun and supported by RedBull to attract a greater audience.
}\vspace{3mm}\\
\DIFadd{The team will continue to stay in touch with companies, authorities and foundations as well as contacting new potential sponsors.
}\vspace{3mm}\\
\DIFadd{The team wants to reach out to other media sources such as Uppsala Nya Tidning, Sveriges Radio and Uppsala university's student website/social media.
}\vspace{3mm}\\
\DIFadd{In a later stage of the project, a documentary movie will be released that follows the team´s project journey and the launch from Kiruna in October 2026.
}



\DIFaddend \section{Risk Register}

\textbf{Risk ID}

TC -- technical/implementation

MS -- mission (operational performance)

SF -- safety

VE -- vehicle

PE -- personnel

EN -- environmental

OR -- organisational

DS -- data related

\textbf{Probability (P)}

A. Minimum -- Almost impossible to occur

B. Low -- Small chance to occur

C. Medium -- Reasonable chance to occur

D. High -- Quite likely to occur

E. Maximum -- Certain to occur, maybe more than once

\textbf{Severity (S)}

1. Negligible -- Minimal or no impact

2. Significant -- Leads to reduced experiment performance

3. Major -- Leads to failure of subsystem or loss of flight data

4. Critical -- Leads to experiment failure or creates minor health hazards

5. Catastrophic -- Leads to termination of the REXUS and/or BEXUS programme, damage to the vehicle or injury to personnel

The rankings for probability (P) and severity (S) are combined to assess the overall risk classification, ranging from very low to very high and being coloured green, yellow, orange or red according to the \DIFdelbegin \DIFdel{\mbox{%DIFAUXCMD
\ac{sed} }\hskip0pt%DIFAUXCMD
}\DIFdelend \DIFaddbegin \DIFadd{\mbox{%DIFAUXCMD
\gls{sed} }\hskip0pt%DIFAUXCMD
}\DIFaddend guidelines.

\begin{longtable}{L{1.5cm}L{3cm}L{1cm}L{1cm}L{2cm}L{6cm}}
\caption{Risk Register}
\label{tab:risk_register}\\

\toprule
\textbf{ID} & \textbf{Risk \& consequence} & \textbf{P} & \textbf{S} &
\textbf{P x S} & \textbf{Prevention and/or Mitigation Action} \\
\midrule
\endfirsthead

\multicolumn{6}{c}{\tablename\ \thetable\ -- \textit{Continued from previous page}}\\
\toprule
\textbf{ID} & \textbf{Risk \& consequence} & \textbf{P} & \textbf{S} &
\textbf{P x S} & \textbf{Prevention and/or Mitigation Action} \\
\midrule
\endhead

\midrule
\multicolumn{6}{r}{\textit{Continued on next page}}\\
\endfoot

\bottomrule
\endlastfoot

\rowcolor{mediumorange} TC10 & Pump failure & C & 4 & medium & \DIFdelbegin \DIFdel{have }\DIFdelend \DIFaddbegin \DIFadd{Have }\DIFaddend multiple pumps in setup, choose appropriate components \\
\DIFaddbegin \rowcolor{lowyellow} \DIFadd{TC11 }& \DIFadd{Pump failure }& \DIFadd{C }& \DIFadd{2 }& \DIFadd{low }&  \DIFadd{Acceptable risk }\\
\DIFaddend \rowcolor{mediumorange} MS10 & \DIFdelbegin \DIFdel{Over frozen }\DIFdelend \DIFaddbegin \DIFadd{Over-frozen }\DIFaddend gas inlet or outlet & C & 4 & medium & Large enough radius \DIFaddbegin \DIFadd{of outlets and inlets}\DIFaddend , heating element at inlet \DIFdelbegin \DIFdel{and warm }\DIFdelend \DIFaddbegin \DIFadd{to heat }\DIFaddend air to appropriate temperature\DIFaddbegin \DIFadd{, move primary expansion point of air to being inside the experiment enclosure }\DIFaddend \\
\DIFaddbegin \rowcolor{verylowgreen} \DIFadd{MS11 }& \DIFadd{Over-frozen gas inlet or outlet }& \DIFadd{A }& \DIFadd{4 }& \DIFadd{very low }& \DIFadd{Acceptable risk }\\
\DIFaddend \rowcolor{lowyellow} DS10 & Loss of data & B & 4 & low & Send part of data to ground regularly \\
\DIFaddbegin \rowcolor{verylowgreen} \DIFadd{DS11 }& \DIFadd{Loss of data }& \DIFadd{A }& \DIFadd{4 }& \DIFadd{very low }& \DIFadd{Acceptable risk }\\
\DIFaddend \rowcolor{mediumorange} TC20 & \DIFdelbegin \DIFdel{Over/}\DIFdelend Underheating of electronics \DIFaddbegin \DIFadd{during ascent }\DIFaddend & C & \DIFdelbegin \DIFdel{4 }\DIFdelend \DIFaddbegin \DIFadd{3 }\DIFaddend & medium & More heating \DIFdelbegin \DIFdel{/cooling measures depending on test results }\DIFdelend \DIFaddbegin \DIFadd{measures such as insulation and heating elements }\DIFaddend \\
\DIFaddbegin \rowcolor{verylowgreen} \DIFadd{TC21 }& \DIFadd{Underheating of electronics during ascent }& \DIFadd{A }& \DIFadd{4 }& \DIFadd{very low }& \DIFadd{Acceptable risk }\\
\DIFaddend \rowcolor{lowyellow} TC30 & Failure of thermal control \DIFaddbegin \DIFadd{due to too low temperatures }\DIFaddend & \DIFdelbegin \DIFdel{C }\DIFdelend \DIFaddbegin \DIFadd{B }\DIFaddend & 3 & low & \DIFdelbegin \DIFdel{Testing and simulations, more heating or cooling depending on results }\DIFdelend \DIFaddbegin \DIFadd{Acceptable risk}\DIFaddend \\
\DIFdelbegin %DIFDELCMD < \rowcolor{lowyellow} %%%
\DIFdelend \DIFaddbegin \rowcolor{verylowgreen} \DIFaddend PE10 & Drop out of Team member & \DIFdelbegin \DIFdel{C }\DIFdelend \DIFaddbegin \DIFadd{B }\DIFaddend & 2 & \DIFaddbegin \DIFadd{very }\DIFaddend low & \DIFdelbegin \DIFdel{Good communication between team members about workload and availability; additional recruiting to reduce workload; selection for team members with sufficient time resources during application }\DIFdelend \DIFaddbegin \DIFadd{Acceptable risk }\DIFaddend \\
\rowcolor{mediumorange} OR10 & Lack of funding & \DIFdelbegin \DIFdel{B }\DIFdelend \DIFaddbegin \DIFadd{C }\DIFaddend & \DIFdelbegin \DIFdel{5 }\DIFdelend \DIFaddbegin \DIFadd{4 }\DIFaddend & medium & \DIFdelbegin \DIFdel{Contacting of companiesin progress }\DIFdelend \DIFaddbegin \DIFadd{Contact more companies, find other sources from university }\DIFaddend \\
\rowcolor{lowyellow} \DIFaddbegin \DIFadd{OR11 }& \DIFadd{Lack of funding }& \DIFadd{B }& \DIFadd{4 }& \DIFadd{low }& \DIFadd{Acceptable risk  }\\
\rowcolor{lowyellow} \DIFaddend MS20 & \DIFdelbegin \DIFdel{\mbox{%DIFAUXCMD
\ac{mcu} }\hskip0pt%DIFAUXCMD
}\DIFdelend \DIFaddbegin \DIFadd{\mbox{%DIFAUXCMD
\gls{mcu} }\hskip0pt%DIFAUXCMD
}\DIFaddend failure & B & 4 & low & Test under pressure and temperature conditions while on ground \DIFaddbegin \DIFadd{and choose different component if it doesn't tolerate}\DIFaddend , use watchdog timer and consider redundant systems\\
\DIFdelbegin %DIFDELCMD < \rowcolor{lowyellow} %%%
\DIFdel{MS30 }%DIFDELCMD < & %%%
\DIFdel{Automatic switching of measurement times fails }%DIFDELCMD < & %%%
\DIFdel{C }%DIFDELCMD < & %%%
\DIFdel{2 }%DIFDELCMD < & %%%
\DIFdel{low }%DIFDELCMD < & %%%
\DIFdel{Add manual option via E-Link as backup }%DIFDELCMD < \\
%DIFDELCMD < %%%
\DIFdelend \rowcolor{verylowgreen} \DIFdelbegin \DIFdel{MS40 }\DIFdelend \DIFaddbegin \DIFadd{MS21 }\DIFaddend & \DIFdelbegin \DIFdel{Experiment doesn't respond to manual commands from ground }\DIFdelend \DIFaddbegin \DIFadd{\mbox{%DIFAUXCMD
\gls{mcu} }\hskip0pt%DIFAUXCMD
failure }\DIFaddend & \DIFdelbegin \DIFdel{C }\DIFdelend \DIFaddbegin \DIFadd{A }\DIFaddend & \DIFdelbegin \DIFdel{1 }\DIFdelend \DIFaddbegin \DIFadd{4 }\DIFaddend & \DIFdelbegin \DIFdel{Very }\DIFdelend \DIFaddbegin \DIFadd{very }\DIFaddend low & \DIFdelbegin \DIFdel{No action }\DIFdelend \DIFaddbegin \DIFadd{Acceptable risk}\DIFaddend \\
\rowcolor{verylowgreen} \DIFdelbegin \DIFdel{MS50 }\DIFdelend \DIFaddbegin \DIFadd{TC40 }\DIFaddend & \DIFdelbegin \DIFdel{Loss of communication with ground stations }\DIFdelend \DIFaddbegin \DIFadd{Outside sensors fail }\DIFaddend & \DIFdelbegin \DIFdel{C }\DIFdelend \DIFaddbegin \DIFadd{B }\DIFaddend & \DIFdelbegin \DIFdel{1 }\DIFdelend \DIFaddbegin \DIFadd{2 }\DIFaddend & \DIFdelbegin \DIFdel{Very }\DIFdelend \DIFaddbegin \DIFadd{very }\DIFaddend low & \DIFdelbegin \DIFdel{No action }\DIFdelend \DIFaddbegin \DIFadd{Acceptable risk }\DIFaddend \\
\rowcolor{lowyellow} \DIFdelbegin \DIFdel{TC40 }\DIFdelend \DIFaddbegin \DIFadd{TC50 }\DIFaddend & \DIFdelbegin \DIFdel{Outside sensors not working due to low temperatures }\DIFdelend \DIFaddbegin \DIFadd{Condensation of water inside the measurement chamber }\DIFaddend & C & 2 & low & \DIFdelbegin \DIFdel{Choose and test appropriate components }\DIFdelend \DIFaddbegin \DIFadd{Implement flushing cycle }\DIFaddend \\
\DIFdelbegin %DIFDELCMD < \rowcolor{mediumorange} %%%
\DIFdel{TC50 }\DIFdelend \DIFaddbegin \rowcolor{verylowgreen} \DIFadd{TC51 }\DIFaddend & Condensation of water inside the measurement chamber & \DIFdelbegin \DIFdel{C }\DIFdelend \DIFaddbegin \DIFadd{B }\DIFaddend & \DIFdelbegin \DIFdel{4 }\DIFdelend \DIFaddbegin \DIFadd{2 }\DIFaddend & \DIFdelbegin \DIFdel{medium }\DIFdelend \DIFaddbegin \DIFadd{very low }\DIFaddend & \DIFdelbegin \DIFdel{Test flushing cycle or find other solutions }\DIFdelend \DIFaddbegin \DIFadd{Acceptable risk }\DIFaddend \\
\DIFdelbegin %DIFDELCMD < \rowcolor{mediumorange} %%%
\DIFdelend \DIFaddbegin \rowcolor{lowyellow} \DIFaddend TC60 & Contamination due to materials of the balloon, gondola, and experiment setup & C & \DIFdelbegin \DIFdel{4 }\DIFdelend \DIFaddbegin \DIFadd{2 }\DIFaddend & \DIFdelbegin \DIFdel{medium }\DIFdelend \DIFaddbegin \DIFadd{low }\DIFaddend & \DIFdelbegin \DIFdel{Check materials used and outgasing potential}\DIFdelend \DIFaddbegin \DIFadd{Carefull selection of non-outgassing materials of the experiment and check with organisers for balloon/gondola materials; direct the inlet away from the ballon and other experiments}\DIFaddend \\
\rowcolor{verylowgreen} \DIFdelbegin \DIFdel{DS20 }\DIFdelend \DIFaddbegin \DIFadd{TC61 }\DIFaddend & \DIFdelbegin \DIFdel{Data storage overflow }\DIFdelend \DIFaddbegin \DIFadd{Contamination due to materials of the balloon, gondola, and experiment setup }\DIFaddend & A & \DIFdelbegin \DIFdel{4 }\DIFdelend \DIFaddbegin \DIFadd{2 }\DIFaddend & \DIFdelbegin \DIFdel{Very }\DIFdelend \DIFaddbegin \DIFadd{very }\DIFaddend low & \DIFdelbegin \DIFdel{SD card should be chosen large enough with enough margin }\DIFdelend \DIFaddbegin \DIFadd{Acceptable risk}\DIFaddend \\
\rowcolor{lowyellow} \DIFdelbegin \DIFdel{MS60 }\DIFdelend \DIFaddbegin \DIFadd{MS30 }\DIFaddend & Software failure \DIFdelbegin \DIFdel{for }\DIFdelend \DIFaddbegin \DIFadd{of }\DIFaddend thermal control & B & 3 & low & \DIFdelbegin \DIFdel{Option }\DIFdelend \DIFaddbegin \DIFadd{Implement a watchdog and an option }\DIFaddend to restart software from ground\DIFaddbegin \DIFadd{. Include thermal switch to disengage heating devices if temperatures are too hot. }\DIFaddend \\
\DIFaddbegin \rowcolor{verylowgreen} \DIFadd{MS31 }& \DIFadd{Software failure of thermal control }& \DIFadd{A }& \DIFadd{3 }& \DIFadd{very low }& \DIFadd{Acceptable risk }\\
\DIFaddend \rowcolor{lowyellow} \DIFdelbegin \DIFdel{MS70 }\DIFdelend \DIFaddbegin \DIFadd{MS40 }\DIFaddend & Software failure \DIFdelbegin \DIFdel{for }\DIFdelend \DIFaddbegin \DIFadd{of }\DIFaddend pressure control & \DIFdelbegin \DIFdel{B }\DIFdelend \DIFaddbegin \DIFadd{C }\DIFaddend & \DIFdelbegin \DIFdel{4 }\DIFdelend \DIFaddbegin \DIFadd{2 }\DIFaddend & low & \DIFdelbegin \DIFdel{Option }\DIFdelend \DIFaddbegin \DIFadd{Implement watchdog and option }\DIFaddend to restart software from ground \\
\DIFaddbegin \rowcolor{verylowgreen} \DIFadd{MS41 }& \DIFadd{Software failure of pressure control }& \DIFadd{A }& \DIFadd{2 }& \DIFadd{very low }& \DIFadd{Acceptable risk }\\
\DIFaddend \end{longtable}
%DIF < 
%DIF < \begin{landscape}
%DIF < \begin{table}[H]
%DIF < \centering
%DIF < \caption{Risk Register}
%DIF < \label{tab:risk_register}
%DIF < \begin{tabular}{p{1.5cm}p{4cm}p{1cm}p{1cm}p{2cm}p{6cm}}
%DIF < \toprule
%DIF < \textbf{ID} & \textbf{Risk \& consequence} & \textbf{P} & \textbf{S} & \textbf{P x S} & \textbf{Prevention and/or Mitigation Action} \\
%DIF < \midrule
%DIF < TC10 & Pump failure & C & 4 & medium & have multiple pumps in setup, chose appropriate components \\
%DIF < MS10 & Over frozen gas inlet or outlet & C & 4 & medium & Large enough radius, heat air to appropriate temperature \\
%DIF < DS10 & Loss of data & B & 4 & low & Send part of data to ground regularly \\
%DIF < TC20 & Over/Underheating of electronics & C & 4 & medium & More heating/cooling measures depending on test results \\
%DIF < TC30 & Failure of thermal control & C & 3 & low & Planning and testing \\
%DIF < PE10 & Drop out of Team member & C & 2 & low & Good communication between team members about workload and availability; additional recruiting to reduce workload; selection for team members with sufficient time resources during application \\
%DIF < OR10 & Lack of funding & B & 5 & medium & Contact companies soon \\
%DIF < MS20 & \ac{mcu} failure & B & 4 & low & Test under pressure and temperature conditions while on ground \\
%DIF < MS30 & Automatic switching of measurement times fails & C & 2 & low & Add manual option as backup \\
%DIF < MS40 & Experiment doesn't respond to manual commands from ground & C & 1 & Very low & No action \\
%DIF < MS50 & Loss of communication with ground stations & C & 1 & Very low & No action \\
%DIF < TC40 & Outside sensors not working due to low temperatures & C & 2 & low & Chose and test appropriate components \\
%DIF < TC50 & Condensation of water inside the \ac{ndir} sensor box. & C & 4 & medium & Add water removal system to design \\
%DIF < DS20 & Data storage overflow & A & 4 & Very low & SD card should be chosen large enough with enough margin \\
%DIF < MS60 & Software failure for thermal control & B & 3 & low & Option to restart software from ground \\
%DIF < MS70 & Software failure for pressure control & B & 4 & low & Option to restart software from ground \\
%DIF < \bottomrule
%DIF < \end{tabular}
%DIF < \end{table}
%DIF < \end{landscape}
%DIF < 
%DIF >  %
%DIF >  %\begin{landscape}
%DIF >  %\begin{table}[H]
%DIF >  %\centering
%DIF >  %\caption{Risk Register}
%DIF >  %\label{tab:risk_register}
%DIF >  %\begin{tabular}{p{1.5cm}p{4cm}p{1cm}p{1cm}p{2cm}p{6cm}}
%DIF >  %\toprule
%DIF >  %\textbf{ID} & \textbf{Risk \& consequence} & \textbf{P} & \textbf{S} & \textbf{P x S} & \textbf{Prevention and/or Mitigation Action} \\
%DIF >  %\midrule
%DIF >  %TC10 & Pump failure & C & 4 & medium & have multiple pumps in setup, chose appropriate components \\
%DIF >  %MS10 & Over frozen gas inlet or outlet & C & 4 & medium & Large enough radius, heat air to appropriate temperature \\
%DIF >  %DS10 & Loss of data & B & 4 & low & Send part of data to ground regularly \\
%DIF >  %TC20 & Over/Underheating of electronics & C & 4 & medium & More heating/cooling measures depending on test results \\
%DIF >  %TC30 & Failure of thermal control & C & 3 & low & Planning and testing \\
%DIF >  %PE10 & Drop out of Team member & C & 2 & low & Good communication between team members about workload and availability; additional recruiting to reduce workload; selection for team members with sufficient time resources during application \\
%DIF >  %OR10 & Lack of funding & B & 5 & medium & Contact companies soon \\
%DIF >  %MS20 & \gls{mcu} failure & B & 4 & low & Test under pressure and temperature conditions while on ground \\
%DIF >  %MS30 & Automatic switching of measurement times fails & C & 2 & low & Add manual option as backup \\
%DIF >  %MS40 & Experiment doesn't respond to manual commands from ground & C & 1 & Very low & No action \\
%DIF >  %MS50 & Loss of communication with ground stations & C & 1 & Very low & No action \\
%DIF >  %TC40 & Outside sensors not working due to low temperatures & C & 2 & low & Chose and test appropriate components \\
%DIF >  %TC50 & Condensation of water inside the \gls{ndir} sensor box. & C & 4 & medium & Add water removal system to design \\
%DIF >  %DS20 & Data storage overflow & A & 4 & Very low & SD card should be chosen large enough with enough margin \\
%DIF >  %MS60 & Software failure for thermal control & B & 3 & low & Option to restart software from ground \\
%DIF >  %MS70 & Software failure for pressure control & B & 4 & low & Option to restart software from ground \\
%DIF >  %\bottomrule
%DIF >  %\end{tabular}
%DIF >  %\end{table}
%DIF >  %\end{landscape}
%DIF >  %

\DIFdelbegin \chapter{\DIFdel{Experiment Description}}
%DIFAUXCMD
\addtocounter{chapter}{-1}%DIFAUXCMD
\DIFdelend %DIF >  %%%%%%%%%%%%%%%%%%%%%%%%%%%%%%%%%%%%%%%%%%%%%%%%%%%%%%%%%%%%%%%%%%%%%%%%%%%%%%%\newpage

\DIFaddbegin \chapter{\DIFadd{EXPERIMENT DESCRIPTION}}
\label{sec:experiment_description}
\DIFaddend \section{Experiment \DIFdelbegin \DIFdel{Setup}\DIFdelend \DIFaddbegin \DIFadd{setup}\DIFaddend }

\DIFdelbegin %DIFDELCMD < \begin{figure}[h]
%DIFDELCMD <     %%%
\DIFdelendFL \DIFaddbeginFL \begin{figure}[htbp]
    \DIFaddendFL \centering
    \DIFdelbeginFL %DIFDELCMD < \includegraphics[width=0.8\textwidth]{Graphics/Figures/Mirage Block schematic.png}
%DIFDELCMD <     %%%
\DIFdelendFL \DIFaddbeginFL \includegraphics[width=0.8\textwidth]{0_images/figures/Mirage Block schematic.png}
    \DIFaddendFL \caption{Experimental Setup}
    \label{fig:experimental_setup}

\end{figure}

The experiment will feature a NDIR sensor (K96), a pressurization system, a thermal system, ambient environment sensors and a processing unit. 
The pressurization system will include a pumping mechanism in order to increase the pressure and the amount of particles \DIFdelbegin \DIFdel{we can measure. The gas }\DIFdelend \DIFaddbegin \DIFadd{that can be measured. The air }\DIFaddend will be heated up and pumped into the measurement chamber with the K96 sensor inside together with temperature, pressure and humidity sensors. The gas will be exchanged continously to exchange the sample while maintaining the pressure of \DIFdelbegin \DIFdel{at least 2 }\DIFdelend \DIFaddbegin \DIFadd{3 }\DIFaddend bar. 
The K96 Sensor will communicate with the processing unit and the processing unit sends the measured digital data for storage to the SDHC card.
The pressure sensor,temperature sensor and humidity sensor communicate with the processing unit via digital interfaces (such as $I^2C$) to ensure that the pump and thermal control system are running with the right parameters to maintain room temperature conditions and a stable pressure \DIFdelbegin \DIFdel{between 2 and 5 }\DIFdelend \DIFaddbegin \DIFadd{of 3 }\DIFaddend bar. Unique addresses are used to identify each sensor's signal, allowing the processing unit to collect and organize data from all ambient sensors efficiently.

%DIF > %%%%%%%%%%%%%%%%%%%%%%%%%%%%%%%%%%%%%%%%%%%%%%%%%%%%%%%%%%%%%%%%%%%%%%%%%%%%%%
\section{Experiment Interfaces}

%DIF > %%%%%%%%%%%%%%%%%%%%%%%%%%%%%%%%%%%%%%%%%%%%%%%%%%%%%%%%%%%%%%%%%%%%%%%%%%%%%%
\subsection{Mechanical \DIFdelbegin %DIFDELCMD < \MBLOCKRIGHTBRACE
%DIFDELCMD < %%%
\DIFdel{The experiment is mounted }\DIFdelend \DIFaddbegin \DIFadd{Interface}}\label{mechanical_interface}

\DIFadd{The \mbox{%DIFAUXCMD
\gls{mirage} }\hskip0pt%DIFAUXCMD
enclosure is mechanically attached }\DIFaddend to the ESCARGO \DIFdelbegin \DIFdel{aluminium frame composed of }\DIFdelend \DIFaddbegin \DIFadd{gondola structure via four mounting points located on its bottom plate. These four connection points interface with two parallel }\DIFaddend 45 $\times$ 45 \DIFdelbegin \DIFdel{~mm Bosch-type }\DIFdelend \DIFaddbegin \DIFadd{mm Bosch aluminium }\DIFaddend T-slot profiles \DIFdelbegin \DIFdel{used by the BEXUS gondola.
}%DIFDELCMD < 

%DIFDELCMD < %%%
\DIFdel{The experiment is housed in a rigid aluminium enclosure and is fixed to a structural aluminium base plate using vertical }\DIFdelend \DIFaddbegin \DIFadd{using }\DIFaddend M6 \DIFdelbegin \DIFdel{bolts. The base plate forms the primary load-bearing interface of the experiment and ensures a rigid connection between the experiment enclosure and the gondola mounting structure. }%DIFDELCMD < 

%DIFDELCMD < %%%
\DIFdel{The aluminium base plate is attached }\DIFdelend \DIFaddbegin \DIFadd{T-bolts, M6 rubber spacers, and M6 flange nuts. The parallel profiles are then mounted }\DIFaddend to the ESCARGO \DIFdelbegin \DIFdel{frame via horizontal }\DIFdelend \DIFaddbegin \DIFadd{gondola structure using standard Rexroth }\DIFaddend M6 \DIFdelbegin \DIFdel{bolts engaging }\DIFdelend \DIFaddbegin \DIFadd{brackets, }\DIFaddend M6 \DIFdelbegin \DIFdel{T-nuts (or }\DIFdelend T-bolts\DIFdelbegin \DIFdel{) inserted into the ESCARGO T-slots. The horizontal bolt orientation allows direct fastening }\DIFdelend \DIFaddbegin \DIFadd{, and M6 flange nuts. This configuration facilitates installation without necessitating structural modifications }\DIFaddend to the ESCARGO \DIFdelbegin \DIFdel{profiles without any modification of the gondola structure.
Washers are used under all bolt heads to ensure proper load distribution}\DIFdelend \DIFaddbegin \DIFadd{frame, ensuring full compatibility with the standard mounting interface.
}

\DIFadd{To reduce vibration transmission, temperature transfer and accommodate minor alignment tolerances during installation, rubber spacers are installed between the enclosure and the mounting brackets}\DIFaddend .

\DIFdelbegin \DIFdel{A total of four independent mounting points are used between the base plate }\DIFdelend \DIFaddbegin \DIFadd{The mounting configuration provides a rigid mechanical interface between the experiment enclosure }\DIFaddend and the ESCARGO frame, \DIFdelbegin \DIFdel{providing redundancy and ensuring that the experiment remains securely attached in the event of a single fastener failure}\DIFdelend \DIFaddbegin \DIFadd{ensuring structural redundancy and compliance with the standard ESCARGO mechanical interface requirements}\DIFaddend .

\DIFdelbegin \DIFdel{No modifications to the ESCARGO frame are required. A schematic diagram of the mechanical interface between the experiment, the aluminium base plate, and the ESCARGO frame is provided.
}%DIFDELCMD < 

%DIFDELCMD < \begin{table}[h]
%DIFDELCMD < %%%
\DIFdelendFL \DIFaddbeginFL \begin{table}[htbp]
\DIFaddendFL \centering
\DIFdelbeginFL %DIFDELCMD < \caption{%
{%DIFAUXCMD
\DIFdelFL{Mechanical interface fasteners and mounting hardware}}
%DIFAUXCMD
%DIFDELCMD < \begin{tabular}{|l|c|l|}
%DIFDELCMD < \hline
%DIFDELCMD < \textbf{Item} & \textbf{Quantity} & \textbf{Specification} \\
%DIFDELCMD < \hline
%DIFDELCMD < Bolts (experiment $\rightarrow$ base plate) & 4 & M6, steel \\
%DIFDELCMD < Bolts (base plate $\rightarrow$ ESCARGO) & 4 & M6, steel \\
%DIFDELCMD < T-nuts / T-bolts & 4 & \makecell[l]{M6, compatible with 45 $\times$ 45~mm \\ ESCARGO T-slots} \\
%DIFDELCMD < Washers & 8 & M6 \\
%DIFDELCMD < Base plate & 1 & Aluminium, structural \\
%DIFDELCMD < \hline
%DIFDELCMD < \end{tabular}
%DIFDELCMD < %%%
\DIFdelendFL \DIFaddbeginFL \begin{tabular}{lcl}
\toprule
\textbf{Item} & \textbf{Quantity} & \textbf{Specification} \\
\midrule
L-bracket & 4 & Bosch Rexroth aluminium M6 bracket \\
T-Bolt & 12 & M6 $\times$ 25 mm \\
Flange nut & 12 & M6 \\
Rubber Spacer/Damper & 4 & M6 \\
\bottomrule
\end{tabular}
\caption{\DIFaddFL{Mounting hardware components}}
\label{tab:mounting_hardware}
\DIFaddendFL \end{table}

\DIFdelbegin \DIFdel{No special tools or non-standard interfaces are required. All mounting interfaces comply with the standard ESCARGO mechanical interface.
}\DIFdelend \DIFaddbegin \begin{figure}[H]
    \centering
    \includegraphics[width=0.45\textwidth]{0_images/figures/mouting.png}
    \includegraphics[width=0.45\textwidth]{0_images/figures/render3.png}
    \caption{\DIFaddFL{CAD images of the mechanical interface.}}
    \label{fig:example_image}
    \DIFaddendFL 

\DIFdelbeginFL \subsection{\DIFdelFL{Electrical}}
%DIFAUXCMD
\addtocounter{subsection}{-1}%DIFAUXCMD
\DIFdelendFL \DIFaddbeginFL \end{figure}
%DIF > %%%%%%%%%%%%%%%%%%%%%%%%%%%%%%%%%%%%%%%%%%%%%%%%%%%%%%%%%%%%%%%%%%%%%%%%%%%%%%
\subsection{\DIFadd{Electrical Interface}}\label{electrical_interface}
\DIFaddend 

\DIFaddbegin \DIFadd{The electrical interface to the Bexus power unit consists of the MIL-C 26482P series 1 connector, as provided by the user manual. Using pin A and B, the supply voltage from the SAFT MP176065 battery pack will be fed to the power distribution board.
The connector to the E-Link of the gondola will be an Amphenol RJF21B and the RJ45 connector will be employed for communication with a RJ45 plug on the MCU PCB.
}

\DIFaddend Since the Bexus power supply is 28.8V \DIFdelbegin \DIFdel{we will use }\DIFdelend two LM2596HVS step-down voltage regulators \DIFdelbegin \DIFdel{on our }\DIFdelend \DIFaddbegin \DIFadd{will be used on the }\DIFaddend power distribution board. \DIFdelbegin \DIFdel{We will use }\DIFdelend \DIFaddbegin \DIFadd{For the MCU and K96 sensor }\DIFaddend 8V \DIFdelbegin \DIFdel{for our Arduino, pressurization system }\DIFdelend \DIFaddbegin \DIFadd{will be used, the pressurization system operates at 24V, the thermal system at 12V }\DIFaddend and \DIFdelbegin \DIFdel{K96 sensor, and }\DIFdelend 3.3V for the \DIFdelbegin \DIFdel{ambient sensors. Our estimated }\DIFdelend \DIFaddbegin \DIFadd{sensors. The estimated average }\DIFaddend current draw from the \DIFdelbegin \DIFdel{power distribution board without the thermal system is 1.15A , which is well below the max current of 3A for the step-down voltage regulators. Regarding our current draw from the Bexus power unit , we estimate 0.35A. The thermal system is still not worked out and }\DIFdelend \DIFaddbegin \DIFadd{BEXUS battery is 0.8A and the average power consumption is 23.11W. The total power consumption from BEXUS power unit }\DIFaddend is \DIFdelbegin \DIFdel{therefore not in the calculations}\DIFdelend \DIFaddbegin \DIFadd{138.67Wh as listed in section 4.7 and the peak current draw is 2.26A. For protecting the BEXUS batteries and other experiments a pre-charge circuit will be implemented in order to limit the currents of the pressurization system and an LC filter to deal with the EMI and switching frequencies of DC/DC converters}\DIFaddend . Distributed single point ground will be implemented.


\DIFdelbegin %DIFDELCMD < \begin{figure}[H]
%DIFDELCMD <     \centering
%DIFDELCMD <     \includegraphics[width=1.0\textwidth]{Graphics/Figures/Skarmbild.png}
%DIFDELCMD <     %%%
%DIFDELCMD < \caption{%
{%DIFAUXCMD
\DIFdelFL{Electrical Schematic}}
    %DIFAUXCMD
%DIFDELCMD < \label{fig:electrical_schematic}
%DIFDELCMD < \end{figure}
%DIFDELCMD < 

%DIFDELCMD < \begin{table}[H]
%DIFDELCMD < \centering
%DIFDELCMD < %%%
%DIFDELCMD < \caption{%
{%DIFAUXCMD
\DIFdelFL{REXUS connector pin out usage}}
%DIFAUXCMD
%DIFDELCMD < \label{tab:rexus_pinout}
%DIFDELCMD < \begin{tabular}{llp{8cm}}
%DIFDELCMD < \toprule
%DIFDELCMD < \textbf{Pin No.} & \textbf{Name} & \textbf{Use} \\
%DIFDELCMD < \midrule
%DIFDELCMD < 1 & +28 V & Experiment Power \\
%DIFDELCMD < 2 & Charging (28 V/1 A) & Not connected -- not used \\
%DIFDELCMD < 3 & SODS & Not connected -- not used \\
%DIFDELCMD < 4 & SOE & Used \\
%DIFDELCMD < 5 & LO & Used \\
%DIFDELCMD < 6 & EXP out+ & Downlink (non-inverted) \\
%DIFDELCMD < 7 & EXP out- & Downlink (inverted) \\
%DIFDELCMD < 8 & 28 V Ground & Ground \\
%DIFDELCMD < 9 & +28 V & Experiment Power \\
%DIFDELCMD < 10 & n.c. & Not connected \\
%DIFDELCMD < 11 & n.c. & Not connected \\
%DIFDELCMD < 12 & Charging Return & Not connected -- not used \\
%DIFDELCMD < 13 & EXP in + & Uplink (non-inverted) \\
%DIFDELCMD < 14 & EXP in - & Uplink (inverted) \\
%DIFDELCMD < 15 & 28 V Ground & Ground \\
%DIFDELCMD < \bottomrule
%DIFDELCMD < \end{tabular}
%DIFDELCMD < \end{table}
%DIFDELCMD < 

%DIFDELCMD < %%%
\DIFdelend %DIF > %%%%%%%%%%%%%%%%%%%%%%%%%%%%%%%%%%%%%%%%%%%%%%%%%%%%%%%%%%%%%%%%%%%%%%%%%%%%%%
\subsection{Thermal}\DIFdelbegin %DIFDELCMD < 

%DIFDELCMD < %%%
\DIFdel{The thermal properties of our }\DIFdelend \DIFaddbegin \label{thermal}
\DIFadd{As outlined in \sref{sec: thermal design}, the thermal behaviour of the }\DIFaddend experiment will be \DIFdelbegin \DIFdel{partially controlled by heating /cooling elements . The specifics remain to be determined once the thermal behaviour of the electronic components can be modelled or tested. Where necessary, we will insulate the whole experiment to avoid an impact of the extreme temperatures, which can go down to -80}%DIFDELCMD < \degree %%%
\DIFdel{C, on the electronics and pumps. This will be also done by rubber bumpers to avoid cooling from the gondola, therefore also mitigating heat transfer to the gondola .
}%DIFDELCMD < 

%DIFDELCMD < %%%
\DIFdel{The thermal impact on the vehicle should be minimal, but this }\DIFdelend \DIFaddbegin \DIFadd{controlled via heating elements and insulation. To avoid thermal coupling to ESCARGO, rubber bumpers between the experiment and the gondola will be used.
The thermal coupling }\DIFaddend will be further tested when \DIFdelbegin \DIFdel{we have all components}\DIFdelend \DIFaddbegin \DIFadd{the experiment is assembled}\DIFaddend . The air inlet and outlet for the air samples to be measured will be directed away from the vehicle and should not heat or cool the vehicle\DIFdelbegin \DIFdel{either}\DIFdelend .

%DIF > %%%%%%%%%%%%%%%%%%%%%%%%%%%%%%%%%%%%%%%%%%%%%%%%%%%%%%%%%%%%%%%%%%%%%%%%%%%%%%
\section{Experiment Components}

\begin{table}[H]
\centering
\caption{Experiment summary table}
\label{tab:experiment_summary}
\DIFdelbeginFL %DIFDELCMD < \begin{tabular}{L{6cm}L{6cm}}
%DIFDELCMD < \toprule
%DIFDELCMD < Experiment mass (in kg): & $\sim$1.45 \\
%DIFDELCMD < Experiment dimensions (in m): & 0.10x0.10x0.05\\
%DIFDELCMD < Experiment expected \ac{cog} position: & Distance from geometric centre: x=4mm y=23mm z=22mm \\
%DIFDELCMD < \bottomrule
%DIFDELCMD < \end{tabular}
%DIFDELCMD < %%%
\DIFdelendFL \DIFaddbeginFL \begin{tabular}{L{6cm}L{6cm}}
\toprule
Experiment mass (in kg): & $\sim$4.0 \\
Experiment dimensions (in m): & 0.20x0.20x0.15\\
Experiment expected \gls{cog} position: & Distance from geometric centre: x=4mm y=23mm z=22mm \\
\bottomrule
\end{tabular}
\DIFaddendFL \end{table}

\begin{table}[H]
\centering
\caption{Electrical components}
\label{tab:electrical_components}
\DIFdelbeginFL %DIFDELCMD < \begin{tabular}{p{4cm}p{4cm}ll}
%DIFDELCMD < \hline
%DIFDELCMD < \textbf{Component} & \textbf{Dimensions (mm)} & \textbf{Mass (g)} & \textbf{Status} \\
%DIFDELCMD < \hline
%DIFDELCMD < Arduino Mega & 54x102x15 & 37 & Preliminary\\
%DIFDELCMD < MS5607 pressure sensor & 3.0x5.0x0.96 & $<5$ & Preliminary\\
%DIFDELCMD < K96 & 42.26x33.9x26.04 & 300 & Delivered \\
%DIFDELCMD < DS18B20 temperature sensor & 4.0x3.3x4.7 & 0.2 & Preliminary \\
%DIFDELCMD < SHT3X-DIS & 2.5x2.5x0.9 & 0.715 & Preliminary \\
%DIFDELCMD < LM2596HVS-ADJ Stepdown DC/DC regulator & 10.1x9.0x4.5 & 1.59 & Premliminary \\
%DIFDELCMD < \ac{pcb} & 40x30x1.6 & 3.6 & Developing phase \\
%DIFDELCMD < Pump & 43x21x14 & 300 & Preliminary \\
%DIFDELCMD < Thermal system & TBD & TBD & Preliminary \\
%DIFDELCMD < Electric valve & TBD & TBD & Preliminary \\
%DIFDELCMD < \hline
%DIFDELCMD < Total & & 648.105 & \\
%DIFDELCMD < \hline
%DIFDELCMD < \end{tabular}
%DIFDELCMD < %%%
\DIFdelendFL \DIFaddbeginFL \begin{tabular}{p{4cm}p{4cm}ll}
\hline
\textbf{Component} & \textbf{Dimensions (mm)} & \textbf{Mass (g)} & \textbf{Quantity}\\
\hline
Vacuum pump & 119$\times$44$\times$73 & 294 & 2 \\
Compressor & 130$\times$60$\times$85 & 445 & 1 \\
Arduino Mega & 54$\times$102$\times$15 & 37 & 1\\
K96 & 42.26$\times$33.9$\times$26.04 & 300 & 1 \\
MS5607 & 3.0$\times$5.0$\times$0.95 & <1 & 1\\
SHT45 & 1.5$\times$1.5$\times$0.5 & <1 & 4 \\
LM2596HVS-ADJ & 10.1$\times$9.0$\times$4.5 & 1.59 & 3 \\
\gls{pcb} & TBD & TBD & 1 \\
Thermal system & TBD & TBD &  \\
\hline
Total & & 1079.59 & \\
\hline
\end{tabular}
\DIFaddendFL \end{table}

%DIF > %%%%%%%%%%%%%%%%%%%%%%%%%%%%%%%%%%%%%%%%%%%%%%%%%%%%%%%%%%%%%%%%%%%%%%%%%%%%%%
\section{Mechanical Design}\DIFaddbegin \label{mechanical_design}
\DIFaddend 

The mechanical design of the \DIFdelbegin \DIFdel{experiment consists of an external protective aluminium housing with dimensions of 100x100x50 }\DIFdelend \DIFaddbegin \DIFadd{\mbox{%DIFAUXCMD
\gls{mirage} }\hskip0pt%DIFAUXCMD
experiment focuses on a compact enclosure housing a pressurized measurement chamber, a multi-stage pump system, supporting electronics, and a heating system. \mbox{%DIFAUXCMD
\gls{mirage} }\hskip0pt%DIFAUXCMD
is designed to fit within the ESCARGO gondola structure while maintaining structural integrity and ensuring safe operation under the environmental conditions expected during the mission.
}

\begin{figure}[htbp]
    \centering
    \includegraphics[width=0.8\textwidth]{0_images/figures/tech_drawing.png}
    \caption{\DIFaddFL{Technical drawing of the system as a whole.}}
    \label{fig:tech_drawing}

\end{figure}

\subsection{\DIFadd{Experiment Enclosure}}

\DIFadd{The experiment is housed in a \mbox{%DIFAUXCMD
\gls{cots} }\hskip0pt%DIFAUXCMD
aluminium enclosure with external dimensions of $200\times 200\times 150$ }\DIFaddend mm and a wall thickness of \DIFaddbegin \DIFadd{expected }\DIFaddend 2 mm\DIFdelbegin \DIFdel{, as well as an internal pressurised chamber with dimensions of 38x33x45 }\DIFdelend \DIFaddbegin \DIFadd{. The enclosure provides mechanical protection for all internal components and acts as the primary structural body of the experiment. A \mbox{%DIFAUXCMD
\gls{cots} }\hskip0pt%DIFAUXCMD
enclosure was selected to reduce manufacturing complexity, minimize cost, and provide a robust structural housing with well-characterized mechanical properties.
The enclosure is directly mounted to the ESCARGO frame as described in \sref{mechanical_interface}. Internally, the enclosure accommodates the pressure chamber, two diaphragm vacuum pumps, one piston compressor, a passive outflow hole, shutter valve, tubing, and electronic control hardware. The enclosure also includes two dedicated holes for air inlet and one for outlet, as well as electrical feedthroughs required for connecting to the BEXUS power and E-link system.
}

\begin{figure}[htbp]
    \centering
    \includegraphics[width=0.45\textwidth]{0_images/figures/tech_drawing_3.png}
    \includegraphics[width=0.45\textwidth]{0_images/figures/tech_drawing_5.png}
    \caption{\DIFaddFL{Technical drawings of experiment enclosure. Drawings are representative and the thickness of both “Main Body” and “Main Lid” will be reduced to 2mm by buying COTS parts.}}
    \label{fig:enclosure}
\end{figure}

\subsection{\DIFadd{Internal Layout}}

\DIFadd{The internal layout of the \mbox{%DIFAUXCMD
\gls{mirage} }\hskip0pt%DIFAUXCMD
experiment is arranged to maintain a compact configuration and minimize the length of the airflow path between the pressurized chamber and the compressor.
Electronic control hardware is mounted along the inner bottom of the enclosure, providing physical isolation between electrical components and the pressurized air system.
This configuration ensures efficient airflow through the chamber while maintaining accessibility for integration and servicing.
}

\begin{figure}[htbp]
    \centering
    \includegraphics[width=0.8\textwidth]{0_images/figures/tech_drawing_6.png}
    \caption{\DIFaddFL{Drawing of the internal layout.}}
    \label{fig:int_layout}

\end{figure}

\subsection{\DIFadd{Pressurized Measurement Chamber}}

\DIFadd{The core element of the experiment is a small pressurized chamber that houses the K96 sensor used for the measurements. The chamber has internal dimensions of approximately $38 \times 33 \times 45$ }\DIFaddend mm and a wall thickness of \DIFdelbegin \DIFdel{15 mmcontaining the K96 sensor, as shown in \fref{fig:CAD}.
The total mass of the mechanical structure, excluding internal components, is approximately 800 g. }\DIFdelend \DIFaddbegin \DIFadd{7 mm. It will be machined from aluminium by the \mbox{%DIFAUXCMD
\gls{mirage} }\hskip0pt%DIFAUXCMD
team at Uppsala University.
}

\DIFadd{The chamber is sealed using a copper gasket to ensure reliable sealing under elevated pressure conditions and temperature changes. Copper gaskets were selected due to their ability to provide a robust metal-to-metal seal capable of maintaining integrity at high pressure differentials.
}\DIFaddend 

\DIFdelbegin \DIFdel{The internal pressurised chamber operates at pressures of up to 5 bar , while }\DIFdelend \DIFaddbegin \DIFadd{The chamber is designed to minimize internal volume while providing sufficient space for the sensor and the necessary airflow path. Reducing the internal volume proportionally lowers the airflow rate required to exchange the chamber atmosphere, which improves the overall performance of the experiment.
}

\begin{figure}[htbp]
    \centering
    \includegraphics[width=0.45\textwidth]{0_images/figures/render2.png}
    \includegraphics[width=0.45\textwidth]{0_images/figures/render1.png}
    \caption{\DIFaddFL{Renders of the pressurized chamber.}}
    \label{fig:pressure_render}
\end{figure}

\begin{figure}[htbp]
    \centering
    \includegraphics[width=0.45\textwidth]{0_images/figures/tech_drawing_2.png}
    \includegraphics[width=0.45\textwidth]{0_images/figures/tech_drawing_4.png}
    \caption{\DIFaddFL{Technical drawings of the pressurized chamber.}}
    \label{fig:pressure_drawing}
\end{figure}

\subsection{\DIFadd{Pump System Architecture}}\label{sec: pump system arch}

\DIFadd{The \mbox{%DIFAUXCMD
\gls{mirage} }\hskip0pt%DIFAUXCMD
experiment utilizes a multi-stage air circulation system to maintain a controlled internal pressure while ensuring a continuous exchange of air within the measurement chamber.
}

\textbf{\DIFadd{System Operation:}}

\begin{enumerate}
    \item \DIFadd{Pre-heating: Ambient intake air is pre-heated before entering the vacuum stage to prevent condensation or freezing within the system components.
    }\item \DIFadd{Vacuum Stage: Two diaphragm vacuum pumps operate continuously from launch up to an altitude of 22 km. These pumps supply the secondary stage with air at approximately atmospheric pressure ($\approx$ 1 bar). Since the volumetric flow rate decreases with decreasing intake pressure, the pumps will operate at higher settings at higher altitudes.
}

    \item \DIFadd{Intercooling: To manage thermal loads generated by compression, the tubing connecting the vacuum pumps to the piston compressor serves as a passive intercooler.
    }\item \DIFadd{Compression Stage: A piston compressor increases the airflow pressure to a maximum of 3 bar before it enters the measurement chamber.
    }\item \DIFadd{Pressure Regulation: The internal chamber pressure is maintained by a passive outflow valve equipped with a shutter mechanism. This valve regulates the exhaust flow rate to the external environment.
    }\item \DIFadd{Monitoring: Integrated pressure and temperature sensors are situated within the tubing between compression stages to provide real-time performance data and ensure stable operation.
}\end{enumerate}

%DIF > % ADD FIGURE %%

\textbf{\DIFadd{Compressor Design Considerations:}} \DIFadd{Two primary design methodologies are being evaluated for the compression stage:
}

\begin{enumerate}
    \item \DIFadd{Continuous Low-Flow Operation: Utilizing a compressor capable of maintaining a continuous, very low flow rate to match the chamber's exchange requirements.
    }\item \DIFadd{Intermittent Buffered Operation: Utilizing a higher flow rate compressor operated in duty cycles. This configuration requires a 1 bar buffer tank between pressurisation stages to accumulate air from the vacuum stage. The compressor activates to deplete the tank and deactivates once the tank pressure reaches a lower threshold.
}\end{enumerate}

\DIFadd{The nominal operating pressure inside the chamber is expected to reach approximately 3 bar above perfect vacuum. At float altitude, }\DIFaddend the external ambient pressure \DIFdelbegin \DIFdel{at float altitude can be as low as }\DIFdelend \DIFaddbegin \DIFadd{may decrease to }\DIFaddend approximately 10 mbar, resulting in a \DIFdelbegin \DIFdel{maximum }\DIFdelend pressure differential of approximately \DIFdelbegin \DIFdel{5 bar which }\DIFdelend \DIFaddbegin \DIFadd{3 bar across the chamber walls. This pressure differential }\DIFaddend is accounted for in the structural design \DIFaddbegin \DIFadd{of the chamber}\DIFaddend .

The \DIFdelbegin \DIFdel{outer housing has one detachable wall fastened with four M2x10 screws, allowing access during integration and testing. The housing includes openings for air inlet and outlet for the pressurisation system, as well as internal mounting provisions for electrical components}\DIFdelend \DIFaddbegin \DIFadd{pressure between the vacuum pumps and the compressor is expected to remain close to atmospheric pressure ($\approx$ 1 bar) during normal operation in order to match expected compressor efficiency.
}

\DIFadd{Temperature conditions at float altitude will require the implementation of pre-heating of the intake air before the vacuum stage to prevent condensation or freezing within the system. Additionally, pressure sensors will be placed between compression stages to monitor system performance and ensure stable pressure control during operation.
}

%DIF >  \subsection{In-line Sensing and Environmental Monitoring}

%DIF >  To ensure the operational integrity of the gas handling system and to facilitate real-time performance analysis, the \gls{mirage} experiment incorporates an integrated sensing suite within the tube circuit. This monitoring architecture is essential for validating the state of the air as it transitions through the various compression and thermal management stages.

%DIF >  \begin{enumerate}
%DIF >      \item Pressure Sensing: High-precision pressure sensors are situated within the tubing between the diaphragm vacuum pumps and the piston compressor. These sensors monitor the inter-stage pressure to ensure that the vacuum pumps, operating continuously up to 22 km, maintain the required 1 bar supply pressure for the compressor.
%DIF >      \item Temperature Sensing: Thermal sensors are embedded directly into the flow path to monitor the effectiveness of the pre-heating and intercooling mechanisms. Monitoring the air temperature post-heating is critical to verify that the intake air has reached a safe threshold before entering the pump heads and K96 sensor chamber, thereby preventing damage from cold or heat
%DIF >  \end{enumerate}

\subsection{\DIFadd{Pump and Fluid Handling Technologies}}

\DIFadd{Small gas circulation and sampling systems used in atmospheric sensing instruments typically rely on compact electrically driven pumps to move air through measurement chambers and sensors. Several pump technologies are commonly used in portable gas analysis systems, including diaphragm pumps, piston compressors, rotary vane pumps, and scroll pumps. Each technology offers different advantages in terms of pressure capability, size, reliability, and maintenance requirements.
}

\DIFadd{Diaphragm pumps are widely used in environmental monitoring and portable gas sampling instruments. These pumps operate by oscillating a flexible diaphragm using an electric motor, creating alternating suction and compression cycles. Because the gas is hermetically separated the mechanical drive components by the diaphragm, diaphragm pumps operate oil-free and are resistant to contamination. This makes them particularly suitable for experiments involving sensitive sensors or atmospheric gas sampling. Furthermore, diaphragm pumps are compact, lightweight, and tolerant of condensation in the airflow, which is expected in the \mbox{%DIFAUXCMD
\gls{mirage} }\hskip0pt%DIFAUXCMD
experiment due to low external temperatures at float altitude.
}

\DIFadd{Piston compressors are commonly used when higher pressures are required. These pumps use a reciprocating piston to compress air within a cylinder, allowing significantly higher pressure ratios than typical diaphragm pumps. However, piston compressors generally contain more moving mechanical parts and may be more sensitive to contamination or thermal effects. For this reason, piston compressors are often used in series with diaphragm pumps in small air circulation systems, where the diaphragm pumps provide initial airflow and the compressor increases the pressure to the required level.
}

\DIFadd{Other pump technologies, such as rotary vane pumps and scroll pumps, are frequently used in laboratory vacuum systems. While these technologies can achieve significantly lower pressures and higher pumping speeds, they typically require lubrication, larger mechanical assemblies, and higher power consumption. Due to these characteristics, they are generally unsuitable for compact balloon payload experiments such as \mbox{%DIFAUXCMD
\gls{mirage}}\hskip0pt%DIFAUXCMD
, where mass, power consumption, and system simplicity are critical design constraints}\DIFaddend .

\DIFdelbegin \DIFdel{The internal pressurised chambercontaining the K96 sensor includes dedicated interfaces for air inlet, air outlet aligned with the openings in the external housing, and electrical feedthroughs.
One wall of the chamber is detachable and fastened with four M2x10 screws to allow sensor installation and servicing}\DIFdelend \DIFaddbegin \DIFadd{Based on this comparison, diaphragm pumps were selected as the primary airflow drivers in the \mbox{%DIFAUXCMD
\gls{mirage} }\hskip0pt%DIFAUXCMD
experiment due to their oil-free operation, compact size, reliability, and resistance to condensation. A piston compressor is used downstream of the diaphragm pumps to achieve the required pressurization of the measurement chamber}\DIFaddend .

\DIFdelbegin \DIFdel{The pressurisation system is based on a }\DIFdelend \DIFaddbegin \subsection{\DIFadd{Tubing and Flow Control Components}}

\DIFadd{Gas flow regulation in compact atmospheric sensing instruments is traditionally achieved through the use of miniature solenoid or needle valves. Solenoid valves are often preferred for their rapid response times and ability to provide high-precision, electronically controlled flow regulation within automated systems.
However, for the \mbox{%DIFAUXCMD
\gls{mirage} }\hskip0pt%DIFAUXCMD
experiment, the design has been optimized to utilize a passive outflow valve equipped with a mechanical shutter. This shift to a passive regulation system is engineered to maintain the required internal chamber pressure of approximately 3 bar by modulating the exhaust flow rate (hole size) based on the pressure differential encountered during the flight profile. In addition the usage of a shutter-based passive mechanism rather than an active solenoid, the system achieves greater mechanical simplicity and safety if the ballon would fall into the water.
The final selection of the hole and the shutter's mechanical properties is based on ensuring a stable airflow path that accommodates the continuous operation of the upstream diaphragm pumps up to the 22 km altitude threshold.
}

\subsection{\DIFadd{Tubing and Fluid Connections}}

\DIFadd{The airflow between the pumps, compressor, and pressurized measurement chamber is transported through flexible tubing. In compact gas circulation systems, tubing connections are typically achieved using barbed fittings, compression fittings, or push-to-connect pneumatic fittings.
}

\DIFadd{Barbed fittings are commonly used with flexible tubing materials such as silicone or polyurethane, where the tubing is pressed onto the fitting and secured through friction or hose clamps. Compression fittings provide a more robust and leak-tight connection and are frequently used in higher-pressure gas systems. Push-to-connect fittings are widely used in pneumatic systems due to their ease of installation and reliable sealing.
}

\DIFadd{For the \mbox{%DIFAUXCMD
\gls{mirage} }\hskip0pt%DIFAUXCMD
experiment, the tubing connections will be selected to ensure leak-tight operation at pressures up to approximately 3 bar while maintaining compatibility with the chosen tubing material and the compact geometry of the experimental setup. Particular attention will be given to minimizing leakage and dead volumes in the system to maintain stable pressure control and reliable airflow through the measurement chamber.
}

\DIFadd{In alignment with the experiment's thermal management requirements, the intake tube lines will be integrated with high-conductivity copper windings to facilitate efficient heat transfer between the enclosure and the vacuum pump interface. This thermal exchange mechanism is designed to rapidly elevate the temperature of the intake air, ensuring the fluid stream remains within the pump's safe operational range. This will also be the case between compressor and the chamber as the heat created by compression could affect the quality of the study. This proactively mitigates the risk of mechanical failure or structural damage to the diaphragm components caused by the extreme low-temperature environment encountered during the ascent to 22 km and high temperatures after compression.
}

\subsection{\DIFadd{Pumps}}

\DIFadd{The air circulation system utilizes commercially available miniature medical pumps. The vacuum pumps are compact DC diaphragm pumps capable of generating the suction required to draw air into the system at 10 mbar. Two pumps are used to provide sufficient airflow at 1 bar for the compressor and to enhance system robustness.
}

\DIFadd{The compressor is a small DC piston air }\DIFaddend pump used to increase the \DIFdelbegin \DIFdel{internal chamber pressure and an electronically controlled valve used to enable controlled flushing of the chamber.
}\DIFdelend \DIFaddbegin \DIFadd{pressure of the air entering the pressurized chamber. Both components were selected based on their cost, compact size, low mass, and compatibility with the power constraints of the BEXUS payload system.
}\DIFaddend 

\DIFaddbegin \subsection{\DIFadd{Structural Considerations}}

\DIFadd{The mechanical structure is designed with sufficient strength and safety margin to withstand internal pressure differentials exceeding the maximum expected value of approximately 3 bar between the pressurized chamber and the external environment.
The pressurized chamber is manufactured from solid machined aluminium with a wall thickness of 7 mm. This thickness provides a structural safety factor far greater than 2 relative to the maximum expected pressure differential, accounting for uncertainties in material properties, manufacturing tolerances, and thermal stresses.
The outer aluminium enclosure provides mechanical protection and structural support for the internal components. A copper gasket is used to ensure reliable sealing and minimize the risk of leakage during operation.
The internal layout maintains a compact configuration while ensuring that mechanical loads generated by the pumps, tubing, and mounting hardware are safely transferred to the enclosure structure.
}

\subsection{\DIFadd{Mass Budget}}

\DIFadd{A preliminary mass budget of the \mbox{%DIFAUXCMD
\gls{mirage} }\hskip0pt%DIFAUXCMD
experiment was derived from the SolidWorks CAD assembly model. The estimated mass distribution of the main structural, fluid handling, electronic, and fastening components is summarized in \tref{tab:mass_budget}. Smaller components such as fasteners and connectors are included to provide a realistic estimate of the total system mass.
}

\begin{table}[htbp]
\centering
\begin{tabular}{lccc}
\toprule
\textbf{Component} & \textbf{Qty} & \textbf{Mass (g)} & \textbf{Total (g)} \\
\midrule

\multicolumn{4}{l}{\textbf{Structural components}} \\

COTS aluminium enclosure & 1 & 1200 (TBD) & 1200 \\
Pressurized chamber body & 1 & 159 & 159 \\
Pressurized chamber lid & 1 & 70 & 70 \\
Vacuum pump arch & 1 & $\sim$30 & 30 \\
Valve bracket & 1 & $\sim$20 & 20 \\

\addlinespace
\multicolumn{4}{l}{\textbf{Fluid system}} \\

Diaphragm vacuum pump & 2 & $\sim$294 & 588 \\
Piston compressor & 1 & $\sim$445 & 445 \\
Shutter valve & 1 & 250 (TBD) & 250 \\
Air tubing assembly & 1 & $\sim$120 & 120 \\

\addlinespace
\multicolumn{4}{l}{\textbf{Electronics}} \\

Arduino Mega 2560 & 1 & 37 & 37 \\
RJF21B connector & 1 & $\sim$15 & 15 \\

\addlinespace
\multicolumn{4}{l}{\textbf{Fasteners}} \\

M3 screws & 13 & $\sim$1 & 13 \\
M3.5 screws & 4 & $\sim$2 & 8 \\
M3 long screws & 4 & $\sim$2 & 8 \\
M6 mounting bolts & 4 & $\sim$10 & 40 \\

\midrule
\textbf{Estimated total} & & & \textbf{$\sim$3003} \\
\bottomrule
\end{tabular}
\caption{\DIFaddFL{Mass budget of system components.}}
\label{tab:mass_budget}
\end{table}

\subsection{\DIFadd{Component Procurement Summary}}

\DIFadd{The experiment design utilizes a combination of \mbox{%DIFAUXCMD
\gls{cots} }\hskip0pt%DIFAUXCMD
components and custom-manufactured parts. \mbox{%DIFAUXCMD
\gls{cots} }\hskip0pt%DIFAUXCMD
components are standard products purchased directly from manufacturers without modification, while custom components are specifically designed and machined for the \mbox{%DIFAUXCMD
\gls{mirage} }\hskip0pt%DIFAUXCMD
experiment.
}

\DIFadd{The \mbox{%DIFAUXCMD
\gls{cots} }\hskip0pt%DIFAUXCMD
components include the aluminium enclosure, diaphragm vacuum pumps, piston compressor, shutter valve, tubing, and standard fasteners. These components were selected in order to reduce manufacturing complexity, cost, and lead time.
}

\DIFadd{Custom-manufactured components include the pressurized measurement chamber (chamber body and lid) and several structural mounting elements such as pump holders, valve holders, and supporting brackets. These parts are machined from aluminium at Uppsala University to meet the structural and pressure requirements of the experiment.
}

\DIFadd{The classification of the main components as \mbox{%DIFAUXCMD
\gls{cots} }\hskip0pt%DIFAUXCMD
or custom-manufactured is summarized in \tref{tab:component_type_summary}.
}

\begin{table}[htbp]
\centering
\begin{tabular}{ll}
\toprule
\textbf{Component} & \textbf{Type} \\
\midrule

Main enclosure & COTS \\
Diaphragm vacuum pumps & COTS \\
Piston compressor & COTS \\
Shutter valve & COTS \\
Arduino Mega 2560 & COTS \\
RJF21B connector & COTS \\

Pressurized chamber body & Custom machined \\
Pressurized chamber lid & Custom machined \\
Valve holder & Custom machined \\
Pump mounting brackets & Custom machined \\
Structural holders and supports & Custom machined \\

Tubing and pneumatic fittings & COTS \\
Fasteners (M3/M6 bolts, screws) & COTS \\

\bottomrule
\end{tabular}
\caption{\DIFaddFL{Component Type Summary}}
\label{tab:component_type_summary}
\end{table}

%DIF > %%%%%%%%%%%%%%%%%%%%%%%%%%%%%%%%%%%%%%%%%%%%%%%%%%%%%%%%%%%%%%%%%%%%%%%%%%%%%%
\DIFaddend \section{Electronics Design}

\DIFdelbegin \subsection{\DIFdel{Processing core unit}}
%DIFAUXCMD
\addtocounter{subsection}{-1}%DIFAUXCMD
\DIFdelend \DIFaddbegin \DIFadd{The overall electronics architecture of the \mbox{%DIFAUXCMD
\gls{mirage} }\hskip0pt%DIFAUXCMD
experiment is shown in \fref{fig:electrical_schematic}. The system consists of a Power Distribution Board, which interfaces with the BEXUS gondola power supply and provides regulated voltages to the experiment subsystems. The processing core is based on an Arduino Mega microcontroller, which manages sensor data acquisition, system control, and communication with onboard storage.
}\DIFaddend 

\DIFaddbegin \DIFadd{Additional subsystems include the pressurization system, the thermal control system, the K96 NDIR sensor core, and external environmental sensors measuring ambient temperature, pressure, and humidity. The ambient sensors communicate with the processing unit via an I²C interface, while the K96 NDIR sensor communicates using a UART interface. The interactions between these subsystems and the processing unit are illustrated in \fref{fig:electrical_schematic}. Each subsystem is described in more detail in the following sections.
}

\DIFaddend \begin{figure}[H]
    \centering
        \DIFdelbeginFL %DIFDELCMD < \includegraphics[width=0.5\textwidth]{Graphics/Figures/Arduino MEGA Rev3.png}
%DIFDELCMD <     %%%
\DIFdelendFL \DIFaddbeginFL \includegraphics[width=1.0\textwidth]{0_images/figures/Mirage_Electrical_Schematic.png}
            \DIFaddendFL \caption{\DIFdelbeginFL \DIFdelFL{Arduino Mega Rev3 WiFi}\DIFdelendFL \DIFaddbeginFL \DIFaddFL{Electrical Schematic}\DIFaddendFL }
                \DIFdelbeginFL %DIFDELCMD < \label{fig:component_arduino}
%DIFDELCMD < %%%
\DIFdelendFL \DIFaddbeginFL \label{fig:electrical_schematic}
                \DIFaddendFL \end{figure}

\DIFdelbegin \DIFdel{An Arduino mega }\DIFdelend \DIFaddbegin \subsection{\DIFadd{Processing core unit}}

%DIF > \begin{figure}[H]
%DIF >     \centering
%DIF >     \includegraphics[width=0.5\textwidth]{0_images/figures/Arduino MEGA Rev3.png}
%DIF >     \caption{Arduino Mega Rev3 WiFi}
%DIF >     \label{fig:component_arduino}
%DIF > \end{figure}


\DIFadd{An Arduino Mega }\DIFaddend will be used for preliminary testing and \DIFdelbegin \DIFdel{is to }\DIFdelend \DIFaddbegin \DIFadd{as a backup solution, and will }\DIFaddend be replaced by a \DIFdelbegin \DIFdel{processing \mbox{%DIFAUXCMD
\ac{pcb} }\hskip0pt%DIFAUXCMD
once assembly starts}\DIFdelend \DIFaddbegin \DIFadd{dedicated processing \mbox{%DIFAUXCMD
\gls{pcb} }\hskip0pt%DIFAUXCMD
once assembly begins}\DIFaddend . The Arduino \DIFdelbegin \DIFdel{mega }\DIFdelend \DIFaddbegin \DIFadd{Mega }\DIFaddend is equipped with an \DIFdelbegin \DIFdel{ethernet }\DIFdelend \DIFaddbegin \DIFadd{Ethernet }\DIFaddend shield, enabling it to communicate with EGSE and save measurement data to an \DIFdelbegin \DIFdel{SD-card}\DIFdelend \DIFaddbegin \DIFadd{SD card}\DIFaddend . The Arduino needs a power input between 7-12V\DIFaddbegin \DIFadd{, }\DIFaddend which will be converted from the 28.8V \DIFdelbegin \DIFdel{from the SAFT batteries }\DIFdelend provided by Bexus\DIFaddbegin \DIFadd{' SAFT batteries }\DIFaddend through a DC/DC stepdown regulator.

\subsection{\DIFdelbegin \DIFdel{Power Distribution Unit}\DIFdelend \DIFaddbegin \DIFadd{MicroSD Card}\DIFaddend }

\DIFdelbegin %DIFDELCMD < \begin{figure}[H]
%DIFDELCMD <     \centering
%DIFDELCMD <     \includegraphics[width=0.3\textwidth]{Graphics/Figures/Stepdown DCDC.png}
%DIFDELCMD <     %%%
%DIFDELCMD < \caption{%
{%DIFAUXCMD
\DIFdelFL{LM2596HVS-ADJ step-down DC/DC regulator}}
    %DIFAUXCMD
%DIFDELCMD < \label{fig:component_LM2596HVS-ADJ_DC/DC_regulator}
%DIFDELCMD < \end{figure}
%DIFDELCMD < %%%
\DIFdelend %DIF > \begin{figure}[H]
%DIF >     \centering
%DIF >     \includegraphics[width=0.3\textwidth]{0_images/figures/MicroSD card.png}
%DIF >     \caption{SDHC Memory Card}
%DIF >     \label{fig:component_SDHC}
%DIF > \end{figure}

\DIFdelbegin \DIFdel{Power distribution }\DIFdelend \DIFaddbegin \DIFadd{Measurement data will be stored on a microSDHC memory card (SP016GBSTHBU1V10SP) with a capacity of 16~GB. The card has maximum read and write speeds of up to 85~MB/s and 15~MB/s, respectively. The device operates within a supply voltage range of 2.7--3.6~V and an operating temperature range of 0}\degree \DIFadd{C to 70}\degree \DIFadd{C.
}

\DIFadd{The card interfaces with the processing core unit via the SPI (Serial Peripheral Interface), using the MOSI, MISO, SCK, and CS signals for data transfer and control.
}

\subsection{\DIFadd{Power Distribution Board}}

%DIF > \begin{figure}[H]
%DIF >     \centering
%DIF >     \includegraphics[width=0.3\textwidth]{0_images/figures/Stepdown DCDC.png}
%DIF >     \caption{LM2596HVS-ADJ step-down DC/DC regulator}
%DIF >     \label{fig:component_LM2596HVS-ADJ_DC/DC_regulator}
%DIF > \end{figure}
\DIFadd{Power distribution within the experiment }\DIFaddend is handled by \DIFdelbegin \DIFdel{multiple }\DIFdelend \DIFaddbegin \DIFadd{a dedicated Power Distribution Board (PDB). The board receives a nominal input voltage of 28.8 V from the BEXUS power system and distributes regulated voltages to the experiment subsystems.
}

\DIFadd{Voltage regulation is implemented using three }\DIFaddend LM2596HVS-ADJ step-down DC/DC \DIFdelbegin \DIFdel{regulators}\DIFdelend \DIFaddbegin \DIFadd{converters}\DIFaddend . These regulators \DIFaddbegin \DIFadd{operate at a switching frequency of 150 kHz, }\DIFaddend support input voltages up to 57 V\DIFaddbegin \DIFadd{, }\DIFaddend and provide adjustable output voltages in the range of 1.2\DIFdelbegin \DIFdel{--37 }\DIFdelend \DIFaddbegin \DIFadd{–37 }\DIFaddend V. The \DIFdelbegin \DIFdel{SAFT battery supplies a nominal voltage of 28.8 V , which is regulated down to the required systemvoltages of 3.3 V and }\DIFdelend \DIFaddbegin \DIFadd{converters generate three regulated voltage rails required by the experiment: 24 V for the thermal and pressurization systems (compressor and vacuum pumps), }\DIFaddend 8 \DIFdelbegin \DIFdel{V}\DIFdelend \DIFaddbegin \DIFadd{V for the K96 NDIR sensor and the MCU, and 3.3 V for the environmental sensors.
}

\DIFadd{To mitigate switching noise introduced by the DC/DC converters on the shared power lines, an LC filter with a cut-off frequency of approximately 4 kHz is implemented on the power distribution board.
}

\DIFadd{Due to the high startup currents expected from the pressurization system, a pre-charge circuit is implemented on the power distribution board to limit inrush currents during system startup}\DIFaddend .



\subsection{K96 Sensor}

\begin{figure}[H]
    \centering
    \DIFdelbeginFL %DIFDELCMD < \includegraphics[width=0.5\textwidth]{Graphics/Figures/K96 Sensor.png}
%DIFDELCMD <     %%%
\DIFdelendFL \DIFaddbeginFL \includegraphics[width=0.5\textwidth]{0_images/figures/K96 Sensor.png}
    \DIFaddendFL \caption{K96 Sensor. Taken from \cite{wastine2022}}
    \label{fig:component_K96_sensor}
\end{figure}

The K96 sensor is a \DIFdelbegin \DIFdel{\mbox{%DIFAUXCMD
\ac{ndir} }\hskip0pt%DIFAUXCMD
}\DIFdelend \DIFaddbegin \DIFadd{\mbox{%DIFAUXCMD
\gls{ndir} }\hskip0pt%DIFAUXCMD
}\DIFaddend gas sensor capable of measuring methane, carbon dioxide, and water vapor. It operates from a 6.0--8.5 V VBB supply and a 2.4--3.6 V IOVCC/RxD logic supply over a temperature range of $-40\degree$C to +85\degree C, with a maximum peak current of 300 mA. The sensor also provides local temperature, pressure, and humidity measurements via an integrated BME280 sensor.

\subsection{\DIFdelbegin \DIFdel{Pump}%DIFDELCMD < \MBLOCKRIGHTBRACE
%DIFDELCMD < %%%
\DIFdel{The pressurization system will include an electrically driven pump used to increase the pressure inside the measurement chamber. The pump will be powered from the experiment power distribution system and controlled }\DIFdelend \DIFaddbegin \DIFadd{Temperature and Humidity Sensor}}

%DIF > \begin{figure}[H]
%DIF >     \centering
%DIF >     \includegraphics[width=0.3\textwidth]{0_images/figures/Humidity Sensor.png}
%DIF >     \caption{SHT3X-DIS Humidity Sensor}
%DIF >     \label{fig:component_SHT3X-DIS_humidity}
%DIF > \end{figure}

\DIFadd{The ambient temperature and humidity sensor is a SHT45, with an operating range of -40 to 125$^\circ$C and a humidity measurement range of 0--100\% RH. The sensor provides a humidity accuracy of $\pm 1.0$ \% RH and a temperature accuracy of $\pm 0.1\,^\circ$C. Communication with the sensor is performed via an I$^2$C digital interface, allowing temperature and humidity data to be read directly }\DIFaddend by the processing unit. \DIFaddbegin \DIFadd{The operating voltage range is 1.08 V to 3.6 V, with a typical average current consumption of approximately 0.4 $\mu$A.
}\DIFaddend 

\DIFdelbegin \DIFdel{At the current design stage, the specific pump model has not yet been selected. Candidate solutions include compact diaphragm- or piston-based pumps capable of operating at low inlet pressures and generating the required pressure differential for chamber pressurization. Final selection will be based on electrical power consumption, achievable pressure ratio, mass, and compatibility with low-temperature operation.
}%DIFDELCMD < 

%DIFDELCMD < %%%
\DIFdel{The pump will be operated in combination with an electronically controlled outlet valve to enable controlled air exchange of the measurement chamber.
Detailed electrical characteristics and control parameters will be defined during the detailed design phase following component selection and testing.
}%DIFDELCMD < 

%DIFDELCMD < %%%
\DIFdelend \subsection{\DIFdelbegin \DIFdel{Ambient }\DIFdelend Pressure Sensor}
\DIFdelbegin %DIFDELCMD < 

%DIFDELCMD < \begin{figure}[H]
%DIFDELCMD <     \centering
%DIFDELCMD <     \includegraphics[width=0.5\textwidth]{Graphics/Figures/MS5607 Pressure Sensor.jpg}
%DIFDELCMD <     %%%
%DIFDELCMD < \caption{%
{%DIFAUXCMD
\DIFdelFL{MS5607 Pressure Sensor}}
    %DIFAUXCMD
%DIFDELCMD < \label{fig:component_MS5607_pressure}
%DIFDELCMD < \end{figure}
%DIFDELCMD < 

%DIFDELCMD < %%%
\DIFdelend The ambient pressure sensor MS5607 is a barometric sensor that provides a highly linear
and precise 24-bit pressure and temperature value. It has \DIFdelbegin \DIFdel{I2C }\DIFdelend \DIFaddbegin \DIFadd{I$^2$C }\DIFaddend and SPI interface up to
20MHz and has can measure between 10-1200mbar with an accuracy of \DIFdelbegin \DIFdel{$\pm$}\DIFdelend \DIFaddbegin \DIFadd{±}\DIFaddend 1.5mbar. The
operating temperature range is \DIFdelbegin \DIFdel{$-40$ }\DIFdelend \DIFaddbegin \DIFadd{-40 }\DIFaddend to 85\degree C and the supply voltage is between 1.8 and
3.6V

\subsection{\DIFdelbegin \DIFdel{Ambient Temperature Sensor}%DIFDELCMD < \MBLOCKRIGHTBRACE
%DIFDELCMD < %%%
\DIFdelend \DIFaddbegin \DIFadd{Pressurization system}}
\DIFadd{The pressurization system is powered through the experiment power distribution board and controlled by the processing unit. The system includes two diaphragm vacuum pumps and a piston compressor, which together provide the airflow and pressure increase required for the measurement chamber. The system operates at 24 V and is supplied through the experiment power distribution system.
}\DIFaddend 

\DIFdelbegin %DIFDELCMD < \begin{figure}[H]
%DIFDELCMD <     \centering
%DIFDELCMD <     \includegraphics[width=0.3\textwidth]{Graphics/Figures/Temperature Sensor.png}
%DIFDELCMD <     %%%
%DIFDELCMD < \caption{%
{%DIFAUXCMD
\DIFdelFL{DS18B20 Temperature Sensor}}
    %DIFAUXCMD
%DIFDELCMD < \label{fig:component_DS18B20_temperature}
%DIFDELCMD < \end{figure}
%DIFDELCMD < 

%DIFDELCMD < %%%
\DIFdel{The temperature sensor DS18B20 is a transistor that operates in the temperature range of $-40$ to 85}%DIFDELCMD < \degree %%%
\DIFdel{C with a 0.5}%DIFDELCMD < \degree %%%
\DIFdel{C accuracy through a 9- or 12-bit line.
It has an operating voltage between 3.0 to 5.5V and an active current of 1.5mA}\DIFdelend \DIFaddbegin \DIFadd{Pressure sensors will be installed between the compression stages and inside the measurement chamber to monitor system performance. These sensors provide feedback to the processing unit, enabling monitoring of the internal pressure and verification of correct system operation}\DIFaddend .

\DIFdelbegin \subsection{\DIFdel{Humidity Sensor}}
%DIFAUXCMD
\addtocounter{subsection}{-1}%DIFAUXCMD
\DIFdelend \DIFaddbegin \DIFadd{A passive outlet is used to allow air to exit the measurement chamber. The outlet is implemented as a fixed flow restriction that enables continuous exhaust of air while the pumps and compressor maintain the required chamber pressure. The internal chamber pressure is therefore primarily regulated by the pressurization system, while the passive outlet ensures continuous gas exchange within the chamber.
}\DIFaddend 

\DIFdelbegin %DIFDELCMD < \begin{figure}[H]
%DIFDELCMD <     \centering
%DIFDELCMD <     \includegraphics[width=0.3\textwidth]{Graphics/Figures/Humidity Sensor.png}
%DIFDELCMD <     %%%
%DIFDELCMD < \caption{%
{%DIFAUXCMD
\DIFdelFL{SHT3X-DIS Humidity Sensor}}
    %DIFAUXCMD
%DIFDELCMD < \label{fig:component_SHT3X-DIS_humidity}
%DIFDELCMD < \end{figure}
%DIFDELCMD < 

%DIFDELCMD < %%%
\DIFdel{The ambient humidity sensor is a SHT3X-DIS with an operating range of $-40$ to 125}%DIFDELCMD < \degree %%%
\DIFdel{C. The accuracy of the sensor is 1.5\% RH down to 0}%DIFDELCMD < \degree %%%
\DIFdel{C, which in our case could be a problem when we measure at lower temperatures. This needs to be tested or an alternative must be found. The operating voltage is 2.15V to 5.5V with typical current draw of 1.7\micro A}\DIFdelend \DIFaddbegin \DIFadd{Detailed electrical interfaces, component specifications, and control parameters will be defined during the detailed design phase following component selection and testing}\DIFaddend .

\DIFdelbegin \subsection{\DIFdel{MicroSD Card}}
%DIFAUXCMD
\addtocounter{subsection}{-1}%DIFAUXCMD
%DIFDELCMD < 

%DIFDELCMD < %%%
\DIFdelend %DIF > %%%%%%%%%%%%%%%%%%%%%%%%%%%%%%%%%%%%%%%%%%%%%%%%%%%%%%%%%%%%%%%%%%%%%%%%%%%%%%
\DIFaddbegin \section{\DIFadd{Thermal Design}}\label{sec: thermal design}
\DIFaddend \begin{figure}[H]
    \centering
    \DIFdelbeginFL %DIFDELCMD < \includegraphics[width=0.3\textwidth]{Graphics/Figures/MicroSD card.png}
%DIFDELCMD <     %%%
%DIFDELCMD < \caption{%
{%DIFAUXCMD
\DIFdelFL{SDHC Memory Card}}
    %DIFAUXCMD
%DIFDELCMD < \label{fig:component_SDHC}
%DIFDELCMD < %%%
\DIFdelendFL \DIFaddbeginFL \includegraphics[width=0.8\textwidth]{0_images/figures/pressure_temperature_profile.png}
    \caption{\DIFaddFL{Averaged pressure and temperature profile obtained by the Sodankylä weather station in Finland over all days of October 2025.}}
    \label{fig:temp_press_profile}
\DIFaddendFL \end{figure}

\DIFdelbegin \DIFdel{To store measurement data, a SanDisk SDSDQAB-008G microSD card will be connected to the processing core unit. It operates at a voltage range of 2.7--3.6 V and supports an operating temperature range of $-25\degree$ C to +85}%DIFDELCMD < \degree %%%
\DIFdel{C.
}%DIFDELCMD < 

%DIFDELCMD < %%%
\section{\DIFdel{Thermal Design}}
%DIFAUXCMD
\addtocounter{section}{-1}%DIFAUXCMD
%DIFDELCMD < 

%DIFDELCMD < %%%
\DIFdelend The temperature conditions in \DIFdelbegin \DIFdel{the stratosphere }\DIFdelend \DIFaddbegin \DIFadd{a stratospheric balloon flight }\DIFaddend are quite extreme with temperatures reaching as low as -80\degree C\DIFdelbegin \DIFdel{at the altitudes of about 30 km we could reach}\DIFdelend . At the same time the air pressure can drop as low as 10 mbar, which also implies that convective heat transfer is very weak\DIFdelbegin \DIFdel{and most heat can be dissipated by radiative transfer only. 
This can lead to overheating of the electrical components if the heat from dissipated power accumulates.Therefore, it is essential to know how much power is dissipated and to take additional cooling measures in case that radiative cooling is insufficient. }\DIFdelend \DIFaddbegin \DIFadd{, which might cause a poorly designed system to overheat. An averaged ambient pressure and temperature profile is shown in \fref{fig:temp_press_profile}. The low temperatures might cause substantial convective cooling at lower altitudes of around 8 km, while the low pressure in combination with solar radiation might cause substantial internal and radiative heating. 
In the following, the worst hot and cold cases will be analysed, and expected temperatures of the experiment will be shown by referencing the expected pump operations schedule.}\par 
\DIFaddend 

\DIFaddbegin \subsection{\DIFadd{Thermal analysis background}}
\DIFadd{A detailed description is given in \aref{app:thermal_background}.
}


\subsection{\DIFadd{Cold case}}\label{sec: thermal cold case}
\DIFadd{The cold case is a steady state assumption, given environmental conditions at a specific altitude. The appropriate altitude is determined as the maximum of the air heat convection $Q_\mathrm{conv}(h)$ profile from the averaged temperature and pressure profiles given in \fref{fig:temp_press_profile}. It is assumed that no solar radiation is impacting the system (except \mbox{%DIFAUXCMD
\gls{ir} }\hskip0pt%DIFAUXCMD
radiation from the Earth). Ambient pressure is taken from \fref{fig:temp_press_profile}, while the ambient temperature is set to be $-80$°C. Notably, the tubes containing pressurised, hot air, are assumed to be perfect thermal insulators, which is far from accurate. However, as the heat flux provided by the heated air is very small, this is acceptable.}\par 
\DIFadd{To estimate impact of different operations and system designs, the analysis takes into account whether pumps are turned on or off, whether the experiment has heaters or not and whether the experiment is insulated or not. The insulation is assumed to be a 5 cm thick box of styrofoam, completely encapsulating the chassis. The pumps are assumed to draw 50 W. This corresponds to the maximum power draw possible by two diaphragm pumps and one piston compressor. The active pumps were chosen according to the operations schedule as outlined in \sref{sec: pump system arch}. Assuming an efficiency of 60\% for the diaphragm pumps and 40\% for the piston compressor, the resulting heat load is 26 W.}\par 

\DIFadd{In \tref{tab:thermal results cold case} the results of the cold case analysis are displayed. 
Without the styrofoam insulation, the chassis reaches dangerously low temperatures, even with an active heating system. The heating system then would require over 75 W, which is unfeasible. Even with the styrofoam insulation, the required heating power is around 22 W, which is unsustainable for longer periods of time. However, when the pumps are being used, their heat alone is enough to keep the chassis and chamber within desired temperatures. The cold case analysis therefore concludes a necessity of insulation and active pumps.
}


\DIFaddend \begin{table}[H]
\centering
\DIFdelbeginFL %DIFDELCMD < \caption{%
{%DIFAUXCMD
\DIFdelFL{Component Temperature and Power Ratings}}
%DIFAUXCMD
%DIFDELCMD < \label{tab:component_specs}
%DIFDELCMD < \begin{tabular}{p{4cm}p{1cm}p{1cm}p{1cm}p{1cm}p{1cm}p{1cm}}
%DIFDELCMD < \toprule
%DIFDELCMD < \textbf{Component} &
%DIFDELCMD < \multicolumn{2}{c}{\makecell[c]{Operating T (°C)}} &
%DIFDELCMD < \multicolumn{2}{c}{\makecell[c]{Survivable T (°C)}} &
%DIFDELCMD < \multicolumn{2}{c}{\makecell[c]{Power Dissipated (W)}} \\
%DIFDELCMD < \cmidrule(lr){2-3} \cmidrule(lr){4-5} \cmidrule(lr){6-7}
%DIFDELCMD < & Min & Max & Min & Max & Mean & Max \\
%DIFDELCMD < \midrule
%DIFDELCMD < Pump &  &  &  &  &  &    \\
%DIFDELCMD < Valve &  &  &  &  &  &    \\
%DIFDELCMD < Arduino Mega & -40 & 85 &  &  &  &  1.95  \\
%DIFDELCMD < Switching Regulator & -40 & 125 &  &  &  & 39   \\
%DIFDELCMD < SD Card & -25 & 85 &  &  &  &    \\
%DIFDELCMD < K96 sensor (\ac{ndir}) & 0 & 40 & -40 & 85 &  & 2.447   \\
%DIFDELCMD < BME sensor & -40 & 85 &  &  &  & 0.001   \\
%DIFDELCMD < MS5607 (P) & -40 & 85 &  &  &  & 0.00004   \\
%DIFDELCMD < DS18B20 (T) & -55 & 125 &  &  &  &    \\
%DIFDELCMD < SHT3X-DIS (RH) & -40 & 125 &  &  &  &    \\
%DIFDELCMD < \bottomrule
%DIFDELCMD < \end{tabular}
%DIFDELCMD < %%%
\DIFdelendFL \DIFaddbeginFL \caption{\DIFaddFL{Thermal Performance Analysis: Cold Case Scenarios at 8.8\,km Altitude}}
\label{tab:thermal results cold case}
\begin{tabular}{@{}lll S[table-format=-2.2] S[table-format=-2.2] S[table-format=-2.2]@{}}
\toprule
\textbf{Insulation} & \textbf{Pumps} & \textbf{Thermal System} & {$T_{Chassis}$ [\si{\celsius}]} & {$T_{Chamber}$ [\si{\celsius}]} & {Power [\si{W}]} \\ \midrule
None & Off & None & -79.27 & -79.27 & {---} \\
None & Off & Active & -29.02 & 21.00 & 75.23 \\ 
Styrofoam & Off & Passive & -78.70 & -78.70 & {---} \\
Styrofoam & Off & Active & 6.52 & 21.00 & 21.79 \\ \midrule
None & Active & None & -61.80 & -61.80 & {---} \\
None & Active & Active & -20.40 & 21.00 & 62.28 \\ 
Styrofoam & Active & Passive & -72.41 & -72.41 & {---} \\
Styrofoam & Active & Active & 21.14 & 21.00 & -0.21 \\ \bottomrule
\addlinespace
\multicolumn{6}{l}{\small \textit{Note: Negative power values indicate active cooling required to maintain target temperature.}}
\end{tabular}
\DIFaddendFL \end{table}

\DIFdelbegin \DIFdel{The experiment can be thermally divided into three parts. The coldest part are the ambient sensors on the outside of the experiment box. They are }\DIFdelend \DIFaddbegin \subsection{\DIFadd{Hot case}}\label{sec: thermal hot case}
\DIFadd{Similar to the cold case, the hot case will follow a lumped system approach. However, the appropriate altitude here is chosen to minimize the impact of heat convection via air and maximizing solar heating. This means that the hot case will assume environmental conditions like those found at 30 km altitude, namely $T=-66$°C, $p=0.01$ bar and cloud-free skies. As in the cold case analysis, the effects of pumps, active heating, and styrofoam insulation are tested for individually. The results are displayed in \tref{tab:thermal results hot case}. Without styrofoam insulation, and no thermal system/heating or active pumps, the chassis temperature will reach around $-3$°C. 
Adding the insulation leads to a counterintuitive temperature reduction to $-47$°C. This is due to the styrofoam reflecting more sunlight than the aluminium, while simultaneously loosing heat more efficiently via \mbox{%DIFAUXCMD
\gls{ir} }\hskip0pt%DIFAUXCMD
radiation. 
As the heat loads involved are comparatively small, an active heating system has to provide around 14 W to raise the chamber temperature to 21°C. However, if pumps are active at 30 km altitude, the chassis will quickly reach temperatures higher than 21°C, requiring cooling of around 9 W instead.}\par 

\DIFadd{In conclusion, the hot case analysis shows that active pumps at high altitudes will necessitate a cooling system. In order to avoid complexity, it is therefore favourable to avoid using the pumps at high altitudes. 
This agrees well with the planned operations for the pumps, as given in \sref{sec: pump system arch}.
}

 


\begin{table}[H]
\centering
\caption{\DIFaddFL{Thermal Performance Analysis: Hot Case Scenarios at 30\,km Altitude}}
\label{tab:thermal results hot case}
\begin{tabular}{@{}lll S[table-format=-2.2] S[table-format=-2.2] S[table-format=-2.2]@{}}
\toprule
\textbf{Insulation} & \textbf{Pumps} & \textbf{Thermal System} & {$T_{Chassis}$ [\si{\celsius}]} & {$T_{Chamber}$ [\si{\celsius}]} & {Power [\si{W}]} \\ \midrule
None & Off & None & -3.16 & -3.16 & {---} \\
None & Off & Active & 16.39 & 21.00 & 6.93 \\ 
Styrofoam & Off & Passive & -47.32 & -47.32 & {---} \\
Styrofoam & Off & Active & 11.89 & 21.00 & 13.71 \\ \midrule
None & Active & None & 65.10 & 65.10 & {---} \\
None & Active & Active & 30.25 & 21.00 & -13.92 \\ 
Styrofoam & Active & Passive & -29.47 & -29.47 & {---} \\
Styrofoam & Active & Active & 26.75 & 21.00 & -8.65 \\ \bottomrule
\addlinespace
\multicolumn{6}{l}{\small \textit{Note: Negative power values indicate active cooling required to maintain desired chamber temperature.}}
\end{tabular}
\end{table}


\subsection{\DIFadd{Continuous operation}}
\DIFadd{As before, the continous operation usses a lumped system approach. However, now the total energy of the system is not assumed to be constant, but transient. This means that the system temperature is changing with the changing environmental and internal heat loads. Therefore, thermal inertia needs to be taken into account via
}\begin{align}
    \DIFadd{mc_{p}\dot{T} =}& \DIFadd{Q_\mathrm{ele}+Q_\mathrm{pump}(h)+Q_\mathrm{transfer} + Q_\mathrm{rad\,in}(h) + Q_\mathrm{insulation} - Q_\mathrm{conv}(h) - Q_\mathrm{rad\,out} \label{eq: heat diff chassis}
}\end{align}
\DIFadd{for the chassis. As before, the power required to keep the pressure chamber at 21°C can be calculated via \eqref{eq:heat_balance_chamber}, since $\dot{T}$ is supposed to be equal to 0 for the chamber. 
In this scenario the pumps operate by the operating plan specified in \sref{sec: pump system arch}. Additionally, pumps and thermal systems are activated 30 minutes before launch to simulate pre-heating }\DIFaddend and \DIFdelbegin \DIFdel{should be completely exposed and therefore are required to withstand temperaturesas low as -80}%DIFDELCMD < \degree %%%
\DIFdelend \DIFaddbegin \DIFadd{pre-pressurisation of the chamber. Otherwise the pumps are assumed to work on lower power draw settings at altitudes under 14 km, adding an estimated heat load of 18 W. Between 14 km and 22 km, the pumps operate at full power, adding a heat load of 26 W. Beyond 22 km the pumps are switched off. The total flight time is estimated to be 6 h long. 
The float phase is assumed to be at constant 26 km altitude. The heating system is assumed to be capable of a maximum heating output of 10 W.
After the pumps are deactivated at 22 km, the heating system aims to stabilise the chamber temperature at $-5$°}\DIFaddend C. \DIFdelbegin \DIFdel{This poses a challenge since most commercially available components are made for ground level conditions and manufacturers usually draw the line at -40}%DIFDELCMD < \degree %%%
\DIFdel{Cfor the temperature conditions. 
If we should get access to a thermal chamber , we can test the reliability of those sensors ourselves at lower temperatures. Since we haven't found sensors with operating ranges going down as low as -80}%DIFDELCMD < \degree %%%
\DIFdel{C we might have to resort to this attempt. 
Those sensors are not central to our experiment. The second part is the electronics compartment, which will likely be not actively heated. The bigger issue for this part might be overheating instead and a careful thermal analysis during the design process will be required to decide on mitigation steps if the temperature ranges of the components are exceeded at the upper or possibly still lower bounds. Later, in the integration process, further tests will follow to verify the functionality of that part at high altitudes. }\DIFdelend \DIFaddbegin \DIFadd{This is done as a proxy to prevent the electronics in the chassis from freezing.}\par 
\DIFaddend 

The \DIFdelbegin \DIFdel{last part is }\DIFdelend \DIFaddbegin \DIFadd{result for an experiment inside a styrofoam insulation box and subject to solar daylight is shown in \fref{fig:continous thermal sim main}. 
%DIF > As expected, heating is especially necessary between ca 8 km and 14 km. 
The chassis temperature never drops much below $-10$°C, while the thermal power requirements do not reach the set maximum of 10 W.  
By integrating the power requirement curve one obtains a total required heating power of $\approx 29.45$ Wh. This result refers to heating energy and not to }\DIFaddend the \DIFdelbegin \DIFdel{measurement chamber of which the temperature needs to be more carefully regulated to ensure that the K96 \mbox{%DIFAUXCMD
\ac{ndir} }\hskip0pt%DIFAUXCMD
sensor}\DIFdelend \DIFaddbegin \DIFadd{energy necessary to power heating components}\DIFaddend , which is \DIFdelbegin \DIFdel{central to our experiment, functions properly.Temperature fluctuations can also lead to output drift and therefore we need }\DIFdelend \DIFaddbegin \DIFadd{expected to be greater.}\par 
\DIFadd{If no insulation were used, the total required heating power would jump to 37.45 Wh. If insulation were used, but no daylight were available, the heating power required would be 36.42 Wh. The temperature and heating curves for both cases are provided in \fref{fig:continous thermal sim side}. If no insulation is used, the heating system is unable }\DIFaddend to keep the \DIFdelbegin \DIFdel{temperature in a narrow range . At the same time, pumps, valves or inlet and outlet tube openings could freeze over if the humidity level is high.
Since we might fly through clouds with very high humidity levels it is essential that the relative humidity doesn't get too high. By increasing the temperature }\DIFdelend \DIFaddbegin \DIFadd{chamber temperature within the required temperature range of 15°C to 40°C during ascend. It also fails at keeping the chassis temperature within a safe range }\DIFaddend of \DIFaddbegin \DIFadd{-15 to 65°C during descend. On the other hand, if insulation is used, but the flight is during night time or particularly cloudy conditions, the effects are rather minimal. The heater would max out its power output and }\DIFaddend the \DIFdelbegin \DIFdel{air, condensation can be mitigated and freezing of pumps, valves and tubes should be prevented as well }\DIFdelend \DIFaddbegin \DIFadd{chassis temperature would drop to around $-15$°C during descend, which is just the lower limit of the safe range and still 10°C away from the strict lower limit of -25°C.
}

\begin{figure}[htbp]
    \centering
    \includegraphics[width=0.8\textwidth]{0_images/figures/thermal_analysis_plot_WITHins_WITHlight.png}
    \caption{\DIFaddFL{Expected temperature of the chassis and chamber and the heating power necessary to sustain 21°C in the chamber. The gray lines indicate the time intervals when the balloon is ascending, floating, and descending. The colors indicate different power settings of the pressurisation system: green - no pumps, yellow - low power, red - high power.}}
    \label{fig:continous thermal sim main}
\end{figure}

\begin{figure}[htbp]
\begin{subfigure}[c]{0.5\linewidth}
\includegraphics[width=\linewidth]{0_images/figures/thermal_analysis_plot_NOins_WITHlight.png}
\caption{\DIFaddFL{No insulation, with daylight.}}
\end{subfigure}
\hfill
\begin{subfigure}[c]{0.5\linewidth}
\includegraphics[width=\linewidth]{0_images/figures/thermal_analysis_plot_WITHins_NOlight.png}
\caption{\DIFaddFL{With insulation, no daylight}}
\end{subfigure}%DIF > 
\caption{\DIFaddFL{Expected temperature of the chassis and chamber and the heating power necessary to sustain 21°C in the chamber. The gray lines indicate the time intervals when the balloon is ascending, floating, and descending. The colors indicate different power settings of the pressurisation system: green - no pumps, yellow - low power, red - high power.}}
\label{fig:continous thermal sim side}
\end{figure}

\subsection{\DIFadd{Conclusion}}\label{sec: thermal conclusion}
\DIFadd{The cold and hot case analysis shows that the experiment if enclosed in a styrofoam insulation is able to stay within the required temperature range depending on the proper use of pumps. In the ascend phase, the pumps are generating enough heat to keep the chamber temperature at 21°C, even without additional heating elements. In fact the transientt analysis shows that temperatures of chassis and chamber might reach close to 40°C. During float phase, the chamber temperature is allowed to drop to $-5$°C. This is done in order to keep the electrical components within the chassis at around $-10$°C. The required total heating power for daytime and nighttime would be $\approx 29.45$ Wh and $\approx 36.42$ Wh, respectively. Depending on the final design of the electrical components, the temperature limit and thus required power might be lowered further.
}

\subsection{\DIFadd{Passive thermal control system}}
\DIFadd{The above analysis shows that white styrofoam is a good candidate for insulation. The chassis is planned to be completely sourrounded by 5 cm thick styrofoam insulation. Further analysis on other insulation materials and implementations, as well as on other passive control systems, will commence with the start of thesis project 6 and 7 (see \tref{tab:thesis}). Additionally, the intake air will be preheated before being fed to the diaphragm pumps. This can be done via a thermally conductive, solenoid shaped tubing housed inside the chassis. If the tubing has higher heat capacity than the incoming air, it can act as a heat sink, heating the incoming air to an acceptable temperature. The exact shape and length of the tube will be determined via tests}\DIFaddend .
\DIFdelbegin \DIFdel{Further tests will be conducted in later stages of experiment integration}\DIFdelend \DIFaddbegin 

\subsection{\DIFadd{Active thermal control system}}
\DIFadd{The above analysis shows that during ascend and before launch, the pumps alone will provide enough heating power to heat the chamber to around 40°C. This is on the edge of the acceptable temperature range as outlined in \tref{tab:design_req}, P.7. Further analysis of overheating mitigation methods, including pump duty cycles and other active thermal control methods will commence with the start of thesis project 6 (see \tref{tab:thesis}).}\par 

\DIFadd{During float and descend, the above analysis shows that a 10 W heater is sufficient to keep the electronics inside the chassis within acceptable limits, see \tref{tab:design_req}, D.1.3. There are multiple common choices for thermal control components at low wattages and high precision. These will be discussed below.
}\paragraph{\DIFadd{\mbox{%DIFAUXCMD
\gls{ptc} }\hskip0pt%DIFAUXCMD
ceramics}} \DIFadd{are commonly used in portable space heaters. They are used via resistive heating. However, due to their \mbox{%DIFAUXCMD
\gls{ptc}}\hskip0pt%DIFAUXCMD
, a rising internal temperature causes their resistivity to increase as well. This behaviour can be configured, such that a desired maximum temperature can be set at which the ceramic resistivity becomes too high for the temperature to rise further. They are therefore self regulating, adding additional security and requiring less controlling. They are also quite cheap. However, because they are made of ceramic, they can not easily fit inside a tight experiment}\DIFaddend . 

\DIFdelbegin \DIFdel{Insulation will play another essential role in the thermal design. The whole experiment should be as much thermally separated from the rest of }\DIFdelend \DIFaddbegin \paragraph{\DIFadd{Peltier devices}} \DIFadd{are heat pumps, that transfer heat from one side of the device to the other when an electrical current is applied. Upon reversal of the current polarity, the heat transfer flow is also reversed. These devices can therefore be used to cool and heat. This makes them great candidates for an active thermal control system that has to cool and heat components depending on mission phase. The downside is that they require heat sinks that the excess heat (or lack thereof) can be dumped to. They also require more sophisticated controlling software and circuits to be able to switch between cooling and heating.
}

\paragraph{\DIFadd{Polyimide film heating elements}} \DIFadd{consist of a thin conductive heating element sourrounded by polyimide film. They are very flexible, sturdy and commonly used in space environment heating solutions. Although they are not self regulating, they can be easily controlled using \mbox{%DIFAUXCMD
\gls{pid} }\hskip0pt%DIFAUXCMD
algorithms. They are unable to provide cooling, but offer a great surface area to enable uniform heating. Due to their flexibility, they are also easy to integrate into a system and do not require complex software or circuits.}\par  

\DIFadd{Because of the low heating power required, the small operating temperature expected, the ease of integration, its big surface area, and its low cost a polyimide film heating element is chosen. The heating element is included in \tref{tab:power_consumption}. It can be attached via \mbox{%DIFAUXCMD
\gls{psa} }\hskip0pt%DIFAUXCMD
and wrapped around the pressure chamber. Since the chamber is made out of aluminium, and aluminium is good thermal conductor, one heating pad should suffice. This will be confirmed by tests. Should tests show that there is indeed no need for heating the pressurisation chamber, the heater may be moved directly to }\DIFaddend the \DIFdelbegin \DIFdel{vehicle as possible. On one hand, this avoids furthercooling of the experimentand on }\DIFdelend \DIFaddbegin \DIFadd{temperature sensitive electrical components like the \mbox{%DIFAUXCMD
\gls{pcb}}\hskip0pt%DIFAUXCMD
.}\par 

\DIFadd{The heating pad will be controlled using a \mbox{%DIFAUXCMD
\gls{pid} }\hskip0pt%DIFAUXCMD
algorithm, as common for heating appliances. The algorithm will run on the \mbox{%DIFAUXCMD
\gls{mcu}}\hskip0pt%DIFAUXCMD
. The \mbox{%DIFAUXCMD
\gls{mcu} }\hskip0pt%DIFAUXCMD
then controls the heating power of the heating pad by adjusting its duty cycle by powering a \mbox{%DIFAUXCMD
\gls{mosfet}}\hskip0pt%DIFAUXCMD
. The exact \mbox{%DIFAUXCMD
\gls{pwm} }\hskip0pt%DIFAUXCMD
setting is to be determined by the final \mbox{%DIFAUXCMD
\gls{mcu} }\hskip0pt%DIFAUXCMD
design. To prevent overheating of }\DIFaddend the \DIFdelbegin \DIFdel{other it avoids impacting other experiment or the vehicle. The measurement chambershould be well insulated as well to not (further) heat up the electronics compartment and to avoid condensation of water on the walls of the chamber}\DIFdelend \DIFaddbegin \DIFadd{chamber, a \mbox{%DIFAUXCMD
\gls{tco} }\hskip0pt%DIFAUXCMD
unit is used, wired in series with the heating pad. It can be directly attached via \mbox{%DIFAUXCMD
\gls{psa} }\hskip0pt%DIFAUXCMD
to the outside of the pressure chamber, monitoring its temperature and breaking the circuit when its temperature reaches above 50°C}\DIFaddend .


%DIF >  This can lead to overheating of the electrical components if the heat from dissipated power accumulates. Therefore, it is essential to know how much power is dissipated and to take additional cooling measures in case that radiative cooling is insufficient.
\DIFaddbegin 

%DIF >  \begin{table}[H]
%DIF >  \centering
%DIF >  \caption{Component Temperature and Power Ratings}
%DIF >  \label{tab:component_specs}
%DIF >  \begin{tabular}{p{4cm}p{1cm}p{1cm}p{1cm}p{1cm}p{1cm}p{1cm}}
%DIF >  \toprule
%DIF >  \textbf{Component} &
%DIF >  \multicolumn{2}{c}{\makecell[c]{Operating T (°C)}} &
%DIF >  \multicolumn{2}{c}{\makecell[c]{Survivable T (°C)}} &
%DIF >  \multicolumn{2}{c}{\makecell[c]{Power Dissipated (W)}} \\
%DIF >  \cmidrule(lr){2-3} \cmidrule(lr){4-5} \cmidrule(lr){6-7}
%DIF >  & Min & Max & Min & Max & Mean & Max \\
%DIF >  \midrule
%DIF >  Pump &  &  &  &  &  &    \\
%DIF >  Valve &  &  &  &  &  &    \\
%DIF >  Arduino Mega & -40 & 85 &  &  &  &  1.95  \\
%DIF >  Switching Regulator & -40 & 125 &  &  &  & 39   \\
%DIF >  SD Card & -25 & 85 &  &  &  &    \\
%DIF >  K96 sensor (\gls{ndir}) & 0 & 40 & -40 & 85 &  & 2.447   \\
%DIF >  BME sensor & -40 & 85 &  &  &  & 0.001   \\
%DIF >  MS5607 (P) & -40 & 85 &  &  &  & 0.00004   \\
%DIF >  DS18B20 (T) & -55 & 125 &  &  &  &    \\
%DIF >  SHT3X-DIS (RH) & -40 & 125 &  &  &  &    \\
%DIF >  \bottomrule
%DIF >  \end{tabular}
%DIF >  \end{table}

%DIF >  The experiment can be thermally divided into three parts. The coldest part are the ambient sensors on the outside of the experiment box. They are and should be completely exposed and therefore are required to withstand temperatures as low as -80\degree C. This poses a challenge since most commercially available components are made for ground level conditions and manufacturers usually draw the line at -40\degree C for the temperature conditions. If we should get access to a thermal chamber, we can test the reliability of those sensors ourselves at lower temperatures. Since we haven't found sensors with operating ranges going down as low as -80\degree C we might have to resort to this attempt. Those sensors are not central to our experiment. The second part is the electronics compartment, which will likely be not actively heated. The bigger issue for this part might be overheating instead and a careful thermal analysis during the design process will be required to decide on mitigation steps if the temperature ranges of the components are exceeded at the upper or possibly still lower bounds. Later, in the integration process, further tests will follow to verify the functionality of that part at high altitudes.

%DIF >  The last part is the measurement chamber of which the temperature needs to be more carefully regulated to ensure that the K96 \gls{ndir} sensor, which is central to our experiment, functions properly. Temperature fluctuations can also lead to output drift and therefore we need to keep the temperature in a narrow range. At the same time, pumps, valves or inlet and outlet tube openings could freeze over if the humidity level is high. Since we might fly through clouds with very high humidity levels it is essential that the relative humidity doesn't get too high. By increasing the temperature of the air, condensation can be mitigated and freezing of pumps, valves and tubes should be prevented as well. Further tests will be conducted in later stages of experiment integration.

%DIF >  Insulation will play another essential role in the thermal design. The whole experiment should be as much thermally separated from the rest of the vehicle as possible. On one hand, this avoids further cooling of the experiment and on the other it avoids impacting other experiment or the vehicle. The measurement chamber should be well insulated as well to not (further) heat up the electronics compartment and to avoid condensation of water on the walls of the chamber.

%DIF > %%%%%%%%%%%%%%%%%%%%%%%%%%%%%%%%%%%%%%%%%%%%%%%%%%%%%%%%%%%%%%%%%%%%%%%%%%%%%%
\DIFaddend \section{Power \DIFdelbegin \DIFdel{System}%DIFDELCMD < \MBLOCKRIGHTBRACE
%DIFDELCMD < %%%
\DIFdel{The \mbox{%DIFAUXCMD
\ac{mirage} }\hskip0pt%DIFAUXCMD
}\DIFdelend \DIFaddbegin \DIFadd{design}}

\DIFadd{The \mbox{%DIFAUXCMD
\gls{mirage} }\hskip0pt%DIFAUXCMD
}\DIFaddend team expects to rely exclusively on the 28.8V SAFT MP176065 which is provided by the gondola.

The experiment \DIFdelbegin \DIFdel{estimates 0.35A from Bexus power unit for all the component excluding the }\DIFdelend \DIFaddbegin \DIFadd{will have an estimated average current of 0.80A and use 138.67Wh in total, see table 4.9, including the thermal system and the pressurization system. This is higher than previous estimations, mainly due to the }\DIFaddend thermal system, \DIFdelbegin \DIFdel{which is yet to be decided}\DIFdelend \DIFaddbegin \DIFadd{compressor and vacuum pumps, but this calculation is based on full capacity so it's expected to be lower than this}\DIFaddend . The peak current draw \DIFdelbegin \DIFdel{from our power distribution board is estimated to be 1.15 A and is calculated with max current draw of each components. The Arduino mega has a max current of 200 mA which is much higher tham what is expected during the experiment. The }\DIFdelend \DIFaddbegin \DIFadd{is high at 2.26A from the BEXUS side for ~30 minutes, see table 4.8, while the compressor and vacuum pumps are working at full capacity, which is close to the requirement D.11.1. 
Since the MCU PCB will be designed by thesis students it is difficult to estimate precisely the current draw for this component, however it was estimated to be 300mA. The }\DIFaddend K96 \DIFdelbegin \DIFdel{sensor }\DIFdelend has a max current of \DIFdelbegin \DIFdel{300 mA }\DIFdelend \DIFaddbegin \DIFadd{300mA }\DIFaddend but the expected average is \DIFdelbegin \DIFdel{much }\DIFdelend lower depending on the time per measurement\DIFdelbegin \DIFdel{, the time for flushing and pressurization . The pump is expected to be supplied with a maximum of 500 mA and the average current will depend on the time per measurement }\DIFdelend . 
\DIFdelbegin \DIFdel{The SDHC card has an average current of 20 mA and a maximum of 150 mA}\DIFdelend \DIFaddbegin 

\DIFadd{Before launch, the pressurization system, e.g. vacuum pumps and compressor currents will be limited using some sort of pre-charge circuit and will thus not have a major impact on the energy consumption as a whole. During the flight phase, the pressurization system will be activated in lesser capacity for $\approx 45$ minutes and in full capacity $\approx 30$ minutes in order to conserve energy for the rest of the experiment. After the measurement phase the pressurization system will be limited or completely shut off.
}

\DIFadd{The thermal system will not be activated while the pressurization system is active in order to not exceed current draw limitations and thus will only be operational as detailed in thermal design section}\DIFaddend .

\begin{table}[H]
\centering
\caption{Max current consumption}
\label{tab:current_consumption}
\DIFdelbeginFL %DIFDELCMD < \begin{tabular}{ll}
%DIFDELCMD < \toprule
%DIFDELCMD < \textbf{Components} & \textbf{Max Current Draw} \\
%DIFDELCMD < \midrule
%DIFDELCMD < Arduino Mega & 200mA \\
%DIFDELCMD < K96 & 300mA \\
%DIFDELCMD < MS5607 (pressure sensor) & 1.4 mA \\
%DIFDELCMD < DS18B20 (temperature sensor) & 1.5 mA \\
%DIFDELCMD < SHT3X-DIS (humidity sensor) & 1.7 \micro A \\
%DIFDELCMD < Pump & 500 mA \\
%DIFDELCMD < Thermal System & TBD \\
%DIFDELCMD < SDHC & 150 mA \\
%DIFDELCMD < \midrule
%DIFDELCMD < Total & 1152 mA \\
%DIFDELCMD < \bottomrule
%DIFDELCMD < \end{tabular}
%DIFDELCMD < %%%
\DIFdelendFL \DIFaddbeginFL \begin{tabular}{llllll}
\toprule
\textbf{Amount} & \textbf{Components}& \textbf{Voltage} & \textbf{Current} &\textbf{Power} &\textbf{Energy}\\
\midrule
1 & MCU & 8V & 300mA & 2.4W & 1.2Wh\\
1 & K96 & 8V & 300mA & 2.4 & 1.2Wh \\
2 & MS5607 & 3.3V& 1µA & 6.6µW & 3.3µWh \\
6 & SHT45 & 3.3V & 0.4µA & 7.92µW & 3.96µWh \\
1 & Compressor & 24V & 1.25A & 30W & 15Wh \\
2 & Vacuum pumps & 24V & 0.42A & 20W & 10Wh\\
1 & SDHC & 3.3V & 150 mA & 0.495W & 0.248Wh \\
\midrule
Total &&& 2.45A & 55.295W & 27.648Wh\\
Total (BEXUS) &&28.8V& 2.26A & 65.05W & 32.526Wh \\
\bottomrule
\end{tabular}
\DIFaddendFL \end{table}

\begin{table}[H]
\centering
\caption{Power consumption from the power distribution board and the Bexus power unit}
\label{tab:power_consumption}
\DIFdelbeginFL %DIFDELCMD < \begin{tabular}{lllll}
%DIFDELCMD < \toprule
%DIFDELCMD < \textbf{Component} & \textbf{Number of parts} & \textbf{Max current} & \text{Voltage} & \textbf{Power} \\
%DIFDELCMD < \midrule
%DIFDELCMD < Arduino Mega & 1 & 200 mA & 8 V & 1.6 W \\
%DIFDELCMD < K96 & 1 & 300 mA & 8 V & 2.4 W \\
%DIFDELCMD < MS5607 & 1 & 1.4 mA & 3.3 V & 4.6 mW \\
%DIFDELCMD < DS18B20 & 1 & 1.5 mA & 3.3 V & 4.95 mW \\
%DIFDELCMD < Thermal System & TBD &TBD &TBD & TBD \\
%DIFDELCMD < SDHC & 1 & 150 mA & 3.3 V & 0.495 W \\
%DIFDELCMD < \midrule
%DIFDELCMD < Total & & 1152mA & & 8.505W \\
%DIFDELCMD < Total Bexus Power & & 0.35A & 28.8V & 10.08W \\
%DIFDELCMD < \bottomrule
%DIFDELCMD < \end{tabular}
%DIFDELCMD < %%%
\DIFdelendFL \DIFaddbeginFL \begin{tabular}{L{1.6cm}llllll}
\toprule
\textbf{Duration} & \textbf{Amount} & \textbf{Component} & \textbf{Voltage} & \textbf{Current} & \textbf{Power} & \textbf{Total Energy} \\
\midrule
6h & 1 & MCU & 8V & 300mA & 2.4 W & 14.4Wh \\
1.25h & 1 & K96 & 8V & 300mA & 2.4 W & 3Wh \\
6h & 6 & SHT45 & 3.3 V & 0.4µA & 1.32 µW & 47.52µWh \\
1.25h & 1 & Compressor & 24V & 1.25A & 30W & 37.5Wh \\
1.25h & 2 & Vacuum Pump & 24V & 0.42A & 10W & 25Wh \\
6h & 2 & Pressure sensor & 3.3V & 1µA & 3.3µW & 39.6µWh \\
6h & 1 &SDHC &3.3V& 150mA & 0.495W & 2.97Wh \\
3.5h & 1 &Heating element &12V& 0.83A & 10W & 35Wh \\

\midrule
Average (Bexus) & & & 28.8V & 0.80A & 23.11W\\
Total Bexus Power & & & & & & 138.67Wh \\
\bottomrule
\end{tabular}
\DIFaddendFL \end{table}

\DIFaddbegin \DIFadd{Table 4.9 presents key components, amount of components, duration of use, voltage levels, currents and power. The table is representing from the experimenters perspective while the average and total BEXUS power is from the BEXUS power unit perspective. This calculations used the assumption of $85\%$ efficiency to calculate the power consumtion on the BEXUS power unit. 
%DIF > %%%%%%%%%%%%%%%%%%%%%%%%%%%%%%%%%%%%%%%%%%%%%%%%%%%%%%%%%%%%%%%%%%%%%%%%%%%%%%
}\DIFaddend \section{Software \DIFdelbegin \DIFdel{Design}\DIFdelend \DIFaddbegin \DIFadd{design}\DIFaddend }

The on-board software is responsible for collecting data from the \DIFdelbegin \DIFdel{K96 sensor core }\DIFdelend \DIFaddbegin \DIFadd{sensors }\DIFaddend and store it on the on-board SD card. It is also responsible for handling communication from and to ground, as well as the control of the internal \DIFdelbegin \DIFdel{pumping }\DIFdelend \DIFaddbegin \DIFadd{pressure }\DIFaddend and thermal systems \DIFaddbegin \DIFadd{and preventing condensation}\DIFaddend .

\DIFaddbegin \DIFadd{For the pressurisation system the software needs to be able to read out the pressure sensor in the measurement chamber and adjust the power of the pumps accordingly to maintain a constant pressure of 3 bar. The mechanical side of the system is explained above. Since the vacuum pumps prepressurise the air that is then compressed to 3 bar, the power of the vacuum pumps and the compressor need to be handled separately. Depending on whether the compressor needs to have a minimum flowrate to work, it might be possible that a pulsed operation is necessary for the compressor, which would put strain on the electrical system. The software needs to be able to adjust these different components accordingly based on the internal pressure sensor. }\\

\DIFadd{For the thermal control system the software needs to take sensor inputs from different places. It should be split into two different programs, where one is responsible for stabilizing the temperature in the measurement chamber at around room temperature. This also includes controlling potential active heating elements at the gas inlet of the outer enclosure and before the air at 3 bar enters the measurement chamber based on at least on the thermal sensor of the K96 but possibly also several thermal sensor in between. The other thermal control loop should monitor components that are at risk of overheating and components that could underheat. There might be also active heating elements at critical components like the MCU, which could be switched on if the temperatures drop dangerously low inside the electrical compartment. This control loop need to take the different temperature outputs of sensors that are close to the pumps and electrical components and if necessary take counter measures via controlling the heating elements or reducing the power of the pumps.}\\
\DIFadd{For humidity the role of the software is to detect to high humidity levels (>80\% r.H.). If the ambient sensors detect high levels of humidity outside the software should close the shutter and stop the pumps as a first step to prevent more air with high humidity levels from entering the measurement chamber without losing pressure. When the humidity levels are low enough, the temperature of the air taken into the measurement chamber should be increased slightly by increasing the temperature of the active heating elements until the internal humidity is lower as well.}\\


\DIFaddend The settings of those systems depend on the readings of the internal and external pressure and temperature sensors. Additionally, correct flight phases need to be engaged. This is done based on longer term changes in the recorded ambient pressure or via a command from the ground, with the command taking priority. A \DIFdelbegin \DIFdel{information flow and process flow block }\DIFdelend \DIFaddbegin \DIFadd{control flow and state }\DIFaddend diagram are provided in Figure \DIFdelbegin \DIFdel{\ref{fig:info_flow} and Figure \ref{fig:process_flow}}\DIFdelend \DIFaddbegin \DIFadd{\ref{fig:cont_flow} and Figure \ref{fig:state_diagram}}\DIFaddend .

\DIFaddbegin \DIFadd{The experiment will run on 2 main modes: Manual and Autonomous. Each of the main modes has a submode for standby and active measurements. When the experiment reaches the point where no more data needs to be collected, the system will move to standby mode unless a command has been given to do otherwise. The modes will also allow for autonomous work even through radio blackouts and loss of communication. If below 22\,km and a loss of connection occurs, the system will automatically switch to autonomous measurements mode; if however the experiment is above 22\,km and a loss of connection occurs then the system will switch to Autonomous standby mode that will ensure the experiment remains safe without exhausting any of the hardware.
}

\DIFaddend To guarantee operations \DIFdelbegin \DIFdel{multiple \mbox{%DIFAUXCMD
\ac{mcu} }\hskip0pt%DIFAUXCMD
will be included in the experiment with one being active while the other serving as a watchdog and can overtake operations if needed}\DIFdelend \DIFaddbegin \DIFadd{watchdog systems will be used to reset the operations if it encountered a fatal error.
The software system will work with different threads working asynchronously.
Data transfer between the threads will happen through producer-consumer queues that are thread-safe.
There will be a thread for reading sensor data, a thread of controlling hardware, a thread sending downlink packages and receiving packages from uplink, a thread to write data on the SD card and finally a thread to control everything with the aid of a watchdog}\DIFaddend .

All transmissions are saved to the on-board storage with timestamps. Similarly, all saved data is saved with timestamps, too. A table detailing the expected data sizes of data to be stored is given in Table \ref{tab:data_rates}. \DIFdelbegin \DIFdel{This estimates around 3 kbps of storeddata for a total duration of around }\DIFdelend \DIFaddbegin \DIFadd{In total 39856 bits per second (40 kbit) will be stored. This means that for the duration of }\DIFaddend 6 hours \DIFdelbegin \DIFdel{, giving a total storage requirement of 64.8 Mbit. Ideally the whole data is sent down in regular intervals once it is recorded. }\DIFdelend \DIFaddbegin \DIFadd{a storage of 108 MB (megabyte) is required to store all the data received from the sensors. The SD cards onboard are more than enough for this task. Some of this data will be downlinked every seconds for redundancy and tracking of the functionality of the experiment. }\DIFaddend Based on the sampling rates in Table \ref{tab:data_rates},\DIFdelbegin \DIFdel{an interval of every 30 seconds seems reasonable. Thus, every 30 seconds, 91.5 kbit need to be transmit, ideally within 10 seconds to provide room for sending checksums and possibly duplicates. Therefore, we estimate a normal data rate of 9.15 kbps}\DIFdelend \DIFaddbegin \DIFadd{the Elink thread can pull 100 datapoints from the telemetry and data queue to downlink each second, resulting in an estimated downlink rate of 32.8 kbps, below the experiment requirements (P.18) of 100 kbps}\DIFaddend . Assuming a maximum connection dropout duration of 10 minutes, this would likewise imply a build-up of data of \DIFdelbegin \DIFdel{1.8 }\DIFdelend \DIFaddbegin \DIFadd{19.68 }\DIFaddend Mbit. Transferring it over the next 5 minutes gives a data rate of \DIFdelbegin \DIFdel{60 kbps }\DIFdelend \DIFaddbegin \DIFadd{65.6 kbps which is feasible}\DIFaddend . With compression it is probable that only half the rates stated need to be transmitted, but a higher data rate could be used to send duplicates or checksums for error mitigation.

The \DIFaddbegin \DIFadd{frequencies are chosen based on the datasheets of the sensors. The SHT45 sensor has multiple modes, fast mode plus (FM+) was chosen for this mission. The }\DIFaddend K96 sensor \DIFaddbegin \DIFadd{relies on the speed of the NDIR sensor within with typically is below 10 Hz, whereas in the MS5607 sensor OSR 4096 mode was chosen. The \mbox{%DIFAUXCMD
\gls{rtc} }\hskip0pt%DIFAUXCMD
frequency was calculated to be equal to all the frequencies added. This is to consider the scenario where none of the sensors data readings collide and all the readings from the 3 sensors are collected at different timestamps, which will require one timestamp for each.
}

\DIFadd{The K96 sensor }\DIFaddend core requires a UART interface and relies on the Modbus protocol\DIFaddbegin \DIFadd{, while most other sensors use I2C}\DIFaddend . The communication protocol will be UDP.

\begin{table}[H]
\centering
\caption{\DIFdelbeginFL \DIFdelFL{Sampling }\DIFdelendFL \DIFaddbeginFL \DIFaddFL{Max sampling }\DIFaddendFL rates, sampling sizes, and data rates of data products to be stored}
\label{tab:data_rates}
\DIFdelbeginFL %DIFDELCMD < \begin{tabular}{p{7cm}p{2cm}p{2cm}p{2cm}}
%DIFDELCMD < \toprule
%DIFDELCMD < \textbf{Component} & \textbf{Sampling rate [Hz]} & \textbf{Sampling size [bit]} & \textbf{Data rate [bps]} \\
%DIFDELCMD < \midrule
%DIFDELCMD < K96 concentration readings & 0.05 & 192 & 9.6 \\
%DIFDELCMD < K96 housekeeping data & 1 & 96 & 96 \\
%DIFDELCMD < Internal pressure sensor & 10 & 32 & 320 \\
%DIFDELCMD < Internal temperature sensor & 10 & 32 & 320 \\
%DIFDELCMD < Internal humidity sensor & 10 & 32 & 320 \\
%DIFDELCMD < Ambient pressure sensor & 10 & 32 & 320 \\
%DIFDELCMD < Ambient temperature sensor & 10 & 32 & 320 \\
%DIFDELCMD < Ambient humidity sensor & 10 & 32 & 320 \\
%DIFDELCMD < General housekeeping data (current for e.g. pumps, data from thermal system, ...) & 1 & 384 & 384 \\
%DIFDELCMD < \midrule
%DIFDELCMD < Total & & & 3049.6 \\
%DIFDELCMD < \bottomrule
%DIFDELCMD < \end{tabular}
%DIFDELCMD < %%%
\DIFdelendFL \DIFaddbeginFL \begin{tabular}{p{7cm}p{2cm}p{2cm}p{2cm}}
\toprule
\textbf{Component} & \textbf{Max. sampling rate [Hz]} & \textbf{Sampling size [bit], includes CRC} & \textbf{Data rate [bps]} \\
\midrule
LPL (K96) & 10 & 32 & 320\\
SPL (K96) & 10 & 32 & 320\\
MPL (K96) & 10 & 32 & 320\\
Pressure (K96) & 10 & 16 & 160\\
Temperature (K96) & 10 & 16 & 160\\
Relative Humidity (K96) & 10 & 16 & 160\\
Internal pressure sensor (MS5607) & 110 & 28 & 3080 \\
Internal temperature sensor (SHT45) & 222 & 24 & 5328 \\
Internal humidity sensor (SHT45) & 222 & 24 & 5328 \\
Ambient pressure sensor (MS5607) & 110 & 28 & 3080 \\
Ambient temperature sensor (SHT45) & 222 & 24 & 5328 \\
Ambient humidity sensor (SHT45) & 222 & 24 & 5328 \\
Real Time Clock & 342 & 32 & 10944\\
% \dos{General housekeeping data} (current for e.g. pumps, data from thermal system, ...) & 1 & 384 & 384 \\
\midrule
Total & & & 39856 \\
\bottomrule
\end{tabular}
\DIFaddendFL \end{table}

\begin{figure}[H]
\centering
\DIFdelbeginFL %DIFDELCMD < \includegraphics[width=.9\textwidth]{Graphics/Figures/Software information flow diagram.png}
%DIFDELCMD < %%%
\DIFdelendFL \DIFaddbeginFL \includegraphics[width=.9\textwidth]{0_images/figures/control_flow.png}
\DIFaddendFL \caption{\DIFdelbeginFL \DIFdelFL{Information }\DIFdelendFL \DIFaddbeginFL \DIFaddFL{Control }\DIFaddendFL flow diagram}
\DIFdelbeginFL %DIFDELCMD < \label{fig:info_flow}
%DIFDELCMD < %%%
\DIFdelendFL \DIFaddbeginFL \label{fig:cont_flow}
\DIFaddendFL \end{figure}

\begin{figure}[H]
\centering
\DIFdelbeginFL %DIFDELCMD < \includegraphics[width=\textwidth]{Graphics/Figures/Software Process Flow Diagram.png}
%DIFDELCMD < %%%
\DIFdelendFL \DIFaddbeginFL \includegraphics[width=\textwidth]{0_images/figures/state_chart.png}
\DIFaddendFL \caption{\DIFdelbeginFL \DIFdelFL{Process flow }\DIFdelendFL \DIFaddbeginFL \DIFaddFL{State }\DIFaddendFL diagram}
\DIFdelbeginFL %DIFDELCMD < \label{fig:process_flow}
%DIFDELCMD < %%%
\DIFdelendFL \DIFaddbeginFL \label{fig:state_diagram}
\DIFaddendFL \end{figure}

%DIF > %%%%%%%%%%%%%%%%%%%%%%%%%%%%%%%%%%%%%%%%%%%%%%%%%%%%%%%%%%%%%%%%%%%%%%%%%%%%%%
\section{Ground \DIFdelbegin \DIFdel{Support Equipment}\DIFdelend \DIFaddbegin \DIFadd{support equipment}\DIFaddend }
\DIFdelbegin \DIFdel{The data will be primarily saved onboard; however for redundancy and system check, the team will use the E-link. ESRANGE is equiped with a LIDAR device, which could be used in order to measure the vertical methane concentration for later validation}\DIFdelend \DIFaddbegin 

\DIFadd{A \mbox{%DIFAUXCMD
\gls{gui} }\hskip0pt%DIFAUXCMD
will be developed for easy control of \mbox{%DIFAUXCMD
\gls{mirage} }\hskip0pt%DIFAUXCMD
during the mission. The \mbox{%DIFAUXCMD
\gls{gui} }\hskip0pt%DIFAUXCMD
will allow ground operators to send commands through uplink and visualize the downlink data. Additionally the team will require basic electronic equipment, chairs and a table to operate from and troubleshoot any issues before flight}\DIFaddend .

%DIF < \section{Changes from Former Flight (for Re-Flights only)}
%DIF > %%%%%%%%%%%%%%%%%%%%%%%%%%%%%%%%%%%%%%%%%%%%%%%%%%%%%%%%%%%%%%%%%%%%%%%%%%%%%%\newpage

\chapter{\DIFdelbegin \DIFdel{Experiment Verification and Testing}\DIFdelend \DIFaddbegin \DIFadd{EXPERIMENT VERIFICATION AND TESTING}\DIFaddend }

\section{Verification Matrix}

\begin{longtable}{L{0.7cm}L{5cm}L{1.5cm}L{2cm}L{2.5cm}L{2cm}}
\caption{Verification Matrix}
\label{tab:verification_matrix} \\

\toprule
\textbf{ID} & \textbf{Requirement text} & \textbf{Method} & \textbf{Reference} & \textbf{Status} & \textbf{Verification Result} \\
\midrule
\endfirsthead

\multicolumn{6}{c}{\tablename\ \thetable\ -- \textit{Continued from previous page}}\\
\toprule
\textbf{ID} & \textbf{Requirement text} & \textbf{Method} & \textbf{Reference} & \textbf{Status} & \textbf{Verification Result} \\
\midrule
\endhead

\midrule
\multicolumn{6}{r}{\textit{Continued on next page}}\\
\endfoot

\bottomrule
\endlastfoot

F.1. & The experiment shall measure the atmospheric \DIFdelbegin \DIFdel{methane, CO2 and H2O }\DIFdelend \DIFaddbegin \DIFadd{CH\textsubscript{4}, CO\textsubscript{2} and H\textsubscript{2}O }\DIFaddend concentrations of ambient air simultaneously with an off-the-shelf \DIFdelbegin \DIFdel{\mbox{%DIFAUXCMD
\ac{ndir} }\hskip0pt%DIFAUXCMD
}\DIFdelend \DIFaddbegin \DIFadd{\mbox{%DIFAUXCMD
\gls{ndir} }\hskip0pt%DIFAUXCMD
}\DIFaddend sensor. & \DIFdelbegin \DIFdel{T, R}\DIFdelend \DIFaddbegin \DIFadd{R, T }\DIFaddend & Test \DIFdelbegin \DIFdel{1, Test 2 }\DIFdelend \DIFaddbegin \DIFadd{3, Test 13 }\DIFaddend & To be done (\DIFdelbegin \DIFdel{earliest test in March}\DIFdelend \DIFaddbegin \DIFadd{June}\DIFaddend ) & Pending \\
F.2. & The measurement chamber shall be pressurised with ambient air. & \DIFdelbegin \DIFdel{T, I }\DIFdelend \DIFaddbegin \DIFadd{R, T }\DIFaddend & Test 2 \DIFdelbegin \DIFdel{, Test 3, Test 5 }\DIFdelend & To be done (\DIFdelbegin \DIFdel{earliest test in April}\DIFdelend \DIFaddbegin \DIFadd{June}\DIFaddend ) & Pending \\
F.3. & The air in the measurement chamber shall be exchanged\DIFdelbegin \DIFdel{continuously}\DIFdelend . & \DIFdelbegin \DIFdel{T,I }\DIFdelend \DIFaddbegin \DIFadd{R,T }\DIFaddend & Test 2, Test 3 \DIFdelbegin \DIFdel{, Test 5 }\DIFdelend & To be done (earliest test in \DIFdelbegin \DIFdel{April}\DIFdelend \DIFaddbegin \DIFadd{June}\DIFaddend ) & Pending \\
F.4. & The experiment shall measure pressure inside of the measurement chamber. & \DIFaddbegin \DIFadd{R,}\DIFaddend T & Test 2, Test \DIFdelbegin \DIFdel{3, Test 5 }\DIFdelend \DIFaddbegin \DIFadd{7 }\DIFaddend & To be done (earliest test in \DIFdelbegin \DIFdel{April}\DIFdelend \DIFaddbegin \DIFadd{June}\DIFaddend ) & Pending \\
F.5. & The experiment shall measure the temperature inside of the measurement chamber. & \DIFaddbegin \DIFadd{R,}\DIFaddend T & Test \DIFdelbegin \DIFdel{2, Test }\DIFdelend 3, Test \DIFdelbegin \DIFdel{5 }\DIFdelend \DIFaddbegin \DIFadd{8  }\DIFaddend & To be done (earliest test in \DIFdelbegin \DIFdel{April}\DIFdelend \DIFaddbegin \DIFadd{June}\DIFaddend ) & Pending \\
F.6. & The experiment shall stabilize the temperature inside the measurement chamber. & \DIFdelbegin \DIFdel{T, A}\DIFdelend \DIFaddbegin \DIFadd{A,T }\DIFaddend & Test \DIFdelbegin \DIFdel{2, Test 3, Test }\DIFdelend 5, Test \DIFdelbegin \DIFdel{13 }\DIFdelend \DIFaddbegin \DIFadd{8 }\DIFaddend & To be done (earliest test in \DIFdelbegin \DIFdel{April}\DIFdelend \DIFaddbegin \DIFadd{June}\DIFaddend ) & Pending \\
F.7. & The experiment shall measure the humidity inside of the measurement chamber. & \DIFaddbegin \DIFadd{R,}\DIFaddend T & Test \DIFdelbegin \DIFdel{2, Test }\DIFdelend 3, Test \DIFdelbegin \DIFdel{5 }\DIFdelend \DIFaddbegin \DIFadd{8 }\DIFaddend & To be done (earliest test in \DIFdelbegin \DIFdel{April}\DIFdelend \DIFaddbegin \DIFadd{June}\DIFaddend ) & Pending \\
F.8 & The experiment shall prevent the condensation of water inside the measurement chamber. & \DIFdelbegin \DIFdel{A}\DIFdelend \DIFaddbegin \DIFadd{R}\DIFaddend , T & Test \DIFdelbegin \DIFdel{2, Test 3 }\DIFdelend \DIFaddbegin \DIFadd{8 }\DIFaddend & To be done (\DIFdelbegin \DIFdel{earliest in April}\DIFdelend \DIFaddbegin \DIFadd{July}\DIFaddend ) & Pending \\
F\DIFdelbegin \DIFdel{.8. }%DIFDELCMD < & %%%
\DIFdel{The experiment should change between measurement times based on the flight phase of the balloon. }%DIFDELCMD < & %%%
\DIFdel{T }%DIFDELCMD < & %%%
\DIFdel{Test 5 }%DIFDELCMD < & %%%
\DIFdel{To be done (May) }%DIFDELCMD < & %%%
\DIFdel{Pending }%DIFDELCMD < \\
%DIFDELCMD < %%%
\DIFdel{F}\DIFdelend .9. & The experiment should measure ambient temperature outside. & \DIFdelbegin \DIFdel{T, R}\DIFdelend \DIFaddbegin \DIFadd{R,T }\DIFaddend & Test \DIFdelbegin \DIFdel{2, Test 3 }\DIFdelend \DIFaddbegin \DIFadd{8, Test 13 }\DIFaddend & To be done (earliest in \DIFdelbegin \DIFdel{April}\DIFdelend \DIFaddbegin \DIFadd{July}\DIFaddend ) & Pending \\
F.10. & The experiment should measure ambient pressure outside. & \DIFdelbegin \DIFdel{T, R}\DIFdelend \DIFaddbegin \DIFadd{R,T }\DIFaddend &  Test \DIFdelbegin \DIFdel{2, Test 3 }\DIFdelend \DIFaddbegin \DIFadd{7, Test 13 }\DIFaddend & To be done (earliest in \DIFdelbegin \DIFdel{April}\DIFdelend \DIFaddbegin \DIFadd{July}\DIFaddend ) & Pending \\
F.11. & The experiment should measure ambient humidity outside. & \DIFdelbegin \DIFdel{T, R}\DIFdelend \DIFaddbegin \DIFadd{R,T }\DIFaddend &  Test \DIFdelbegin \DIFdel{2, Test 3 }\DIFdelend \DIFaddbegin \DIFadd{8, Test 13 }\DIFaddend & To be done (earliest in \DIFdelbegin \DIFdel{April}\DIFdelend \DIFaddbegin \DIFadd{July}\DIFaddend ) & Pending \\
F.12. & The experiment shall send data via E-Link. & T & Test 6 & To be done (\DIFdelbegin \DIFdel{May}\DIFdelend \DIFaddbegin \DIFadd{June}\DIFaddend ) & Pending \\
F.13. & The experiment \DIFdelbegin \DIFdel{should }\DIFdelend \DIFaddbegin \DIFadd{shall }\DIFaddend receive data via E-Link. & T & Test 6 & To be done (\DIFdelbegin \DIFdel{May}\DIFdelend \DIFaddbegin \DIFadd{June}\DIFaddend ) & Pending \\
P.1.1 & The \DIFdelbegin \DIFdel{\mbox{%DIFAUXCMD
\ac{ndir} }\hskip0pt%DIFAUXCMD
}\DIFdelend \DIFaddbegin \DIFadd{\mbox{%DIFAUXCMD
\gls{ndir} }\hskip0pt%DIFAUXCMD
}\DIFaddend sensor shall measure the concentration of methane with a range of 0.4 to 3.0 ppm. & \DIFdelbegin \DIFdel{T, R}\DIFdelend \DIFaddbegin \DIFadd{R, T }\DIFaddend & Test \DIFdelbegin \DIFdel{1 }\DIFdelend \DIFaddbegin \DIFadd{3, Test 13 }\DIFaddend & To be done (\DIFdelbegin \DIFdel{March}\DIFdelend \DIFaddbegin \DIFadd{June}\DIFaddend ) & Pending \\
P.1.2 \DIFdelbegin \DIFdel{. }\DIFdelend & The \DIFdelbegin \DIFdel{\mbox{%DIFAUXCMD
\ac{ndir} }\hskip0pt%DIFAUXCMD
}\DIFdelend \DIFaddbegin \DIFadd{\mbox{%DIFAUXCMD
\gls{ndir} }\hskip0pt%DIFAUXCMD
}\DIFaddend sensor shall measure the concentration of methane with an uncertainty of \DIFdelbegin \DIFdel{0.1 }\DIFdelend \DIFaddbegin \DIFadd{0.5 }\DIFaddend ppm or less. & \DIFdelbegin \DIFdel{A, R}\DIFdelend \DIFaddbegin \DIFadd{R, T }\DIFaddend & \DIFaddbegin \DIFadd{Test 3, Test 13 }\DIFaddend & \DIFaddbegin \DIFadd{To be done (June) }\DIFaddend & \DIFaddbegin \DIFadd{Pending}\DIFaddend \\
P.2.1 & The \DIFdelbegin \DIFdel{\mbox{%DIFAUXCMD
\ac{ndir} }\hskip0pt%DIFAUXCMD
}\DIFdelend \DIFaddbegin \DIFadd{\mbox{%DIFAUXCMD
\gls{ndir} }\hskip0pt%DIFAUXCMD
}\DIFaddend sensor shall measure the concentration of \DIFdelbegin \DIFdel{CO\textsubscript{2} }\DIFdelend \DIFaddbegin \DIFadd{carbondioxide }\DIFaddend with a range of 200 to 500 ppm. & \DIFdelbegin \DIFdel{T, R}\DIFdelend \DIFaddbegin \DIFadd{R, T }\DIFaddend & Test \DIFdelbegin \DIFdel{1 }\DIFdelend \DIFaddbegin \DIFadd{3, Test 13  }\DIFaddend & To be done (\DIFdelbegin \DIFdel{March}\DIFdelend \DIFaddbegin \DIFadd{June}\DIFaddend ) & Pending \\
P.2.2 & The \DIFdelbegin \DIFdel{\mbox{%DIFAUXCMD
\ac{ndir} }\hskip0pt%DIFAUXCMD
}\DIFdelend \DIFaddbegin \DIFadd{\mbox{%DIFAUXCMD
\gls{ndir} }\hskip0pt%DIFAUXCMD
}\DIFaddend sensor shall measure the concentration of \DIFdelbegin \DIFdel{CO\textsubscript{2} }\DIFdelend \DIFaddbegin \DIFadd{carbondioxide }\DIFaddend with an uncertainty of \DIFdelbegin \DIFdel{1ppm }\DIFdelend \DIFaddbegin \DIFadd{5 ppm }\DIFaddend or less. & \DIFdelbegin \DIFdel{A, R}\DIFdelend \DIFaddbegin \DIFadd{R, T }\DIFaddend & \DIFaddbegin \DIFadd{Test 3, Test 13 }\DIFaddend & \DIFaddbegin \DIFadd{To be done (June) }\DIFaddend & \DIFaddbegin \DIFadd{Pending}\DIFaddend \\
P.3.1 & The \DIFdelbegin \DIFdel{\mbox{%DIFAUXCMD
\ac{ndir} }\hskip0pt%DIFAUXCMD
}\DIFdelend \DIFaddbegin \DIFadd{\mbox{%DIFAUXCMD
\gls{ndir} }\hskip0pt%DIFAUXCMD
}\DIFaddend sensor shall measure the concentration of \DIFdelbegin \DIFdel{H\textsubscript{2}O }\DIFdelend \DIFaddbegin \DIFadd{water }\DIFaddend with a range of 1 to 30,\DIFdelbegin \DIFdel{000ppm}\DIFdelend \DIFaddbegin \DIFadd{000 ppm}\DIFaddend . & \DIFdelbegin \DIFdel{T, R}\DIFdelend \DIFaddbegin \DIFadd{R, T }\DIFaddend & Test \DIFdelbegin \DIFdel{1 }\DIFdelend \DIFaddbegin \DIFadd{3, Test 13  }\DIFaddend & To be done (\DIFdelbegin \DIFdel{March}\DIFdelend \DIFaddbegin \DIFadd{June}\DIFaddend ) & Pending \\
P.3.2 & The \DIFdelbegin \DIFdel{\mbox{%DIFAUXCMD
\ac{ndir} }\hskip0pt%DIFAUXCMD
}\DIFdelend \DIFaddbegin \DIFadd{\mbox{%DIFAUXCMD
\gls{ndir} }\hskip0pt%DIFAUXCMD
}\DIFaddend sensor shall measure the concentration of \DIFdelbegin \DIFdel{H\textsubscript{2}O }\DIFdelend \DIFaddbegin \DIFadd{water }\DIFaddend with an uncertainty of \DIFdelbegin \DIFdel{1ppm or lessat all altitudes}\DIFdelend \DIFaddbegin \DIFadd{5 ppm or less}\DIFaddend . & \DIFdelbegin \DIFdel{A, R}\DIFdelend \DIFaddbegin \DIFadd{R, T }\DIFaddend & \DIFaddbegin \DIFadd{Test 3, Test 13  }\DIFaddend & \DIFaddbegin \DIFadd{To be done (June) }\DIFaddend & \DIFaddbegin \DIFadd{Pending }\DIFaddend \\
P.4\DIFaddbegin \DIFadd{.1 }\DIFaddend & The pressure inside \DIFaddbegin \DIFadd{of }\DIFaddend the measurement chamber \DIFdelbegin \DIFdel{during methane measurements shall be at a stable pressure of 2-5 barduring measurements}\DIFdelend \DIFaddbegin \DIFadd{shall stay within a range of 2.5 and 3.2 bar}\DIFaddend .  & T & Test 2, Test \DIFdelbegin \DIFdel{5, Test 12 }\DIFdelend \DIFaddbegin \DIFadd{7 }\DIFaddend & To be done (earliest test in \DIFdelbegin \DIFdel{April}\DIFdelend \DIFaddbegin \DIFadd{June}\DIFaddend ) & Pending \\
P\DIFdelbegin \DIFdel{.5 }\DIFdelend \DIFaddbegin \DIFadd{.4.2 }\DIFaddend & The \DIFdelbegin \DIFdel{air }\DIFdelend \DIFaddbegin \DIFadd{pressure }\DIFaddend inside of the measurement chamber \DIFdelbegin \DIFdel{shall be exchanged within 2 minutes before each measurement}\DIFdelend \DIFaddbegin \DIFadd{should stay within a range of 2.8 and 3 bar}\DIFaddend . & T \DIFdelbegin \DIFdel{, A, R }\DIFdelend & Test 2, Test \DIFdelbegin \DIFdel{5 }\DIFdelend \DIFaddbegin \DIFadd{7 }\DIFaddend & To be done (earliest test in \DIFdelbegin \DIFdel{April}\DIFdelend \DIFaddbegin \DIFadd{June}\DIFaddend ) & Pending \\
P\DIFdelbegin \DIFdel{.6.1 }\DIFdelend \DIFaddbegin \DIFadd{.5 }\DIFaddend & The \DIFdelbegin \DIFdel{experiment shall measure the pressure inside the measurement chamber in a range of 300 to 1100 mbar}\DIFdelend \DIFaddbegin \DIFadd{pressurisation system shall have an air throughput of at least 1 l/min}\DIFaddend . & \DIFaddbegin \DIFadd{R, }\DIFaddend T & Test \DIFdelbegin \DIFdel{5 }\DIFdelend \DIFaddbegin \DIFadd{2 }\DIFaddend & To be done (\DIFdelbegin \DIFdel{May}\DIFdelend \DIFaddbegin \DIFadd{earliest test in June}\DIFaddend ) & Pending \\
P.6\DIFdelbegin \DIFdel{.2 }\DIFdelend \DIFaddbegin \DIFadd{.1 }\DIFaddend & The experiment shall measure the pressure inside \DIFdelbegin \DIFdel{the measurement chamber with a maximum uncertainty of 0.02 mbar. }%DIFDELCMD < & %%%
\DIFdel{R }%DIFDELCMD < & & & \\
%DIFDELCMD < %%%
\DIFdel{P.7.1 }%DIFDELCMD < & %%%
\DIFdel{The experiment shall measure the temperature inside the }\DIFdelend \DIFaddbegin \DIFadd{of the }\DIFaddend measurement chamber in a range of \DIFdelbegin \DIFdel{-40 to +85 }%DIFDELCMD < \degree %%%
\DIFdel{C}\DIFdelend \DIFaddbegin \DIFadd{0.3 to 4 bar}\DIFaddend . & R\DIFaddbegin \DIFadd{, T }\DIFaddend & \DIFaddbegin \DIFadd{Test 2, Test 7 }\DIFaddend & \DIFaddbegin \DIFadd{To be done (June) }\DIFaddend & \DIFaddbegin \DIFadd{Pending }\DIFaddend \\
P\DIFdelbegin \DIFdel{.7}\DIFdelend \DIFaddbegin \DIFadd{.6}\DIFaddend .2 & The experiment shall measure the \DIFdelbegin \DIFdel{temperature inside }\DIFdelend \DIFaddbegin \DIFadd{pressure inside of }\DIFaddend the measurement chamber with a maximum uncertainty of \DIFdelbegin \DIFdel{1.5}%DIFDELCMD < \degree %%%
\DIFdel{C}\DIFdelend \DIFaddbegin \DIFadd{0.02 bar}\DIFaddend . & R\DIFaddbegin \DIFadd{, T }\DIFaddend & \DIFaddbegin \DIFadd{Test 2, Test 7 }\DIFaddend & \DIFaddbegin \DIFadd{To be done (June) }\DIFaddend & \DIFaddbegin \DIFadd{Pending }\DIFaddend \\
P\DIFdelbegin \DIFdel{.8.1 }\DIFdelend \DIFaddbegin \DIFadd{.7 }\DIFaddend & The \DIFdelbegin \DIFdel{experiment shall measure the humidity inside }\DIFdelend \DIFaddbegin \DIFadd{temperature of the air inside of }\DIFaddend the measurement chamber \DIFdelbegin \DIFdel{in }\DIFdelend \DIFaddbegin \DIFadd{shall stay within }\DIFaddend a range of \DIFdelbegin \DIFdel{0 to 100\% r. H. }\DIFdelend \DIFaddbegin \DIFadd{15 and 40}\degree \DIFadd{C while active. }\DIFaddend & \DIFdelbegin \DIFdel{T, R }\DIFdelend \DIFaddbegin \DIFadd{A, T }\DIFaddend & Test \DIFdelbegin \DIFdel{1, Test 2, Test 3 }\DIFdelend \DIFaddbegin \DIFadd{5, Test 8 }\DIFaddend & To be done (\DIFdelbegin \DIFdel{earliest test in March}\DIFdelend \DIFaddbegin \DIFadd{June}\DIFaddend )& Pending\\
P.8 \DIFdelbegin \DIFdel{.2 }\DIFdelend & The \DIFdelbegin \DIFdel{experiment shall measure the humidity inside }\DIFdelend \DIFaddbegin \DIFadd{temperature of the air inside of }\DIFaddend the measurement chamber \DIFdelbegin \DIFdel{with a maximum uncertainty of 3\% r. H. }\DIFdelend \DIFaddbegin \DIFadd{should stay within a range of 19 and 24}\degree \DIFadd{C while active. }\DIFaddend & \DIFdelbegin \DIFdel{R }\DIFdelend \DIFaddbegin \DIFadd{A, T }\DIFaddend & \DIFaddbegin \DIFadd{Test 5, Test 8 }\DIFaddend & \DIFaddbegin \DIFadd{To be done (June)}\DIFaddend & \DIFaddbegin \DIFadd{Pending }\DIFaddend \\
P.9\DIFaddbegin \DIFadd{.1 }\DIFaddend & The \DIFdelbegin \DIFdel{temperature of the air inside the measurement chamber shall stay within }\DIFdelend \DIFaddbegin \DIFadd{experiment shall measure the temperature inside of the measurement chamber in }\DIFaddend a range of -40 \DIFdelbegin \DIFdel{and 60}\DIFdelend \DIFaddbegin \DIFadd{to +85 }\DIFaddend \degree C.  & \DIFdelbegin \DIFdel{A, }\DIFdelend R, T & Test \DIFdelbegin \DIFdel{3, Test 5 }\DIFdelend \DIFaddbegin \DIFadd{5, Test 8 }\DIFaddend & To be done (earliest test in \DIFdelbegin \DIFdel{May}\DIFdelend \DIFaddbegin \DIFadd{June}\DIFaddend ) & Pending \\
P\DIFdelbegin \DIFdel{.10 }\DIFdelend \DIFaddbegin \DIFadd{.9.2 }\DIFaddend & The experiment \DIFdelbegin \DIFdel{should keep the temperature of the air inside the measurement chamber within a range of 20 and 23}\DIFdelend \DIFaddbegin \DIFadd{shall measure the temperature inside of the measurement chamber with a maximum uncertainty of 1.5}\DIFaddend \degree C. & \DIFdelbegin \DIFdel{T, A }\DIFdelend \DIFaddbegin \DIFadd{R, T }\DIFaddend & Test \DIFdelbegin \DIFdel{3, Test 5 }\DIFdelend \DIFaddbegin \DIFadd{5, Test 8 }\DIFaddend & To be done (earliest test in \DIFdelbegin \DIFdel{May}\DIFdelend \DIFaddbegin \DIFadd{June}\DIFaddend ) & Pending \\
P\DIFdelbegin \DIFdel{.11 }\DIFdelend \DIFaddbegin \DIFadd{.10 }\DIFaddend & The experiment shall keep the humidity inside of the measurement chamber below 90\% r.H. & T \DIFdelbegin \DIFdel{, A }\DIFdelend & Test \DIFdelbegin \DIFdel{3 }\DIFdelend \DIFaddbegin \DIFadd{8 }\DIFaddend & To be done (\DIFdelbegin \DIFdel{May}\DIFdelend \DIFaddbegin \DIFadd{earliest test in July}\DIFaddend ) & Pending \\
P.11\DIFaddbegin \DIFadd{.1 }\DIFaddend & The experiment \DIFdelbegin \DIFdel{should measure the temperature, pressure and }\DIFdelend \DIFaddbegin \DIFadd{shall measure the }\DIFaddend humidity inside of the measurement chamber \DIFdelbegin \DIFdel{at a rate of 1Hz.}\DIFdelend \DIFaddbegin \DIFadd{in a range of 0 to 100\% r.H. }\DIFaddend & \DIFaddbegin \DIFadd{R, }\DIFaddend T & Test \DIFdelbegin \DIFdel{2, Test }\DIFdelend 3, Test \DIFdelbegin \DIFdel{5, Test 6 }\DIFdelend \DIFaddbegin \DIFadd{8 }\DIFaddend & To be done (earliest test in \DIFdelbegin \DIFdel{April}\DIFdelend \DIFaddbegin \DIFadd{June}\DIFaddend ) & Pending \\
P\DIFdelbegin \DIFdel{.12.1 }\DIFdelend \DIFaddbegin \DIFadd{.11.2 }\DIFaddend & The \DIFdelbegin \DIFdel{\mbox{%DIFAUXCMD
\ac{ndir} }\hskip0pt%DIFAUXCMD
sensor shall have a measurement time for the CH4, CO2 and H2O concentrations of 20s during ascent and descent.}\DIFdelend \DIFaddbegin \DIFadd{experiment shall measure the humidity inside of the measurement chamber with a maximum uncertainty of 3\% r.H. }\DIFaddend & \DIFaddbegin \DIFadd{R, }\DIFaddend T & Test \DIFdelbegin \DIFdel{2, Test 5 }\DIFdelend \DIFaddbegin \DIFadd{3, Test 8 }\DIFaddend & To be done (earliest test in \DIFdelbegin \DIFdel{April}\DIFdelend \DIFaddbegin \DIFadd{June}\DIFaddend ) & Pending \\
\DIFaddbegin 

\DIFaddend P.12 \DIFdelbegin \DIFdel{.2 }\DIFdelend & The \DIFdelbegin \DIFdel{\mbox{%DIFAUXCMD
\ac{ndir} }\hskip0pt%DIFAUXCMD
sensor shall have a measurement time for the CH4, CO2 and H2O concentrations of 50s during the floating phase}\DIFdelend \DIFaddbegin \DIFadd{experiment should measure the temperature, pressure and humidity inside of the measurement chamber at a rate over 10 Hz}\DIFaddend . & \DIFaddbegin \DIFadd{R, }\DIFaddend T & Test \DIFdelbegin \DIFdel{2, Test 5 }\DIFdelend \DIFaddbegin \DIFadd{8, Test 9 }\DIFaddend & To be done (earliest test in \DIFdelbegin \DIFdel{April}\DIFdelend \DIFaddbegin \DIFadd{July}\DIFaddend ) & Pending \\
P.13.1 & The experiment should measure the ambient pressure outside in a range of 20 to 1100 mbar. & \DIFdelbegin \DIFdel{T, R}\DIFdelend \DIFaddbegin \DIFadd{R, T }\DIFaddend & Test \DIFdelbegin \DIFdel{5 }\DIFdelend \DIFaddbegin \DIFadd{13 }\DIFaddend & To be done (\DIFdelbegin \DIFdel{May}\DIFdelend \DIFaddbegin \DIFadd{August}\DIFaddend ) & Pending \\
P.13.2 & The experiment should measure the ambient pressure outside with a maximum uncertainty of \DIFdelbegin \DIFdel{0.05 }\DIFdelend \DIFaddbegin \DIFadd{20 }\DIFaddend mbar. & R\DIFaddbegin \DIFadd{, T }\DIFaddend & \DIFaddbegin \DIFadd{Test 13 }\DIFaddend & \DIFaddbegin \DIFadd{To be done (August) }\DIFaddend & \DIFaddbegin \DIFadd{Pending  }\DIFaddend \\
P.14.1 & The experiment should measure the ambient temperature outside in a range of \DIFdelbegin \DIFdel{-80 }\DIFdelend \DIFaddbegin \DIFadd{-70 }\DIFaddend to +60 \degree C. & \DIFdelbegin \DIFdel{T, R}\DIFdelend \DIFaddbegin \DIFadd{R, T }\DIFaddend & Test \DIFdelbegin \DIFdel{2, Test 3 }\DIFdelend \DIFaddbegin \DIFadd{13, Test 8 }\DIFaddend & To be done (earliest test in \DIFdelbegin \DIFdel{April}\DIFdelend \DIFaddbegin \DIFadd{July}\DIFaddend ) & Pending \\
P.14.2 & The experiment should measure the ambient temperature outside with a maximum uncertainty of 1.5\degree C. & R\DIFaddbegin \DIFadd{, T }\DIFaddend & \DIFaddbegin \DIFadd{Test 13, Test 8 }\DIFaddend & \DIFaddbegin \DIFadd{To be done (earliest test in July) }\DIFaddend & \DIFaddbegin \DIFadd{Pending }\DIFaddend \\
P.15.1 & The experiment should measure the ambient humidity outside in a range of 0 to 100\DIFaddbegin \DIFadd{\% }\DIFaddend r.H. & R\DIFaddbegin \DIFadd{, T }\DIFaddend & \DIFaddbegin \DIFadd{Test 13, Test 8 }\DIFaddend & \DIFaddbegin \DIFadd{To be done (earliest test in July) }\DIFaddend & \DIFaddbegin \DIFadd{Pending }\DIFaddend \\
P.15.2 & The experiment should measure the \DIFdelbegin \DIFdel{humidity outside the sensor box }\DIFdelend \DIFaddbegin \DIFadd{ambient humidity outside }\DIFaddend with a maximum uncertainty of 3\% r.H. & R\DIFaddbegin \DIFadd{, T }\DIFaddend & \DIFaddbegin \DIFadd{Test 13, Test 8 }\DIFaddend & \DIFaddbegin \DIFadd{To be done (earliest test in July) }\DIFaddend & \DIFaddbegin \DIFadd{Pending }\DIFaddend \\
P.16 & The experiment should measure the ambient temperature, pressure and humidity outside at a rate \DIFdelbegin \DIFdel{of 1Hz}\DIFdelend \DIFaddbegin \DIFadd{over 10 Hz}\DIFaddend . & \DIFaddbegin \DIFadd{R, }\DIFaddend T & Test \DIFdelbegin \DIFdel{2, Test 6 }\DIFdelend \DIFaddbegin \DIFadd{13, Test 9 }\DIFaddend & To be done (earliest test in \DIFdelbegin \DIFdel{April}\DIFdelend \DIFaddbegin \DIFadd{July}\DIFaddend ) & Pending \\
P.17 & The experiment should receive data via E-Link with an uplink rate \DIFdelbegin \DIFdel{of 10kbps}\DIFdelend \DIFaddbegin \DIFadd{under 10 kbps}\DIFaddend . & \DIFaddbegin \DIFadd{R, }\DIFaddend T & Test \DIFdelbegin \DIFdel{6 }\DIFdelend \DIFaddbegin \DIFadd{12, Test 14 }\DIFaddend & To be done (\DIFdelbegin \DIFdel{May}\DIFdelend \DIFaddbegin \DIFadd{July}\DIFaddend ) & Pending \\
P.18 & The experiment shall send the obtained data with the E-Link with a downlink rate \DIFdelbegin \DIFdel{of }\DIFdelend \DIFaddbegin \DIFadd{under }\DIFaddend 100 kbps. & \DIFaddbegin \DIFadd{R, }\DIFaddend T & Test \DIFdelbegin \DIFdel{6 }\DIFdelend \DIFaddbegin \DIFadd{12, Test 14 }\DIFaddend & To be done (\DIFdelbegin \DIFdel{May}\DIFdelend \DIFaddbegin \DIFadd{July}\DIFaddend )& Pending \\
\DIFaddbegin 

\DIFaddend D.1\DIFaddbegin \DIFadd{.1 }\DIFaddend & The experiment shall operate in the temperature profile of the BEXUS vehicle flight and launch. & \DIFaddbegin \DIFadd{A, T }& \DIFadd{Test 8 }& \DIFadd{To be done (July) }& \DIFadd{Pending }\\
\DIFadd{D.1.2 }& \DIFadd{The temperature in the electronic compartment shall stay within a range of -25 and 85°C. }& \DIFadd{A, }\DIFaddend T \DIFdelbegin \DIFdel{, A, R }\DIFdelend & Test \DIFdelbegin \DIFdel{3 }\DIFdelend \DIFaddbegin \DIFadd{8 }\DIFaddend & \DIFaddbegin \DIFadd{To be done (July) }& \DIFadd{Pending}\\
\DIFadd{D.1.3 }& \DIFadd{The temperature in the electronic compartment should stay within a range of -15 and 60°C.}& \DIFadd{A, T }& \DIFadd{Test 8 }& \DIFaddend To be done (\DIFdelbegin \DIFdel{May}\DIFdelend \DIFaddbegin \DIFadd{July}\DIFaddend ) & Pending\\
D.2\DIFaddbegin \DIFadd{.1 }\DIFaddend & The experiment shall operate in the pressure profile of the BEXUS vehicle flight and launch. & \DIFdelbegin \DIFdel{T, A, R }\DIFdelend \DIFaddbegin \DIFadd{A, T }\DIFaddend & Test \DIFaddbegin \DIFadd{1, Test 2, Test }\DIFaddend 4, Test \DIFdelbegin \DIFdel{5, }\DIFdelend \DIFaddbegin \DIFadd{7 }& \DIFaddend Test \DIFdelbegin \DIFdel{12 }\DIFdelend \DIFaddbegin \DIFadd{1 completed successfully  }\DIFaddend & \DIFaddbegin \DIFadd{Pending }\\
\DIFadd{D.2.2 }& \DIFadd{The measurement chamber shall withstand a pressure difference of up to 4 bar.  }& \DIFadd{A, T }& \DIFadd{Test 7 }& \DIFaddend To be done (\DIFdelbegin \DIFdel{earliest test in April}\DIFdelend \DIFaddbegin \DIFadd{July}\DIFaddend ) & Pending \\
D.3 & The experiment shall operate in the vibration profile of the BEXUS vehicle flight and launch. & T & Test \DIFdelbegin \DIFdel{7 }\DIFdelend \DIFaddbegin \DIFadd{10 }\DIFaddend & To be done (July) & Pending \\
D.4 & The experiment \DIFaddbegin \DIFadd{shall operate in the shock profile of the BEXUS vehicle flight and launch. }& \DIFadd{T }& \DIFadd{Test 11 }& \DIFadd{To be done (July) }& \DIFadd{Pending }\\
\DIFadd{D.5 }& \DIFadd{The experiment }\DIFaddend power connector shall be an Amphenol PT02E8-4P, Code A. & \DIFdelbegin \DIFdel{I }\DIFdelend \DIFaddbegin \DIFadd{R }\DIFaddend & & & \\
D\DIFdelbegin \DIFdel{.5 }\DIFdelend \DIFaddbegin \DIFadd{.6 }\DIFaddend & The experiment \DIFdelbegin \DIFdel{shall }\DIFdelend \DIFaddbegin \DIFadd{data connector shall be an Amphenol RJF21B connector.  }& \DIFadd{R }& & & \\
\DIFadd{D.7 }& \DIFadd{The experiment shall }\DIFaddend not disturb or harm the launch vehicle.  & \DIFdelbegin \DIFdel{A, T }\DIFdelend \DIFaddbegin \DIFadd{R, A }\DIFaddend &  \DIFdelbegin \DIFdel{Test 8, Test 10, Test 7 }\DIFdelend &  \DIFdelbegin \DIFdel{To be done (earliest test in May) }\DIFdelend & Pending \\
D\DIFdelbegin \DIFdel{.6 }\DIFdelend \DIFaddbegin \DIFadd{.8 }\DIFaddend & \DIFdelbegin \DIFdel{The \mbox{%DIFAUXCMD
\ac{ndir} }\hskip0pt%DIFAUXCMD
sensor shall be separate from the other electronics in the measurement chamber of which the air will be exchanged continuously}\DIFdelend \DIFaddbegin \DIFadd{Only the measurement chamber containing the \mbox{%DIFAUXCMD
\gls{ndir} }\hskip0pt%DIFAUXCMD
sensor core and an additional pressure sensor shall be pressurised}\DIFaddend . & I\DIFaddbegin \DIFadd{, R }\DIFaddend & & & \\
D\DIFdelbegin \DIFdel{.7 }\DIFdelend \DIFaddbegin \DIFadd{.9 }\DIFaddend & The \DIFdelbegin \DIFdel{measurement chamber shall withstand the pressure difference between inside and outside over an extended period}\DIFdelend \DIFaddbegin \DIFadd{experiment shall consume less than 150 Wh during the flight}\DIFaddend . & T, A & Test \DIFdelbegin \DIFdel{5, Test 9, Test 12 }\DIFdelend \DIFaddbegin \DIFadd{7, Test 8, Test 14 }\DIFaddend & To be done (earliest test in \DIFdelbegin \DIFdel{May}\DIFdelend \DIFaddbegin \DIFadd{July}\DIFaddend ) & Pending \\
D\DIFdelbegin \DIFdel{.8 }\DIFdelend \DIFaddbegin \DIFadd{.10 }\DIFaddend & The experiment shall \DIFdelbegin \DIFdel{consume less than 150 Wh during the flight}\DIFdelend \DIFaddbegin \DIFadd{operate under 28.8\,V}\DIFaddend .  & \DIFaddbegin \DIFadd{R, }\DIFaddend T & Test \DIFdelbegin \DIFdel{5, Test 9 }\DIFdelend \DIFaddbegin \DIFadd{7, Test 8, Test 14 }\DIFaddend & To be done (earliest test in \DIFdelbegin \DIFdel{May}\DIFdelend \DIFaddbegin \DIFadd{July}\DIFaddend ) & Pending \\
D\DIFdelbegin \DIFdel{.9 }\DIFdelend \DIFaddbegin \DIFadd{.11.1 }\DIFaddend & The experiment shall have no currents exceeding \DIFaddbegin \DIFadd{2.4 A. }& \DIFadd{R, T }& \DIFadd{Test 7, Test 8, Test 14 }& \DIFadd{To be done (earliest test in July) }& \DIFadd{Pending }\\
\DIFadd{D.11.2 }& \DIFadd{The experiment should have no currents exceeding }\DIFaddend 1.4 A\DIFdelbegin \DIFdel{at any time during the flight}\DIFdelend .  & \DIFaddbegin \DIFadd{R, }\DIFaddend T & Test \DIFdelbegin \DIFdel{5, Test 9 }\DIFdelend \DIFaddbegin \DIFadd{7, Test 8, Test 14 }\DIFaddend & To be done (earliest test in \DIFdelbegin \DIFdel{May}\DIFdelend \DIFaddbegin \DIFadd{July}\DIFaddend ) & Pending \\
D\DIFdelbegin \DIFdel{.10 }\DIFdelend \DIFaddbegin \DIFadd{.12 }\DIFaddend & The experiment should stay within a size limit of \DIFdelbegin \DIFdel{20x20x20 }\DIFdelend \DIFaddbegin \DIFadd{15x20x20 }\DIFaddend cm\textsuperscript{3}. & \DIFdelbegin \DIFdel{I }\DIFdelend \DIFaddbegin \DIFadd{T }\DIFaddend & \DIFaddbegin \DIFadd{Test 9 }\DIFaddend & \DIFaddbegin \DIFadd{To be done (July) }\DIFaddend & \DIFaddbegin \DIFadd{Pending  }\DIFaddend \\
D\DIFdelbegin \DIFdel{.11 }\DIFdelend \DIFaddbegin \DIFadd{.13 }\DIFaddend & The experiment \DIFdelbegin \DIFdel{should not exceed a mass of 2 }\DIFdelend \DIFaddbegin \DIFadd{shall stay within a weight limit of 5 }\DIFaddend kg. & \DIFaddbegin \DIFadd{A,}\DIFaddend T & Test \DIFdelbegin \DIFdel{11 }\DIFdelend \DIFaddbegin \DIFadd{9 }\DIFaddend & \DIFaddbegin \DIFadd{To be done (July) }\DIFaddend & \DIFaddbegin \DIFadd{Pending }\DIFaddend \\
D\DIFdelbegin \DIFdel{.12 }\DIFdelend \DIFaddbegin \DIFadd{.14 }\DIFaddend & The air sample to be measured shall not be changed in its composition by the experiment, the vehicle or other experiments. & \DIFaddbegin \DIFadd{R, }\DIFaddend T & Test \DIFdelbegin \DIFdel{2 }\DIFdelend \DIFaddbegin \DIFadd{13 }\DIFaddend & To be done (\DIFdelbegin \DIFdel{April}\DIFdelend \DIFaddbegin \DIFadd{August}\DIFaddend ) & Pending \\
\DIFdelbegin \DIFdel{D.13 }\DIFdelend \DIFaddbegin 

\DIFadd{D.15 }\DIFaddend & The experiment shall have access to fresh air from outside. & \DIFaddbegin \DIFadd{R, }\DIFaddend I & & & \\
D\DIFdelbegin \DIFdel{.14 }\DIFdelend \DIFaddbegin \DIFadd{.16 }\DIFaddend & The data storage shall be \DIFaddbegin \DIFadd{able to store at least 2GB. }& \DIFadd{R }& & & \\
\DIFadd{D.17 }& \DIFadd{The data storage shall be }\DIFaddend protected from water due to condensation or during landing. & R & & & \\
D\DIFdelbegin \DIFdel{.15 }\DIFdelend \DIFaddbegin \DIFadd{.18 }\DIFaddend & The \DIFdelbegin \DIFdel{\mbox{%DIFAUXCMD
\ac{ndir} }\hskip0pt%DIFAUXCMD
}\DIFdelend \DIFaddbegin \DIFadd{\mbox{%DIFAUXCMD
\gls{ndir} }\hskip0pt%DIFAUXCMD
}\DIFaddend sensor should still function after recovery post flight. & T & Test 9, Test 10 & To be done (earliest test in \DIFdelbegin \DIFdel{June}\DIFdelend \DIFaddbegin \DIFadd{July}\DIFaddend ) & Pending \\
D\DIFdelbegin \DIFdel{.16 }\DIFdelend \DIFaddbegin \DIFadd{.19 }\DIFaddend & The budget for all components except the K96 sensor should stay below \DIFdelbegin \DIFdel{500}\DIFdelend \DIFaddbegin \DIFadd{2000}\DIFaddend \texteuro.  & R & & & \\
O.1 & The experiment shall be able to conduct measurements autonomously in case connection with the ground segment is lost. & T & Test \DIFdelbegin \DIFdel{2, Test 5, Test 9 }\DIFdelend \DIFaddbegin \DIFadd{12, Test 14 }\DIFaddend & To be done (earliest test in \DIFdelbegin \DIFdel{April}\DIFdelend \DIFaddbegin \DIFadd{July}\DIFaddend ) & Pending \\
O.2 & The experiment shall accept a request for radio silence at any time while on the launch pad.  & T & Test \DIFdelbegin \DIFdel{6, Test 9 }\DIFdelend \DIFaddbegin \DIFadd{12, Test 14 }\DIFaddend & To be done (earliest test in \DIFdelbegin \DIFdel{May}\DIFdelend \DIFaddbegin \DIFadd{July}\DIFaddend ) & Pending \\
O.3 & The experiment \DIFdelbegin \DIFdel{should switch automatically to a longer or shorter measurement time depending on the change in altitude. }%DIFDELCMD < & %%%
\DIFdel{T }%DIFDELCMD < & %%%
\DIFdel{Test 5 }%DIFDELCMD < & %%%
\DIFdel{To be done (May) }%DIFDELCMD < & %%%
\DIFdel{Pending }%DIFDELCMD < \\
%DIFDELCMD < %%%
\DIFdel{O.4 }%DIFDELCMD < & %%%
\DIFdel{The experiment shall respond to manual changes of measurement time from the ground. }%DIFDELCMD < & %%%
\DIFdel{T }%DIFDELCMD < & %%%
\DIFdel{Test 6, Test 9 }%DIFDELCMD < & %%%
\DIFdel{To be done (earliest test in May) }%DIFDELCMD < & %%%
\DIFdel{Pending }%DIFDELCMD < \\
%DIFDELCMD < %%%
\DIFdel{O.5 }%DIFDELCMD < & %%%
\DIFdel{The experiment shall }\DIFdelend \DIFaddbegin \DIFadd{shall }\DIFaddend store all data on board and continuously send \DIFdelbegin \DIFdel{condensed data sets }\DIFdelend \DIFaddbegin \DIFadd{data }\DIFaddend to ground via E-Link.  & T & Test \DIFdelbegin \DIFdel{6, Test 9 }\DIFdelend \DIFaddbegin \DIFadd{12 }\DIFaddend & To be done (\DIFdelbegin \DIFdel{earliest test in May}\DIFdelend \DIFaddbegin \DIFadd{July}\DIFaddend ) & Pending \\
O\DIFdelbegin \DIFdel{.6 }\DIFdelend \DIFaddbegin \DIFadd{.4 }\DIFaddend & The experiment shall automatically \DIFdelbegin \DIFdel{commence a full flushing cycle in the case of }\DIFdelend \DIFaddbegin \DIFadd{register and respond to }\DIFaddend high r.H. inside the chamber \DIFaddbegin \DIFadd{in order to lower the humidity}\DIFaddend . & T & Test \DIFdelbegin \DIFdel{3 }\DIFdelend \DIFaddbegin \DIFadd{8, Test 13 }\DIFaddend & To be done (\DIFdelbegin \DIFdel{May}\DIFdelend \DIFaddbegin \DIFadd{earliest test in July}\DIFaddend ) & Pending \\
\DIFaddbegin 

\DIFaddend \end{longtable}

\section{Verification Plan}

\DIFdelbegin %DIFDELCMD < \begin{table}[h]
%DIFDELCMD < %%%
\DIFdelendFL %DIF > \begin{table}[htbp]
%DIF > \centering
%DIF > \caption{Controlled functionality testing and calibration of K96 sensor %only}
%DIF > \label{tab:test1}
%DIF > \begin{tabular}{L{5cm}L{10cm}}
%DIF > \toprule
%DIF > Test number & 1 \\
%DIF > Test type & Sensor functionality \\
%DIF > Test facility & Uppsala University \\
%DIF > Tested item & K96 sensor \\
%DIF > Model & Prototype Model \\
%DIF > Procedure, Test level and duration & Performance test, compare sensor %readings to known composition of gases tested with varying humidity %levels for measurement times of 20, 50 and 100s. \\
%DIF > Test campaign duration & 2 days \\
%DIF > Test campaign date & March \\
%DIF > Test completed & No \\
%DIF > Requirements verified & \\
%DIF > \bottomrule
%DIF > \end{tabular}
%DIF > \end{table}
\DIFaddbeginFL \begin{table}[H]
\DIFaddendFL \centering
\caption{\DIFdelbeginFL \DIFdelFL{Controlled functionality testing and calibration of K96 sensor only}\DIFdelendFL \DIFaddbeginFL \DIFaddFL{Pressure test light bulb}\DIFaddendFL }
\DIFdelbeginFL %DIFDELCMD < \label{tab:test1}
%DIFDELCMD < \begin{tabular}{L{5cm}L{10cm}}
%DIFDELCMD < \toprule
%DIFDELCMD < Test number & 1 \\
%DIFDELCMD < Test type & Sensor functionality \\
%DIFDELCMD < Test facility & Uppsala University \\
%DIFDELCMD < Tested item & K96 sensor \\
%DIFDELCMD < Model & Prototype Model \\
%DIFDELCMD < Procedure, Test level and duration & Performance test, compare sensor readings to known composition of gases tested with varying humidity levels for measurement times of 20, 50 and 100s. \\
%DIFDELCMD < Test campaign duration & 2 days \\
%DIFDELCMD < Test campaign date & March \\
%DIFDELCMD < Test completed & No \\
%DIFDELCMD < Requirements verified & \\
%DIFDELCMD < \bottomrule
%DIFDELCMD < \end{tabular}
%DIFDELCMD < %%%
\DIFdelendFL \DIFaddbeginFL \label{tab:test_pressure_light_bulb}
\begin{tabular}{L{5cm}L{10cm}}
\toprule
Test number & 1 \\
Test type & Pressure test \\
Test facility & Uppsala University \\
Tested item & Light bulb of K96 sensor \\
Model & Prototype model \\
Procedure, Test level and duration & First part: Put Light bulb into vacuum chamber and drop pressure slowly to 10 mbar. \\
& Second part: Put light bulb into airtight container and increase pressure slowly until at least 4 bar are reached or until the light bulb bursts. \\
Test campaign duration & 2 days \\
Test campaign date & March \\
Test completed & Partially \\
Requirements verified & High pressure test sucessfully conducted\\
\bottomrule
\end{tabular}
\DIFaddendFL \end{table}

\begin{table}[H]
\centering
\caption{\DIFdelbeginFL \DIFdelFL{Measurement chamber functionality}\DIFdelendFL \DIFaddbeginFL \DIFaddFL{Test of pressurisation system}\DIFaddendFL }
\DIFdelbeginFL %DIFDELCMD < \label{tab:test2}
%DIFDELCMD < \begin{tabular}{L{5cm}L{10cm}}
%DIFDELCMD < \toprule
%DIFDELCMD < Test number & 2 \\
%DIFDELCMD < Test type & Thermal, pressure, sensor functionality \\
%DIFDELCMD < Test facility & ICOS, Norunda \\
%DIFDELCMD < Tested item & Measurement Chamber \\
%DIFDELCMD < Model & Prototype Model \\
%DIFDELCMD < Procedure, Test level and duration & Place measurement chamber close to atmospheric sensing facilities at the location and cross reference data by K96 collected without pressurisation and with flushing and pressurisation by pumps over longer time. Conduct with and without thermal control. Compare ambient sensor readings with measurement done by ICOS as well. \\
%DIFDELCMD < Test campaign duration & 4 days (1 day measurement chamber without pressurisation and thermal control, 1 day with pressurisation, 1 day with thermal control, 1 day full system test) \\
%DIFDELCMD < Test campaign date & April \\
%DIFDELCMD < Test completed & No \\
%DIFDELCMD < Requirements verified & \\
%DIFDELCMD < \bottomrule
%DIFDELCMD < \end{tabular}
%DIFDELCMD < %%%
\DIFdelendFL \DIFaddbeginFL \label{tab:test_pressurisation_system}
\begin{tabular}{L{5cm}L{10cm}}
\toprule
Test number & 2 \\
Test type & Pressure test \\
Test facility & Uppsala University, Vacuum chamber \\
Tested item & Measurement chamber without K96\\
Model & Prototype model \\
Procedure, Test level and duration & Slowly lower pressure from 1 bar to 10 mbar and observe stability of pressure inside of the box. Simulate balloon ascent in duration and rate of pressure decline.  \\
Test campaign duration & 1 day \\
Test campaign date & June \\
Test completed & No \\
Requirements verified & \\
\bottomrule
\end{tabular}
\DIFaddendFL \end{table}


\begin{table}[H]
\centering
\caption{\DIFdelbeginFL \DIFdelFL{Temperature testing}\DIFdelendFL \DIFaddbeginFL \DIFaddFL{Measurement chamber functionality}\DIFaddendFL }
\label{tab:test3}
\DIFdelbeginFL %DIFDELCMD < \begin{tabular}{L{5cm}L{10cm}}
%DIFDELCMD < \toprule
%DIFDELCMD < Test number & 3 \\
%DIFDELCMD < Test type & Thermal \\
%DIFDELCMD < Test facility & \ac{irf}, Thermal Chamber \\
%DIFDELCMD < Tested item & Whole experiment \\
%DIFDELCMD < Model & Proto flight model \\
%DIFDELCMD < Procedure, Test level and duration & Vary temperature and test functionality, 20 minutes at temperature steps of 5\degree C between -40\degree C and 10\degree C \\
%DIFDELCMD < Test campaign duration & 2 days (1 day build-up, 1 day testing and packing) \\
%DIFDELCMD < Test campaign date & May \\
%DIFDELCMD < Test completed & No \\
%DIFDELCMD < Requirements verified & \\
%DIFDELCMD < \bottomrule
%DIFDELCMD < \end{tabular}
%DIFDELCMD < %%%
\DIFdelendFL \DIFaddbeginFL \begin{tabular}{L{5cm}L{10cm}}
\toprule
Test number & 3 \\
Test type & Sensor functionality \\
Test facility & ICOS, Norunda \\
Tested item & Measurement Chamber \\
Model & Prototype Model \\
Procedure, Test level and duration & Put measurement chamber close to atmospheric sensing facilities at the location and cross reference collected data. Compare ambient sensor readings with measurements done by ICOS as well. \\
Test campaign duration & 2 days (setup and long term data collection) \\
Test campaign date & June \\
Test completed & No \\
Requirements verified & \\
\bottomrule
\end{tabular}
\DIFaddendFL \end{table}


\begin{table}[H]
\centering
\caption{Vacuum test of the electronic compartment}
\DIFdelbeginFL %DIFDELCMD < \label{tab:test2}
%DIFDELCMD < \begin{tabular}{L{5cm}L{10cm}}
%DIFDELCMD < \toprule
%DIFDELCMD < Test number & 4 \\
%DIFDELCMD < Test type & Pressure test \\
%DIFDELCMD < Test facility & Uppsala University, Vacuum chamber \\
%DIFDELCMD < Tested item & Electronic compartment \\
%DIFDELCMD < Model & Prototype model \\
%DIFDELCMD < Procedure, Test level and duration & Slowly lower pressure from 1 bar to 10 mbar and observe stability of circuit and temperature inside and outside of box. Simulate ascent in duration and rate of pressure decline. Keep at 10 mbar for 2 hours at least. \\
%DIFDELCMD < Test campaign duration & 1 day \\
%DIFDELCMD < Test campaign date & April \\
%DIFDELCMD < Test completed & No \\
%DIFDELCMD < Requirements verified & \\
%DIFDELCMD < \bottomrule
%DIFDELCMD < \end{tabular}
%DIFDELCMD < %%%
\DIFdelendFL \DIFaddbeginFL \label{tab:test4}
\begin{tabular}{L{5cm}L{10cm}}
\toprule
Test number & 4 \\
Test type & Pressure test \\
Test facility & Uppsala University, Vacuum chamber \\
Tested item & Electronic compartment \\
Model & Prototype model \\
Procedure, Test level and duration & Slowly lower pressure from 1 bar to 10 mbar and observe stability of circuit and temperature inside and outside of box. Simulate balloon ascent in duration and rate of pressure decline. Keep at 10 mbar for 2 hours at least. \\
Test campaign duration & 1 day \\
Test campaign date & June \\
Test completed & No \\
Requirements verified & \\
\bottomrule
\end{tabular}
\DIFaddendFL \end{table}


\begin{table}[H]
\centering
\caption{\DIFdelbeginFL \DIFdelFL{Vacuum }\DIFdelendFL \DIFaddbeginFL \DIFaddFL{Thermal control }\DIFaddendFL test\DIFdelbeginFL \DIFdelFL{of whole experiment}\DIFdelendFL }
\DIFdelbeginFL %DIFDELCMD < \label{tab:vacuum_test}
%DIFDELCMD < \begin{tabular}{L{5cm}L{10cm}}
%DIFDELCMD < \toprule
%DIFDELCMD < Test number & 5 \\
%DIFDELCMD < Test type & Pressure test \\
%DIFDELCMD < Test facility & Uppsala University, Vacuum chamber \\
%DIFDELCMD < Tested item & Whole experiment \\
%DIFDELCMD < Model & Protoflight model \\
%DIFDELCMD < Procedure, Test level and duration & Slowly lower pressure from 1 bar to 1 mbar and observe stability of circuit and temperature inside and outside of box. Simulate ascent in duration and rate of pressure decline. Keep at 1 mbar for at least 15 minutes. Perform cyclic pressurisation of measurement chamber during the whole time with measurement time corresponding to pressure level. \\
%DIFDELCMD < Test campaign duration & 2 days (setup and execution) \\
%DIFDELCMD < Test campaign date & May \\
%DIFDELCMD < Test completed & No \\
%DIFDELCMD < Requirements verified & \\
%DIFDELCMD < \bottomrule
%DIFDELCMD < \end{tabular}
%DIFDELCMD < %%%
\DIFdelendFL \DIFaddbeginFL \label{tab:test_thermal}
\begin{tabular}{L{5cm}L{10cm}}
\toprule
Test number & 5 \\
Test type & Thermal \\
Test facility & Thermal chamber and vacuum chamber \\
Tested item & Measurement chamber without K96 sensor \\
Model & Prototype model \\
Procedure, Test level and duration & In Thermal chamber: maintain room temperature of air in measurement chamber while decreasing temperature down to -40°C by autonomous current regulation. \\
& In Vacuum chamber: maintain room temperature of air in measurement chamber while decreasing pressure down to 10\,mbar by autonomous current regulation.  \\
Test campaign duration & 2 days \\
Test campaign date & June \\
Test completed & No \\
Requirements verified & \\
\bottomrule
\end{tabular}
\DIFaddendFL \end{table}




\begin{table}[H]
\centering
\caption{\DIFdelbeginFL \DIFdelFL{Software testing E-Link}\DIFdelendFL \DIFaddbeginFL \DIFaddFL{Static loads test}\DIFaddendFL }
\DIFdelbeginFL %DIFDELCMD < \label{tab:test6}
%DIFDELCMD < \begin{tabular}{L{5cm}L{10cm}}
%DIFDELCMD < \toprule
%DIFDELCMD < Test number & 6 \\
%DIFDELCMD < Test type & Software test \\
%DIFDELCMD < Test facility & Uppsala University \\
%DIFDELCMD < Tested item & E-link, EGSE and software, measurement chamber \\
%DIFDELCMD < Model & Prototype model \\
%DIFDELCMD < Procedure, Test level and duration & Send and receive data via E-link, test manual control of pumps, thermal control and measurement time \\
%DIFDELCMD < Test campaign duration & 1 day \\
%DIFDELCMD < Test campaign date & May \\
%DIFDELCMD < Test completed & No \\
%DIFDELCMD < Requirements verified & \\
%DIFDELCMD < \bottomrule
%DIFDELCMD < \end{tabular}
%DIFDELCMD < %%%
\DIFdelendFL \DIFaddbeginFL \label{tab:test_static_loads}
\begin{tabular}{L{5cm}L{10cm}}
\toprule
Test number & 6 \\
Test type & Mechanical test \\
Test facility & Uppsala University \\
Tested item & Main structure of experiment \\
Model & Qualification model \\
Procedure, Test level and duration & Mounting of experiment should simulate the gondola connection. Then the experiment is loaded with 10 to 30 times the experiment's own weight. \\
Test campaign duration & 1 day \\
Test campaign date & June (latest July) \\
Test completed & No \\
Requirements verified & \\
\bottomrule
\end{tabular}
\DIFaddendFL \end{table}



\begin{table}[H]
\centering
\caption{\DIFdelbeginFL \DIFdelFL{Vibration }\DIFdelendFL \DIFaddbeginFL \DIFaddFL{Vacuum }\DIFaddendFL test \DIFaddbeginFL \DIFaddFL{of whole experiment}\DIFaddendFL }
\DIFdelbeginFL %DIFDELCMD < \label{tab:test_vibration}
%DIFDELCMD < \begin{tabular}{L{5cm}L{10cm}}
%DIFDELCMD < \toprule
%DIFDELCMD < Test number & 7 \\
%DIFDELCMD < Test type & Mechanical test \\
%DIFDELCMD < Test facility & Outside \\
%DIFDELCMD < Tested item & Whole experiment \\
%DIFDELCMD < Model & Qualification model \\
%DIFDELCMD < Procedure, Test level and duration & Mount experiment in the back of a car and drive through rough terrain, inspect condition of experiment afterwards and test full functionality of experiment. \\
%DIFDELCMD < Test campaign duration & 2 days \\
%DIFDELCMD < Test campaign date & July \\
%DIFDELCMD < Test completed & No \\
%DIFDELCMD < Requirements verified & \\
%DIFDELCMD < \bottomrule
%DIFDELCMD < \end{tabular}
%DIFDELCMD < %%%
\DIFdelendFL \DIFaddbeginFL \label{tab:vacuum_test}
\begin{tabular}{L{5cm}L{10cm}}
\toprule
Test number & 7  \\
Test type & Pressure test \\
Test facility & Uppsala University, Vacuum chamber \\
Tested item & Whole experiment \\
Model & Protoflight model \\
Procedure, Test level and duration & Slowly lower pressure from 1 bar to 10 mbar and observe stability of circuit and temperature inside and outside of box. Simulate ascent in duration and rate of pressure decline. Keep at 10 mbar for at least 15 minutes. Pressurise measurement chamber during the whole time. \\
Test campaign duration & 2 days (setup and execution) \\
Test campaign date & July \\
Test completed & No \\
Requirements verified & \\
\bottomrule
\end{tabular}
\DIFaddendFL \end{table}


\begin{table}[H]
\centering
\caption{\DIFdelbeginFL \DIFdelFL{Static loads }\DIFdelendFL \DIFaddbeginFL \DIFaddFL{Temperature and humidity }\DIFaddendFL test \DIFaddbeginFL \DIFaddFL{of whole experiment}\DIFaddendFL }
\DIFdelbeginFL %DIFDELCMD < \label{tab:test_static_loads}
%DIFDELCMD < \begin{tabular}{L{5cm}L{10cm}}
%DIFDELCMD < \toprule
%DIFDELCMD < Test number & 8 \\
%DIFDELCMD < Test type & Mechanical test \\
%DIFDELCMD < Test facility & Uppsala University \\
%DIFDELCMD < Tested item & Main structure of experiment \\
%DIFDELCMD < Model & Qualification model \\
%DIFDELCMD < Procedure, Test level and duration & Mounting of experiment should simulate the gondola connection. Then the experiment is loaded with 10 to 30 times the experiment's own weight. \\
%DIFDELCMD < Test campaign duration & 1 day \\
%DIFDELCMD < Test campaign date & May (latest July) \\
%DIFDELCMD < Test completed & No \\
%DIFDELCMD < Requirements verified & \\
%DIFDELCMD < \bottomrule
%DIFDELCMD < \end{tabular}
%DIFDELCMD < %%%
\DIFdelendFL \DIFaddbeginFL \label{tab:test_temperature_humidity}
\begin{tabular}{L{5cm}L{10cm}}
\toprule
Test number & 8 \\
Test type & Thermal and humidity \\
Test facility & \gls{irf}, Thermal Chamber \\
Tested item & Whole experiment \\
Model & Proto flight model \\
Procedure, Test level and duration & Vary temperature and test functionality, 20 minutes at temperature steps of 5\degree C between -40\degree C and 10\degree C. Vary humidity between 5 and 100\% r.H. at each temperature step. \\
Test campaign duration & 3 days (1 day build-up, 2 day testing and packing) \\
Test campaign date & July \\
Test completed & No \\
Requirements verified & \\
\bottomrule
\end{tabular}
\DIFaddendFL \end{table}


\begin{table}[H]
\centering
\caption{\DIFdelbeginFL \DIFdelFL{Bench }\DIFdelendFL \DIFaddbeginFL \DIFaddFL{Weight and dimensions }\DIFaddendFL test}
\DIFdelbeginFL %DIFDELCMD < \label{tab:test_bench}
%DIFDELCMD < \begin{tabular}{L{5cm}L{10cm}}
%DIFDELCMD < \toprule
%DIFDELCMD < Test number & 9 \\
%DIFDELCMD < Test type & Long term test \\
%DIFDELCMD < Test facility & Uppsala University \\
%DIFDELCMD < Tested item & Whole experiment \\
%DIFDELCMD < Model & Flight Model \\
%DIFDELCMD < Procedure, Test level and duration & Assemble as for flight and run through simulated countdown of about 2 hours. Run experiment for 6 hours simulating ascent, float and descent. Simulate E-link dropouts as well. Consider recovery time as well for at least 24 hours. \\
%DIFDELCMD < Test campaign duration & 2 days \\
%DIFDELCMD < Test campaign date & August \\
%DIFDELCMD < Test completed & No \\
%DIFDELCMD < Requirements verified & \\
%DIFDELCMD < \bottomrule
%DIFDELCMD < \end{tabular}
%DIFDELCMD < %%%
\DIFdelendFL \DIFaddbeginFL \label{tab:weight_test}
\begin{tabular}{L{5cm}L{10cm}}
\toprule
Test number & 9 \\
Test type & Mechanical \\
Test facility & Uppsala University \\
Tested item & Whole experiment \\
Model & Qualification model \\
Procedure, Test level and duration & Put the experiment on a scale. Take a measurement tape and make sure that the outside dimensions are as psecified in the design. \\
Test campaign duration & 10 minutes \\
Test campaign date & July \\
Test completed & No \\
Requirements verified & \\
\bottomrule
\end{tabular}
\DIFaddendFL \end{table}


\begin{table}[H]
\centering
\DIFaddbeginFL \caption{\DIFaddFL{Vibration test}}
\label{tab:test_vibration}
\begin{tabular}{L{5cm}L{10cm}}
\toprule
Test number & 10 \\
Test type & Mechanical test \\
Test facility & Outside \\
Tested item & Whole experiment \\
Model & Qualification model \\
Procedure, Test level and duration & Mount experiment in the back of a car and drive through rough terrain, inspect condition of experiment afterwards and test full functionality of experiment. \\
Test campaign duration & 2 days \\
Test campaign date & July \\
Test completed & No \\
Requirements verified & \\
\bottomrule
\end{tabular}
\end{table}


\begin{table}[H]
\centering
\DIFaddendFL \caption{Schock test}
\DIFdelbeginFL %DIFDELCMD < \label{tab:test2}
%DIFDELCMD < \begin{tabular}{L{5cm}L{10cm}}
%DIFDELCMD < \toprule
%DIFDELCMD < Test number & 10 \\
%DIFDELCMD < Test type & Mechanical \\
%DIFDELCMD < Test facility & Uppsala University \\
%DIFDELCMD < Tested item & Whole experiment \\
%DIFDELCMD < Model & Qualification model \\
%DIFDELCMD < Procedure, Test level and duration & Drop experiment from 1 to 3 meters onto a padded surface and check visually and for full functionality afterwards. \\
%DIFDELCMD < Test campaign duration & 1 day \\
%DIFDELCMD < Test campaign date & June \\
%DIFDELCMD < Test completed & No \\
%DIFDELCMD < Requirements verified & \\
%DIFDELCMD < \bottomrule
%DIFDELCMD < \end{tabular}
%DIFDELCMD < %%%
\DIFdelendFL \DIFaddbeginFL \label{tab:test_shock}
\begin{tabular}{L{5cm}L{10cm}}
\toprule
Test number & 11 \\
Test type & Mechanical \\
Test facility & Uppsala University \\
Tested item & Whole experiment \\
Model & Qualification model \\
Procedure, Test level and duration & Drop experiment from 1 to 3 meters onto a padded surface and check visually and for full functionality afterwards. \\
Test campaign duration & 1 day \\
Test campaign date & July \\
Test completed & No \\
Requirements verified & \\
\bottomrule
\end{tabular}
\DIFaddendFL \end{table}


\begin{table}[H]
\centering
\caption{\DIFdelbeginFL \DIFdelFL{Weight test}\DIFdelendFL \DIFaddbeginFL \DIFaddFL{Software testing E-Link}\DIFaddendFL }
\DIFdelbeginFL %DIFDELCMD < \label{tab:weight_test}
%DIFDELCMD < \begin{tabular}{L{5cm}L{10cm}}
%DIFDELCMD < \toprule
%DIFDELCMD < Test number & 11 \\
%DIFDELCMD < Test type & Mechanical \\
%DIFDELCMD < Test facility & Uppsala University \\
%DIFDELCMD < Tested item & Whole experiment \\
%DIFDELCMD < Model & Qualification model \\
%DIFDELCMD < Procedure, Test level and duration & Put the experiment on a scale. \\
%DIFDELCMD < Test campaign duration & 5 minutes \\
%DIFDELCMD < Test campaign date & June \\
%DIFDELCMD < Test completed & No \\
%DIFDELCMD < Requirements verified & \\
%DIFDELCMD < \bottomrule
%DIFDELCMD < \end{tabular}
%DIFDELCMD < %%%
\DIFdelendFL \DIFaddbeginFL \label{tab:test6}
\begin{tabular}{L{5cm}L{10cm}}
\toprule
Test number & 12 \\
Test type & Software test \\
Test facility & Uppsala University \\
Tested item & E-link, EGSE and software, measurement chamber \\
Model & Whole experiment \\
Procedure, Test level and duration & Send and receive data via E-link, test manual control of pumps, thermal control and measurement time \\
Test campaign duration & 1 day \\
Test campaign date & July \\
Test completed & No \\
Requirements verified & \\
\bottomrule
\end{tabular}
\DIFaddendFL \end{table}


\begin{table}[H]
\centering
\caption{\DIFdelbeginFL \DIFdelFL{Pressure }\DIFdelendFL \DIFaddbeginFL \DIFaddFL{Functionality }\DIFaddendFL test \DIFdelbeginFL \DIFdelFL{light bulb}\DIFdelendFL \DIFaddbeginFL \DIFaddFL{with reference measurement}\DIFaddendFL }
\DIFdelbeginFL %DIFDELCMD < \label{tab:test_pressure_light_bulb}
%DIFDELCMD < \begin{tabular}{L{5cm}L{10cm}}
%DIFDELCMD < \toprule
%DIFDELCMD < Test number & 12 \\
%DIFDELCMD < Test type & Pressure test \\
%DIFDELCMD < Test facility & Uppsala University \\
%DIFDELCMD < Tested item & Light bulb of K96 sensor \\
%DIFDELCMD < Model & Prototype model \\
%DIFDELCMD < Procedure, Test level and duration & First part: Put Light bulb into vacuum chamber and drop pressure slowly to 1 mbar. \\
%DIFDELCMD < & Second part: Put light bulb into airtight container and increase pressure slowly until the pump limit or until the light bulb bursts. \\
%DIFDELCMD < Test campaign duration & 2 days \\
%DIFDELCMD < Test campaign date & March \\
%DIFDELCMD < Test completed & No \\
%DIFDELCMD < Requirements verified & \\
%DIFDELCMD < \bottomrule
%DIFDELCMD < \end{tabular}
%DIFDELCMD < %%%
\DIFdelendFL \DIFaddbeginFL \label{tab:test13}
\begin{tabular}{L{5cm}L{10cm}}
\toprule
Test number & 13 \\
Test type & Full referenced functionality test \\
Test facility & ICOS, Norunda \\
Tested item & Whole Experiment \\
Model & Flight Model \\
Procedure, Test level and duration & Put measurement chamber close to atmospheric sensing facilities at the location and cross reference collected data. Compare ambient sensor readings with measurements done by ICOS as well. \\
Test campaign duration & 2 days (setup and long term data collection) \\
Test campaign date & August \\
Test completed & No \\
Requirements verified & \\
\bottomrule
\end{tabular}
\DIFaddendFL \end{table}

\begin{table}[H]
\centering
\caption{\DIFdelbeginFL \DIFdelFL{Thermal control }\DIFdelendFL \DIFaddbeginFL \DIFaddFL{Bench }\DIFaddendFL test}
\DIFdelbeginFL %DIFDELCMD < \label{tab:test_thermal}
%DIFDELCMD < \begin{tabular}{L{5cm}L{10cm}}
%DIFDELCMD < \toprule
%DIFDELCMD < Test number & 13 \\
%DIFDELCMD < Test type & Thermal \\
%DIFDELCMD < Test facility & Thermal chamber and vacuum chamber \\
%DIFDELCMD < Tested item & Measurement chamber without K96 sensor \\
%DIFDELCMD < Model & Prototype model \\
%DIFDELCMD < Procedure, Test level and duration & In Thermal chamber: maintain room temperature of air in measurement chamber while decreasing temperature down to -40°C by autonomous current regulation. \\
%DIFDELCMD < & In Vacuum chamber: maintain room temperature of air in measurement chamber while decreasing pressure down to 5\,mbar by autonomous current regulation. Draw air from outside and put outlet also to outside of the vacuum chamber. \\
%DIFDELCMD < Test campaign duration & 2 days \\
%DIFDELCMD < Test campaign date & June \\
%DIFDELCMD < Test completed & No \\
%DIFDELCMD < Requirements verified & \\
%DIFDELCMD < \bottomrule
%DIFDELCMD < \end{tabular}
%DIFDELCMD < %%%
\DIFdelendFL \DIFaddbeginFL \label{tab:test_bench}
\begin{tabular}{L{5cm}L{10cm}}
\toprule
Test number & 14 \\
Test type & Long term test \\
Test facility & Uppsala University \\
Tested item & Whole experiment \\
Model & Flight Model \\
Procedure, Test level and duration & Assemble as for flight and run through simulated countdown of about 2 hours. Run experiment for 6 hours simulating ascent, float and descent. Simulate E-link dropouts as well. Consider recovery time as well for at least 24 hours. \\
Test campaign duration & 2 days \\
Test campaign date & August \\
Test completed & No \\
Requirements verified & \\
\bottomrule
\end{tabular}
\DIFaddendFL \end{table}


%DIF <  ... Additional test tables would follow similarly ...
%DIF >  % ... Additional test tables would follow similarly ...

\DIFdelbegin \section{\DIFdel{Verification Results}}
%DIFAUXCMD
\addtocounter{section}{-1}%DIFAUXCMD
\DIFdel{Test will be conducted later on and the results will be listed here in later versions of the SED.
%DIF < 
%DIF < \begin{table}[h]
%DIF < \centering
%DIF < \caption{Example of a test result overview}
%DIF < \label{tab:test_result}
%DIF < \begin{tabular}{L{5cm}L{10cm}}
%DIF < \toprule
%DIF < \textbf{Verification number} & \\
%DIF < \textbf{Type of test} & \\
%DIF < \textbf{Facility} & \\
%DIF < \textbf{Verified item} & \\
%DIF < \textbf{Verification description} & \\
%DIF < \textbf{Expected results} & \\
%DIF < \textbf{Obtained Results} & \\
%DIF < \textbf{Conclusions} & \\
%DIF < \bottomrule
%DIF < \end{tabular}
%DIF < \end{table}
%DIF < 
}\DIFdelend %DIF >  \section{Verification Results}
%DIF >  Test will be conducted later on and the results will be listed here in later versions of the SED.
%DIF >  %
%DIF >  %\begin{table}[htbp]
%DIF >  %\centering
%DIF >  %\caption{Example of a test result overview}
%DIF >  %\label{tab:test_result}
%DIF >  %\begin{tabular}{L{5cm}L{10cm}}
%DIF >  %\toprule
%DIF >  %\textbf{Verification number} & \\
%DIF >  %\textbf{Type of test} & \\
%DIF >  %\textbf{Facility} & \\
%DIF >  %\textbf{Verified item} & \\
%DIF >  %\textbf{Verification description} & \\
%DIF >  %\textbf{Expected results} & \\
%DIF >  %\textbf{Obtained Results} & \\
%DIF >  %\textbf{Conclusions} & \\
%DIF >  %\bottomrule
%DIF >  %\end{tabular}
%DIF >  %\end{table}
%DIF >  %\newpage

\chapter{\DIFdelbegin \DIFdel{Launch Campaign Preparation}\DIFdelend \DIFaddbegin \DIFadd{LAUNCH CAMPAIGN PREPARATIONS}\DIFaddend }
\DIFaddbegin \label{sec:campaign_preparations}
\section{\DIFadd{Input for the Campaign / Flight Requirement Plan}}
\label{sec:input_FRP}
\DIFaddend 

\DIFdelbegin \section{\DIFdel{Input for the Campaign / Flight Requirement Plans}}
%DIFAUXCMD
\addtocounter{section}{-1}%DIFAUXCMD
%DIFDELCMD < 

%DIFDELCMD < %%%
\DIFdelend %DIF > %%%%%%%%%%%%%%%%%%%%%%%%%%%%%%%%%%%%%%%%%%%%%%%%%%%%%%%%%%%%%%%%%%%%%%%%%%%%%%
\subsection{Dimensions and Mass}

\DIFdelbegin %DIFDELCMD < \begin{table}[h]
%DIFDELCMD < %%%
\DIFdelendFL \DIFaddbeginFL \begin{table}[htbp]
\DIFaddendFL \centering
\caption{Experiment Mass and Volume}
\label{tab:mass_volume}
\DIFdelbeginFL %DIFDELCMD < \begin{tabular}{L{6cm}L{6cm}}
%DIFDELCMD < \toprule
%DIFDELCMD < Experiment mass (in kg): & $\sim$1.45\\
%DIFDELCMD < Experiment footprint area (for BEXUS) (in $m^2$): & 0.01\\
%DIFDELCMD < Experiment volume (in $m^3$): & 0.0005\\
%DIFDELCMD < Experiment expected \ac{cog} position: & Distance from geometric centre: x=4mm y=23mm z=22mm \\
%DIFDELCMD < \bottomrule
%DIFDELCMD < \end{tabular}
%DIFDELCMD < %%%
\DIFdelendFL \DIFaddbeginFL \begin{tabular}{L{6cm}L{6cm}}
\toprule
Experiment mass (in kg): & $\sim$4\\
Experiment footprint area (for BEXUS) (in $m^2$): & 0.04\\
Experiment volume (in $m^3$): & 0.006\\
Experiment expected \gls{cog} position: & Distance from geometric centre: x=4mm y=23mm z=22mm \\
\bottomrule
\end{tabular}
\DIFaddendFL \end{table}
\clearpage 

%DIF > %%%%%%%%%%%%%%%%%%%%%%%%%%%%%%%%%%%%%%%%%%%%%%%%%%%%%%%%%%%%%%%%%%%%%%%%%%%%%%
\subsection{Ground and Flight Safety Risks}

\DIFdelbegin %DIFDELCMD < \begin{table}[h]
%DIFDELCMD < %%%
\DIFdelendFL \DIFaddbeginFL \begin{table}[htbp]
\DIFaddendFL \centering
\caption{Ground and flight safety risks}
\label{tab:safety_risks}
\DIFdelbeginFL %DIFDELCMD < \begin{tabular}{p{4cm} p{5cm} p{4cm}}
%DIFDELCMD < \toprule
%DIFDELCMD < \textbf{Risk} & \textbf{Key Characteristics} & \textbf{Mitigation} \\
%DIFDELCMD < \midrule
%DIFDELCMD < Use of a pressurised vessel with up to 5 bar difference beween outside and inside. & Material: Aluminium & Make appropriate simulations and test thoroughly before flight. \\
%DIFDELCMD < \bottomrule
%DIFDELCMD < \end{tabular}
%DIFDELCMD < %%%
\DIFdelendFL \DIFaddbeginFL \begin{tabular}{p{4cm} p{5cm} p{4cm}}
\toprule
\textbf{Risk} & \textbf{Key Characteristics} & \textbf{Mitigation} \\
\midrule
Use of a pressurised vessel with 3 bar difference beween outside and inside. & Material: Aluminium & Pressure test will be conducted up to 4 bar in line with the Esrange safety manual. \\
\bottomrule
\end{tabular}
\DIFaddendFL \end{table}
\DIFdelbegin %DIFDELCMD < 

%DIFDELCMD < %%%
\DIFdel{It's possible that our thermal control system will pose a safety risk by reaching above 60°C. We will verify until the next version (before CDR) whether we need mitigation methods.
}%DIFDELCMD < 

%DIFDELCMD < %%%
\DIFdelend \clearpage

%DIF > %%%%%%%%%%%%%%%%%%%%%%%%%%%%%%%%%%%%%%%%%%%%%%%%%%%%%%%%%%%%%%%%%%%%%%%%%%%%%%
\subsection{Electrical Interfaces}

\DIFdelbegin %DIFDELCMD < \begin{table}[h]
%DIFDELCMD < %%%
\DIFdelendFL \DIFaddbeginFL \begin{table}[htbp]
\DIFaddendFL \centering
\caption{Electrical interfaces applicable to BEXUS}
\label{tab:electrical_interfaces_bexus}
\DIFdelbeginFL %DIFDELCMD < \begin{tabular}{p{10cm} p{5cm}}
%DIFDELCMD < \toprule
%DIFDELCMD < \multicolumn{2}{l}{\textbf{BEXUS Electrical Interfaces}} \\
%DIFDELCMD < \midrule
%DIFDELCMD < E-Link Interface: E-Link required? Yes/No & Yes\\
%DIFDELCMD < Number of E-Link interfaces: & 1 \\
%DIFDELCMD < Number of required IP addresses: & 2 (one for the experiment and one for the ground station)\\
%DIFDELCMD < Data rate -- downlink (max. and average): & 100 Kbit/s \\
%DIFDELCMD < Data rate -- uplink (max. and average): & 10 Kbit/s \\
%DIFDELCMD < Interface type (RS-232, Ethernet): & Ethernet \\
%DIFDELCMD < \midrule
%DIFDELCMD < Power system: Gondola power required? Yes/No & Yes\\
%DIFDELCMD < Peak power and current consumption: & 10.08W and 0.35A \\
%DIFDELCMD < Average power and current consumption: & 8.64W and 0.3A \\
%DIFDELCMD < Total power and current consumption after lift-off & 60.48Wh and 2.1Ah \\
%DIFDELCMD < \midrule
%DIFDELCMD < Power system: Experiment includes batteries? Yes/No & No\\
%DIFDELCMD < Type of batteries: & \\
%DIFDELCMD < Number of batteries (and S x P): & \\
%DIFDELCMD < Capacity (1 battery): & Ah \\
%DIFDELCMD < Voltage (1 battery): & V \\
%DIFDELCMD < \bottomrule
%DIFDELCMD < \end{tabular}
%DIFDELCMD < %%%
\DIFdelendFL \DIFaddbeginFL \begin{tabular}{p{10cm} p{5cm}}
\toprule
\multicolumn{2}{l}{\textbf{BEXUS Electrical Interfaces}} \\
\midrule
E-Link Interface: E-Link required? Yes/No & Yes\\
Number of E-Link interfaces: & 1 \\
Number of required IP addresses: & 2 (one for the experiment and one for the ground station)\\
Data rate -- downlink (max. and average): & max: 100 Kbit/s  (average: 33 Kbit/s) \\
Data rate -- uplink (max. and average): & max: 10 Kbit/s \\
Interface type (RS-232, Ethernet): & Ethernet \\
\midrule
Power system: Gondola power required? Yes/No & Yes\\
Peak power and current consumption: & 65.05 W and 2.26 A \\
Average power and current consumption: & 23.11 W and 0.8 A \\
Total power and current consumption after lift-off & 138.67 Wh and 8.59 Ah \\
\midrule
Power system: Experiment includes batteries? Yes/No & No\\
Type of batteries: & \\
Number of batteries (and S x P): & \\
Capacity (1 battery): & Ah \\
Voltage (1 battery): & V \\
\bottomrule
\end{tabular}
\DIFaddendFL \end{table}
\clearpage

%DIF > %%%%%%%%%%%%%%%%%%%%%%%%%%%%%%%%%%%%%%%%%%%%%%%%%%%%%%%%%%%%%%%%%%%%%%%%%%%%%%
\subsection{Launch Site Requirements}
This is a very preliminary assessment and might change in later versions of the SED.
\DIFdelbegin %DIFDELCMD < \begin{table}[h]
%DIFDELCMD < %%%
\DIFdelendFL \DIFaddbeginFL \begin{table}[htbp]
\DIFaddendFL \centering
\caption{Launch Site Requirements}
\label{tab:launch_site_req}
\begin{tabular}{lll}
\toprule
\textbf{Category} & \textbf{Item} & \textbf{Amount} \\
\midrule
Workstation&Table & 3 \\
 &Chair & 10 \\
 &Power Supply 0-30V & 3\\
 &Power Extension Cord & 2 \\

\bottomrule
\end{tabular}
\end{table}
\clearpage

%DIF > %%%%%%%%%%%%%%%%%%%%%%%%%%%%%%%%%%%%%%%%%%%%%%%%%%%%%%%%%%%%%%%%%%%%%%%%%%%%%%
\subsection{Flight Requirements}
\DIFdelbegin %DIFDELCMD < 

%DIFDELCMD < \begin{table}[h]
%DIFDELCMD < %%%
\DIFdelendFL \DIFaddbeginFL \begin{table}[htbp]
\DIFaddendFL \centering
\caption{Flight Requirements applicable to BEXUS}
\label{tab:flight_req_bexus}
\DIFdelbeginFL %DIFDELCMD < \begin{tabular}{p{5cm} p{10cm}}
%DIFDELCMD < \toprule
%DIFDELCMD < \textbf{Optimal altitude} & $>$25km\\
%DIFDELCMD < \textbf{Preferred path} & ascent \\
%DIFDELCMD < \textbf{Minimum float time} & 2 hours\\
%DIFDELCMD < \textbf{Ground track length} & $>$200km\\
%DIFDELCMD < \textbf{Light/dark condition} & no preference but should be uniform suring flight\\
%DIFDELCMD < \bottomrule
%DIFDELCMD < \end{tabular}
%DIFDELCMD < %%%
\DIFdelendFL \DIFaddbeginFL \begin{tabular}{p{5cm} p{10cm}}
\toprule
\textbf{Optimal altitude} & 22km \\
\textbf{Preferred path} & ascent \\
\textbf{Minimum float time} & not important\\
\textbf{Ground track length} & not important\\
\textbf{Light/dark condition} & no preference but should be uniform during ascent if possible\\
\bottomrule
\end{tabular}
\DIFaddendFL \end{table}
\clearpage

%DIF > %%%%%%%%%%%%%%%%%%%%%%%%%%%%%%%%%%%%%%%%%%%%%%%%%%%%%%%%%%%%%%%%%%%%%%%%%%%%%%
\subsection{Accommodation Requirements}
\DIFaddbegin 

\DIFaddend The experiment needs access to fresh air and the gas in- and outlets should not be covered by the canvas\DIFdelbegin \DIFdel{and they }\DIFdelend \DIFaddbegin \DIFadd{. They }\DIFaddend should be directed away from the vehicle and other experiments to minimize interference \DIFdelbegin \DIFdel{. 
}\DIFdelend \DIFaddbegin \DIFadd{and contamination of the intaken air. Therefore the experiment should be accommodated in a corner of the gondola.
}\DIFaddend 

%DIF > %%%%%%%%%%%%%%%%%%%%%%%%%%%%%%%%%%%%%%%%%%%%%%%%%%%%%%%%%%%%%%%%%%%%%%%%%%%%%%
\section{Preparation and Test \DIFdelbegin \DIFdel{Activities }\DIFdelend \DIFaddbegin \DIFadd{Activity }\DIFaddend at Esrange}
\DIFaddbegin \label{sec:ESRANGE_activity}

\DIFaddend This will be filled out in later versions of the SED.


%DIF < \begin{table}[h]
%DIF > \begin{table}[htbp]
%\centering
%\caption{List of planned activities}
%\label{tab:planned_activities}
%\begin{tabular}{lllll}
%\toprule
%\textbf{Time/Day} & \textbf{Main Task} & \textbf{Description} & \textbf{Responsible} & \textbf{Duration [h:m]} \\
%\midrule
%& & & & \\
%\bottomrule
%\end{tabular}
%\end{table}
\clearpage


%DIF > %%%%%%%%%%%%%%%%%%%%%%%%%%%%%%%%%%%%%%%%%%%%%%%%%%%%%%%%%%%%%%%%%%%%%%%%%%%%%%
\section{Timeline for Countdown and Flight}
\DIFaddbegin \label{sec:timeline}

\DIFaddend This will be filled out in later versions of the SED.

%DIF < \begin{table}[h]
%DIF > \begin{table}[htbp]
%\centering
%\caption{Example of timeline of the experiment events}
%\label{tab:timeline}
%\begin{tabular}{lll}
%\toprule
%\textbf{Time [s]} & \textbf{EVENT} & \textbf{Function} \\
%\midrule
%T - .. s & POWER ON & \\
%& & \\
%& & \\
%& & \\
%\bottomrule
%\end{tabular}
%\end{table}
\clearpage



\section{Post-Flight Activities}
\DIFaddbegin 

\DIFaddend This will be filled out in later versions of the SED.

%Celebrate!

\DIFdelbegin \chapter{\DIFdel{Data Analysis And Results}}
%DIFAUXCMD
\addtocounter{chapter}{-1}%DIFAUXCMD
\DIFdelend %DIF > %%%%%%%%%%%%%%%%%%%%%%%%%%%%%%%%%%%%%%%%%%%%%%%%%%%%%%%%%%%%%%%%%%%%%%%%%%%%%%\newpage

\DIFaddbegin \chapter{\DIFadd{DATA ANALYSIS AND RESULTS}}
\label{sec:data_analysis_results}

\DIFaddend \section{Data Analysis Plan}

The focus of attention for \DIFdelbegin \DIFdel{\mbox{%DIFAUXCMD
\ac{mirage} }\hskip0pt%DIFAUXCMD
}\DIFdelend \DIFaddbegin \DIFadd{\mbox{%DIFAUXCMD
\gls{mirage} }\hskip0pt%DIFAUXCMD
}\DIFaddend is to demonstrate the applicability of the \DIFdelbegin \DIFdel{\mbox{%DIFAUXCMD
\ac{ndir} }\hskip0pt%DIFAUXCMD
}\DIFdelend \DIFaddbegin \DIFadd{\mbox{%DIFAUXCMD
\gls{ndir} }\hskip0pt%DIFAUXCMD
}\DIFaddend sensor for low-cost and low-complexity methane measurements in the stratosphere.
Since the biggest trade-off is resolution in this endeavour, \DIFdelbegin \DIFdel{we will assess }\DIFdelend the uncertainty of \DIFdelbegin \DIFdel{our measurement }\DIFdelend \DIFaddbegin \DIFadd{the measurement will be assessed }\DIFaddend with great care.
A lot of this will be done before the flight by carefully calibrating the sensor in the temperature and pressure range \DIFdelbegin \DIFdel{we plan on setting in }\DIFdelend \DIFaddbegin \DIFadd{of }\DIFaddend the measurement chamber for different methane concentrations and humidity levels.
If possible, \DIFdelbegin \DIFdel{we will also train }\DIFdelend a machine learning algorithm \DIFdelbegin \DIFdel{in this process }\DIFdelend \DIFaddbegin \DIFadd{will be trained }\DIFaddend to reduce the impact of noise on \DIFdelbegin \DIFdel{our measurements.
}%DIFDELCMD < 

%DIFDELCMD < %%%
\DIFdelend \DIFaddbegin \DIFadd{the measurements, similar to \mbox{%DIFAUXCMD
\cite{gia_accurate_2025}}\hskip0pt%DIFAUXCMD
. }\DIFaddend This is an important prerequisite for the quality of the data \DIFdelbegin \DIFdel{we can collect }\DIFdelend \DIFaddbegin \DIFadd{collected }\DIFaddend during the flight.

To further demonstrate the reliability of \DIFdelbegin \DIFdel{our }\DIFdelend \DIFaddbegin \DIFadd{the }\DIFaddend measurement data, the focus of post-flight data analysis is on comparison of the vertical methane profile from ascent and descent with other sources of methane data.
For this \DIFdelbegin \DIFdel{we want }\DIFdelend \DIFaddbegin \DIFadd{it is intended }\DIFaddend to use:

\begin{enumerate}
    \item Radiative transfer models (like RTTOV)
    \item Chemistry and Transport models (like MOCAGE)
    \item Satellite data
    \item In-situ measurement data from other groups
\DIFdelbegin %DIFDELCMD < \item %%%
\item%DIFAUXCMD
\DIFdel{In-situ measurement with LIDAR at ESRANGE
}\DIFdelend \end{enumerate}

\DIFdelbegin \DIFdel{\mbox{%DIFAUXCMD
\ac{rtm} }\hskip0pt%DIFAUXCMD
}\DIFdelend \DIFaddbegin \DIFadd{\mbox{%DIFAUXCMD
\gls{rtm} }\hskip0pt%DIFAUXCMD
}\DIFaddend like RTTOV can provide reference concentration profiles of atmospheric gases like methane and have been compared to \DIFdelbegin \DIFdel{\mbox{%DIFAUXCMD
\ac{tdlas} }\hskip0pt%DIFAUXCMD
}\DIFdelend \DIFaddbegin \DIFadd{\mbox{%DIFAUXCMD
\gls{tdlas} }\hskip0pt%DIFAUXCMD
}\DIFaddend measurements before.
\DIFdelbegin \DIFdel{\mbox{%DIFAUXCMD
\ac{ctm} }\hskip0pt%DIFAUXCMD
}\DIFdelend \DIFaddbegin \DIFadd{\mbox{%DIFAUXCMD
\gls{ctm} }\hskip0pt%DIFAUXCMD
}\DIFaddend like MOCAGE assimilate satellite data into their calculations and give an additional gauge of what \DIFdelbegin \DIFdel{we should expect for our data }\DIFdelend \DIFaddbegin \DIFadd{should be expected for the data collected by MIRAGE}\DIFaddend , although they have shown worse correspondence with in-situ measurements \cite{joly2020}.
The use of satellite data in addition to their indirect use via \DIFdelbegin \DIFdel{\mbox{%DIFAUXCMD
\ac{ctm} }\hskip0pt%DIFAUXCMD
}\DIFdelend \DIFaddbegin \DIFadd{\mbox{%DIFAUXCMD
\gls{ctm} }\hskip0pt%DIFAUXCMD
}\DIFaddend needs to be assessed.
This and the familiarisation with \DIFdelbegin \DIFdel{\mbox{%DIFAUXCMD
\ac{rtm} }\hskip0pt%DIFAUXCMD
and \mbox{%DIFAUXCMD
\ac{ctm} }\hskip0pt%DIFAUXCMD
}\DIFdelend \DIFaddbegin \DIFadd{\mbox{%DIFAUXCMD
\gls{rtm} }\hskip0pt%DIFAUXCMD
and \mbox{%DIFAUXCMD
\gls{ctm} }\hskip0pt%DIFAUXCMD
}\DIFaddend will be part of a bachelor thesis project within \DIFdelbegin \DIFdel{\mbox{%DIFAUXCMD
\ac{mirage}}\hskip0pt%DIFAUXCMD
}\DIFdelend \DIFaddbegin \DIFadd{\mbox{%DIFAUXCMD
\gls{mirage}}\hskip0pt%DIFAUXCMD
}\DIFaddend , that is planned for several team members.

In-situ measurements were conducted with \DIFdelbegin \DIFdel{\mbox{%DIFAUXCMD
\ac{tdlas} }\hskip0pt%DIFAUXCMD
}\DIFdelend \DIFaddbegin \DIFadd{\mbox{%DIFAUXCMD
\gls{tdlas} }\hskip0pt%DIFAUXCMD
}\DIFaddend from ESRANGE in 2016 and provide a direct comparison with more expensive but also more accurate laser-based methods \cite{Khair2017}.
The comparison is of course limited because of the different times at which measurements were done.
\DIFdelbegin \DIFdel{We still expect }\DIFdelend \DIFaddbegin \DIFadd{It is still expected }\DIFaddend to get some interesting insights by including these sources into \DIFdelbegin \DIFdel{our }\DIFdelend \DIFaddbegin \DIFadd{the }\DIFaddend data analysis.
Another in-situ measurement campaign conducted over France with AirCore in 2018 compared their data to \DIFdelbegin \DIFdel{\mbox{%DIFAUXCMD
\ac{ftir} }\hskip0pt%DIFAUXCMD
}\DIFdelend \DIFaddbegin \DIFadd{\mbox{%DIFAUXCMD
\gls{ftir} }\hskip0pt%DIFAUXCMD
}\DIFaddend measurements \cite{RINGO}.
\DIFdelbegin \DIFdel{This is an interesting study for us as well, especially if ESRANGE can support us with LIDAR methane measurements. Whether this is possible is not clear yet but will be asked for.
}\DIFdelend %DIF > This is an interesting study for us as well, especially if ESRANGE can support us with LIDAR methane measurements.
%DIF > Whether this is possible is not clear yet but will be asked for.

\DIFdelbegin \DIFdel{As a secondary goal we want to investigate, whether a correlation between different landscapes and methane levels in the stratosphere can be drawn. Depending on the quality of our data and number of data points we collected during the floating phase of the flight, we will study fluctuations and compare them with the GPS data collected by the balloon. Since freshwater bodies, wetlands and industrial areas have higher methane emissions than for example forests, we are interested whether methane fluctuations in the stratosphere can be related to those different landscapes. Machine learning might be employed to help recognising those fluctuations.
}\DIFdelend %DIF > As a secondary goal we want to investigate, whether a correlation between %different landscapes and methane levels in the stratosphere can be drawn.
%DIF > Depending on the quality of our data and number of data points we collected %during the floating phase of the flight, we will study fluctuations and %compare them with the GPS data collected by the balloon.
%DIF > Since freshwater bodies, wetlands and industrial areas have higher methane %emissions than for example forests, we are interested whether methane %fluctuations in the stratosphere can be related to those different landscapes.
%DIF > Machine learning might be employed to help recognising those fluctuations.

%DIF > %%%%%%%%%%%%%%%%%%%%%%%%%%%%%%%%%%%%%%%%%%%%%%%%%%%%%%%%%%%%%%%%%%%%%%%%%%%%%%
\section{Launch Campaign}
\DIFaddbegin 

\DIFaddend This will be filled out in later versions of the \DIFdelbegin \DIFdel{\mbox{%DIFAUXCMD
\ac{sed}}\hskip0pt%DIFAUXCMD
}\DIFdelend \DIFaddbegin \DIFadd{\mbox{%DIFAUXCMD
\gls{sed}}\hskip0pt%DIFAUXCMD
}\DIFaddend .

%DIF > %%%%%%%%%%%%%%%%%%%%%%%%%%%%%%%%%%%%%%%%%%%%%%%%%%%%%%%%%%%%%%%%%%%%%%%%%%%%%%
\section{Results}

It is expected that \DIFdelbegin \DIFdel{we will receive }\DIFdelend less data points with a quality lower than other in-situ missions \DIFaddbegin \DIFadd{will be received by MIRAGE}\DIFaddend . However, if the calibration was done with the intended thoroughness, \DIFdelbegin \DIFdel{we expect that }\DIFdelend the results should be comparable and correspond well to \DIFdelbegin \DIFdel{\mbox{%DIFAUXCMD
\ac{rtm}}\hskip0pt%DIFAUXCMD
. We expect to get good insights into current }\DIFdelend \DIFaddbegin \DIFadd{\mbox{%DIFAUXCMD
\gls{rtm}}\hskip0pt%DIFAUXCMD
. In that case, MIRAGE would have demonstrated that NDIR technology can measure }\DIFaddend methane levels in the stratosphere and \DIFdelbegin \DIFdel{hope to }\DIFdelend \DIFaddbegin \DIFadd{hopefully }\DIFaddend inspire the further development and application of \DIFdelbegin \DIFdel{\mbox{%DIFAUXCMD
\ac{ndir} }\hskip0pt%DIFAUXCMD
}\DIFdelend \DIFaddbegin \DIFadd{\mbox{%DIFAUXCMD
\gls{ndir} }\hskip0pt%DIFAUXCMD
}\DIFaddend sensors for atmospheric monitoring.

\DIFdelbegin \DIFdel{There is a small chance that fluctuations of methane in the stratosphere can be connected to the emissions at ground level even if very precise measurement techniques are given. Especially, if we fly over different large areas like wetlands or lakes and forests, we might be able to detect changes in concentrations. Even if this goal fails, we will have collected methane concentration data with higher accuracy at about 25-30 km height due to the longer measurement time.
}\DIFdelend %DIF > There is a small chance that fluctuations of methane in the stratosphere can %be connected to the emissions at ground level even if very precise %measurement techniques are given. Especially, if we fly over different large %areas like wetlands or lakes and forests, we might be able to detect changes %in concentrations. Even if this goal fails, we will have collected methane %concentration data with higher accuracy at about 25-30 km height due to the %longer measurement time.

%DIF > %%%%%%%%%%%%%%%%%%%%%%%%%%%%%%%%%%%%%%%%%%%%%%%%%%%%%%%%%%%%%%%%%%%%%%%%%%%%%%
\section{Input to Campaign Report}
\DIFaddbegin 

\DIFaddend This section will be filled after the Campaign.

%DIF > %%%%%%%%%%%%%%%%%%%%%%%%%%%%%%%%%%%%%%%%%%%%%%%%%%%%%%%%%%%%%%%%%%%%%%%%%%%%%%
\section{Lessons Learned}
\DIFaddbegin 

\DIFaddend As a fresh and still growing team there were some lessons \DIFdelbegin \DIFdel{we could take away so far}\DIFdelend \DIFaddbegin \DIFadd{learned already}\DIFaddend . 
\begin{enumerate}
  \item \DIFdelbegin \DIFdel{We }\DIFdelend \DIFaddbegin \DIFadd{It }\DIFaddend took too much time \DIFdelbegin \DIFdel{to }\DIFdelend recruit new team members and \DIFdelbegin \DIFdel{still miss }\DIFdelend key team members \DIFaddbegin \DIFadd{are still missing; }\DIFaddend for example in software design.
  \item \DIFdelbegin \DIFdel{We found out that we made wrong }\DIFdelend \DIFaddbegin \DIFadd{Wrong }\DIFaddend assumptions regarding the accuracy of the K96 sensor \DIFaddbegin \DIFadd{were made early one}\DIFaddend . This should have been prevented by asking more questions earlier on and closer communication with our contacts from Senseair \DIFaddbegin \DIFadd{AB}\DIFaddend .
  \item \DIFdelbegin \DIFdel{We need to take holidays into consideration and make }\DIFdelend \DIFaddbegin \DIFadd{Holidays and }\DIFaddend more generous time estimates \DIFdelbegin \DIFdel{to }\DIFdelend \DIFaddbegin \DIFadd{need to be taken into consideration to }\DIFaddend prevent finishing crucial project steps like the \DIFdelbegin \DIFdel{\mbox{%DIFAUXCMD
\ac{sed} }\hskip0pt%DIFAUXCMD
}\DIFdelend \DIFaddbegin \DIFadd{\mbox{%DIFAUXCMD
\gls{sed} }\hskip0pt%DIFAUXCMD
}\DIFaddend under pressure, which can lead to worse quality.
  \item \DIFdelbegin \DIFdel{We have }\DIFdelend \DIFaddbegin \DIFadd{It it important }\DIFaddend to keep track of progress and check in with each sub team or team members working on specific tasks to prevent delays.
\end{enumerate}
Of course there are many more lessons of more practical nature. Everyone knows a little more about pumps and their limitations now, \DIFdelbegin \DIFdel{all of us }\DIFdelend \DIFaddbegin \DIFadd{everyone }\DIFaddend learned something about methane sensing and many more things like that. However, the biggest lesson is probably that the documentation for such a project takes a lot of time and thought but also makes you aware of many aspects that can be easily overlooked.

\DIFdelbegin \chapter{\DIFdel{Abbreviations}}
%DIFAUXCMD
\addtocounter{chapter}{-1}%DIFAUXCMD
\DIFdelend %DIF > %%%%%%%%%%%%%%%%%%%%%%%%%%%%%%%%%%%%%%%%%%%%%%%%%%%%%%%%%%%%%%%%%%%%%%%%%%%%%%\newpage

%DIF <  When adjusting the acronyms, make sure to add them according to the template below 
\DIFdelbegin %DIFDELCMD < \begin{acronym}[LONGESTACRONYM]
%DIFDELCMD < \acro{mirage}[MIRAGE]{Methane InfraRed Absorption Gas Experiment}
%DIFDELCMD < \acro{wbs}[WBS]{Work Breakdown Structure}
%DIFDELCMD < \acro{sed}[SED]{Student Experiment Document}
%DIFDELCMD < \acro{tdlas}[TDLAS]{Tunable Diode Laser Absorption Spectroscopy}
%DIFDELCMD < \acro{ftir}[FTIR]{Fourier Transform Infrared Spectroscopy}
%DIFDELCMD < \acro{ndir}[NDIR]{non-dispersive infrared}
%DIFDELCMD < \acro{ir}[IR]{infrared}
%DIFDELCMD < \acro{mcu}[MCU]{microcontroller unit}
%DIFDELCMD < \acro{irf}[IRF]{Institutet För Rymdfysik}
%DIFDELCMD < \acro{cog}[COG]{centre of gravity}
%DIFDELCMD < \acro{tbd}[TBD]{to be decided}
%DIFDELCMD < \acro{pcb}[PCB]{printed circuit board}
%DIFDELCMD < \acro{rtm}[RTMs]{radiative transfer models}
%DIFDELCMD < \acro{ctm}[CTMs]{chemistry and transport models}
%DIFDELCMD < \end{acronym}
%DIFDELCMD < %%%
\DIFdelend \DIFaddbegin \chapter{\DIFadd{ABBREVIATIONS AND REFERENCES}}
\label{sec:abbreviations_references}
\section{\DIFadd{Abbreviations}}
%DIF > \glsaddall %treats all acronyms as used
\printglossary[type=\acronymtype, style=compactacronyms, title={}, toctitle={}]
 %DIF > prints used acronyms
\DIFaddend 

%DIF < AIT -- Assembly, Integration and Test
%DIF < 
%DIF < ASAP -- as soon as possible
%DIF < 
%DIF < BoM -- Budget of Material
%DIF < 
%DIF < CDR -- Critical Design Review
%DIF < 
%DIF < COG -- Centre of Gravity
%DIF < 
%DIF < CRP -- Campaign Requirement Plan
%DIF < 
%DIF < DLR -- Deutsches Zentrum für Luft- und Raumfahrt
%DIF < 
%DIF < EAT -- Experiment Acceptance Test
%DIF < 
%DIF < EAR -- Experiment Acceptance Review
%DIF < 
%DIF < ECTS -- European Credit Transfer System
%DIF < 
%DIF < EIT -- Electrical Interface Test
%DIF < 
%DIF < EPM -- Esrange Project Manager
%DIF < 
%DIF < ESA -- European Space Agency
%DIF < 
%DIF < Esrange -- Esrange Space Center
%DIF < 
%DIF < ESTEC -- European Space Research and Technology Centre, ESA (NL)
%DIF < 
%DIF < ESW -- Experiment Selection Workshop
%DIF < 
%DIF < FAR -- Flight Acceptance Review
%DIF < 
%DIF < FST -- Flight Simulation Test
%DIF < 
%DIF < FRP -- Flight Requirement Plan
%DIF < 
%DIF < FRR -- Flight Readiness Review
%DIF < 
%DIF < GSE -- Ground Support Equipment
%DIF < 
%DIF < HK -- House Keeping
%DIF < 
%DIF < H/W -- Hardware
%DIF < 
%DIF < ICD -- Interface Control Document
%DIF < 
%DIF < I/F -- Interface
%DIF < 
%DIF < IPR -- Integration Progress Review
%DIF < 
%DIF < LO -- Lift Off
%DIF < 
%DIF < LT -- Local Time
%DIF < 
%DIF < LOS -- Line of Sight
%DIF < 
%DIF < Mbps -- Mega Bits per second
%DIF < 
%DIF < MCU -- Microcontroller Unit
%DIF < 
%DIF < MFH -- Mission Flight Handbook
%DIF < 
%DIF < MORABA -- Mobile Raketen Basis (DLR, EuroLaunch)
%DIF < 
%DIF < NDIR -- Non-Dispersive Infra-Red
%DIF < 
%DIF < OP -- Oberpfaffenhofen, DLR Center
%DIF < 
%DIF < PCB -- Printed Circuit Board (electronic card)
%DIF < 
%DIF < PDR -- Preliminary Design Review
%DIF < 
%DIF < PST -- Payload System Test
%DIF < 
%DIF < RBF -- Remove Before Flight
%DIF < 
%DIF < SED -- Student Experiment Documentation
%DIF < 
%DIF < SNSA -- Swedish National Space Agency
%DIF < 
%DIF < SODS -- Start Of Data Storage
%DIF < 
%DIF < SOE -- Start Of Experiment
%DIF < 
%DIF < SSC -- Swedish Space Corporation
%DIF < 
%DIF < STW -- Student Training Week
%DIF < 
%DIF < S/W -- Software
%DIF < 
%DIF < T -- Time before and after launch noted with + or -
%DIF < 
%DIF < TBC -- To be confirmed
%DIF < 
%DIF < TBD -- To be determined
%DIF < 
%DIF < WBS -- Work Breakdown Structure
%DIF < 
%DIF < ZARM -- Zentrum für Angewandte Raumfahrttechnologie und Mikrogravitation
%DIF > use \gls{asap} to get accronym, first acronym will show it spelled out, after only prints acronym with link to this list

\DIFdelbegin %DIFDELCMD < \bibliographystyle{plain}
%DIFDELCMD < \bibliography{ref}
%DIFDELCMD < %%%
\DIFdelend %DIF > %%%%%%%%%%%%%%%%%%%%%%%%%%%%%%%%%%%%%%%%%%%%%%%%%%%%%%%%%%%%%%%%%%%%%%%%%%%%%%
\DIFaddbegin \section{\DIFadd{References}}
\nocite{*}
\printbibliography[heading=none]
\DIFaddend 



%DIF <  Appendices
%DIF > References have to be cited to appear, this can be done with \cite{example_book}. 
%DIF > If you want references to appear without citing it first use \nocite{*} 
%DIF > above \printbibliography like above\newpage
%DIF > %%%%%%%%%%%%%%%%%%%%%%%%%%%%%%%%%%%%%%%%%%%%%%%%%%%%%%%%%%%%%%%%%%%%%%%%%%%%%%
\DIFaddbegin 

%DIF > appendix
\DIFaddend \appendix

\chapter{EXPERIMENT REVIEWS}
\DIFaddbegin \label{app:experiment_reviews}
%DIF > \section{Preliminary Design Review - PDR}
%DIF > \label{app:exprev_PDR}
\DIFaddend 

\DIFdelbegin \section{\DIFdel{Preliminary Design Review -- PDR}}
%DIFAUXCMD
\addtocounter{section}{-1}%DIFAUXCMD
\DIFdelend %DIF > \section{Critical Design Review - CDR}
%DIF > \label{app:exprev_CDR}

\DIFdelbegin \section{\DIFdel{Critical Design Review -- CDR}}
%DIFAUXCMD
\addtocounter{section}{-1}%DIFAUXCMD
\DIFdelend %DIF > \section{Integration Process Review - IPR}
%DIF > \label{app:exprev_IPR}

\DIFdelbegin \section{\DIFdel{Integration Progress Review -- IPR}}
%DIFAUXCMD
\addtocounter{section}{-1}%DIFAUXCMD
\DIFdelend %DIF > \section{Experiment Acceptance Review - EAR}
%DIF > \label{app:exprev_EAR}\newpage

\DIFdelbegin \section{\DIFdel{Experiment Acceptance Review -- EAR}}
%DIFAUXCMD
\addtocounter{section}{-1}%DIFAUXCMD
%DIFDELCMD < 

%DIFDELCMD < %%%
\DIFdelend \chapter{OUTREACH AND MEDIA COVERAGE}
\DIFaddbegin \label{app:outreach}

\DIFaddend \textbf{Website} \\
Website Coming soon

\textbf{Instagram:}
\url{https://www.instagram.com/riseupp_uu/?igsh=MXhsM3lsZmMyMjh4ZA%3D%3D#}

\textbf{LinkedIn:}
\url{https://www.linkedin.com/company/riseupp-uu/posts/?feedView=all}


\textbf{Logo:} \\
The \DIFaddbegin \DIFadd{RiseUpp }\DIFaddend logo was designed in July \DIFdelbegin \DIFdel{2025
}\DIFdelend \DIFaddbegin \DIFadd{2025. The \mbox{%DIFAUXCMD
\gls{mirage} }\hskip0pt%DIFAUXCMD
logo is under development.
}\DIFaddend 

\begin{figure}[H]
    \centering
    \DIFdelbeginFL %DIFDELCMD < 
\includegraphics[width=0.5\textwidth]{Graphics/Logos/RiseUpp_logo.png}
%DIFDELCMD <     %%%
\DIFdelendFL \DIFaddbeginFL 
\includegraphics[width=0.5\textwidth]{0_images/figures/RiseUpp_logo.png}
    \DIFaddendFL \caption{RiseUpp Logo}
    \label{fig:riseupp_logo}
\end{figure}

The RiseUpp association is \DIFdelbegin \DIFdel{very soon }\DIFdelend an official member of Astronomisk Ungdom association and a website will be created under their domain.

\DIFdelbegin \DIFdel{\mbox{%DIFAUXCMD
\ac{mirage} }\hskip0pt%DIFAUXCMD
}\DIFdelend \DIFaddbegin \DIFadd{\mbox{%DIFAUXCMD
\gls{mirage} }\hskip0pt%DIFAUXCMD
}\DIFaddend patches and merchandise, such as shirts, hoodies and caps will be designed in the future.

\DIFdelbegin \DIFdel{\mbox{%DIFAUXCMD
\ac{mirage} }\hskip0pt%DIFAUXCMD
}\DIFdelend \DIFaddbegin \DIFadd{\mbox{%DIFAUXCMD
\gls{mirage} }\hskip0pt%DIFAUXCMD
}\DIFaddend aims to appear in local newsoutlets, such as Techna (student magazine), Uppsala Nya Tidning (local newspaper) and Student radio (local radiostation)

The \DIFdelbegin \DIFdel{\mbox{%DIFAUXCMD
\ac{mirage} }\hskip0pt%DIFAUXCMD
}\DIFdelend \DIFaddbegin \DIFadd{\mbox{%DIFAUXCMD
\gls{mirage} }\hskip0pt%DIFAUXCMD
}\DIFaddend team will be creating a project movie documenting the progress of the project.\DIFaddbegin \newpage

\DIFaddend \chapter{ADDITIONAL TECHNICAL INFORMATION}
\DIFaddbegin \label{app:add_tech_info}

\DIFaddend \begin{figure}[htbp]
    \centering
    \DIFdelbeginFL %DIFDELCMD < \includegraphics[angle=90, width=0.60\textheight]{Graphics/Figures/Mirage.CADdrawing.v1.png}
%DIFDELCMD <     %%%
\DIFdelendFL \DIFaddbeginFL \includegraphics[angle=90, width=0.60\textheight]{0_images/figures/tech_drawing.png}
    \DIFaddendFL \caption{\DIFdelbeginFL \DIFdelFL{Mechanical }\DIFdelendFL \DIFaddbeginFL \DIFaddFL{Technical }\DIFaddendFL drawing of the \DIFdelbeginFL \DIFdelFL{experiment structure}\DIFdelendFL \DIFaddbeginFL \DIFaddFL{system as a whole.}\DIFaddendFL }
    \DIFdelbeginFL %DIFDELCMD < \label{fig:CAD}
%DIFDELCMD < %%%
\DIFdelendFL \DIFaddbeginFL \label{fig:cad_system_overview}
\DIFaddendFL \end{figure}

\DIFaddbegin \begin{figure}[htbp]
    \centering
    \includegraphics[angle=90, width=0.60\textheight]{0_images/figures/tech_drawing_2.png}
    \caption{\DIFaddFL{Technical drawing of the pressurized chamber.}}
    \label{fig:cad_pressure_chamber_1}
\end{figure}



\begin{figure}[htbp]
    \centering
    \includegraphics[angle=90, width=0.60\textheight]{0_images/figures/tech_drawing_4.png}
    \caption{\DIFaddFL{Technical drawing of the pressurized chamber.}}
    \label{fig:cad_pressure_chamber_2}
\end{figure}

\begin{figure}[htbp]
    \centering
    \includegraphics[angle=90, width=0.60\textheight]{0_images/figures/tech_drawing_3.png}
    \caption{\DIFaddFL{Technical drawing of experiment enclosure. Drawings are representative and the thickness of both “Main Body” and “Main Lid” will be reduced to 2 mm by buying COTS parts.}}
    \label{fig:cad_enclosure_1}
\end{figure}

\begin{figure}[htbp]
    \centering
    \includegraphics[angle=90, width=0.60\textheight]{0_images/figures/tech_drawing_5.png}
    \caption{\DIFaddFL{Technical drawing of experiment enclosure. Drawings are representative and the thickness of both “Main Body” and “Main Lid” will be reduced to 2 mm by buying COTS parts.}}
    \label{fig:cad_enclosure_2}
\end{figure}

\begin{figure}[htbp]
    \centering
    \includegraphics[angle=90, width=0.60\textheight]{0_images/figures/tech_drawing_6.png}
    \caption{\DIFaddFL{Drawing of the internal layout.}}
    \label{fig:cad_internal_layout}
\end{figure}\newpage

\DIFaddend \chapter{CHECKLISTS}
\DIFaddbegin \label{app:checklists}
\newpage

\chapter{\DIFadd{PLANNING}}\newpage

\chapter{\DIFadd{MATERIAL SAFETY DATA SHEETS}}\newpage 

\chapter{\DIFadd{THERMAL ANALYSIS BACKGROUND}}
\label{app:thermal_background}
\DIFadd{Assuming energy balance
}\begin{align}
    \DIFadd{\dot{E}_{\mathrm{system}}=\dot{E}_{\mathrm{in}}-\dot{E}_{out}\qc
}\end{align}
\DIFadd{the first law of thermodynamics 
}\begin{align}
    \DIFadd{\Delta E_\mathrm{system}=Q-W\qc
}\end{align}
\DIFadd{can be used to relate heat fluxes in and out of a system to each other. Here, $Q$ is the net heat added to the system and $W$ is the work done on the system. Assuming that no work is being done (no volume expansion), and that only the internal energy (as opposed to e.g. kinetic energy) is of interest, this becomes
}\begin{align}
    \DIFadd{\Delta }{\DIFadd{E}}\DIFadd{_\mathrm{system}}&\DIFadd{=Q}\\ 
    &\DIFadd{=m_s\cdot c_{p,s} \cdot \Delta }{\DIFadd{T}}\DIFadd{_s\qc
}\end{align}
\DIFadd{where the subscript $_s$ signals system parameters, $m$ is the mass, $c_p$ the specific heat capacity and $T$ the temperature. Finally, this gives 
}\begin{align}
    \DIFadd{m_s\cdot c_{p,s} \cdot \dot{T}_s }&\DIFadd{= \dot{E}_{\mathrm{in}}-\dot{E}_{out}}\\ 
    &\DIFadd{= \dot{Q}_\mathrm{in}-\dot{Q}_\mathrm{out}
}\end{align}
\DIFadd{For simplicity, in the following the $\dot{}$ is dropped from $Q$. 
}

\DIFadd{Potential contributors to the heat flux into the system $Q_\mathrm{in}$ are 
}\begin{itemize}
    \item \DIFadd{heat generated by electronics 
    }\item \DIFadd{heat generated by pumps 
    }\item \DIFadd{heat generated by pressurisation of air 
    }\item \DIFadd{heat through absorbance of solar radiation,
}\end{itemize}
\DIFadd{while contributors to heat flux out of the system $Q_\mathrm{out}$ are 
}\begin{itemize}
    \item \DIFadd{convective cooling via ambient air
    }\item \DIFadd{radiative cooling via emittance of infrared radiation.
}\end{itemize}
\DIFadd{In the following, a derivation of the individual contributions will follow. 
}

\paragraph{\DIFadd{The heat generated by compressed air}} \DIFadd{must be analysed in two stages, as there are two compression stages in the system. The first is facilitated by the vacuum pumps, pressurising ambient air to 1 bar, the second by the compressor, pressurising to 3 bar. The discharge temperature is defined as 
}\begin{align}
    \DIFadd{T_\mathrm{out}}&\DIFadd{=T_\mathrm{in}\mbox{%DIFAUXCMD
\qty}\hskip0pt%DIFAUXCMD
(\frac{p_\mathrm{out}}{p_\mathrm{in}})^{\frac{n-1}{n}}\qc
}\end{align}
\DIFadd{where $T_\mathrm{out}$ is the discharge temperature (air temperature leaving the pumps), $T_\mathrm{in}$ is the temperature of the air on the suction side, $p_\mathrm{in}$ is the air pressure on the suction side and $p_\mathrm{out}$ is the discharge pressure. $n$ is the polytropic factor, determined by the pump and gas. If the process were isentropic, the factor would be equal to the isentropic factor of air $\gamma=1.4$. For the diaphragm pump, due to negligible friction, $n<\gamma$ as the air is dumping part of its heat into the pump.
Conversely, the compressor, as a piston pump, might generate significant heat, further heating the air, such that $n>\gamma$.
Although the discharge temperature can reach up to 500°C (assuming $n=\gamma$), the actual heat capacity of that air diminishes, as the air density decreases. Thus the total heat energy carried by the air can be calculated as  
}\begin{align}
    \DIFadd{Q_\mathrm{heat} }&\DIFadd{= \dot{m}c_p\Delta T}\\ 
    &\DIFadd{= \rho I_\mathrm{pump} c_p\Delta T\qc
}\end{align}
\DIFadd{where $\rho$ is the air density, $I_{pump}$ the volumetric pump flow rate, $c_p$ the specific heat of air, and $\Delta T$ the temperature difference between suction side and discharge side. To determine the final discharge temperature, it needs to be calculated for the diaphragm pumps and based on that for the compressor.
}

\paragraph{\DIFadd{The heat generated by the pumps}} \DIFadd{themselves can be estimated from their power requirements and an efficiency factor. The diaphragm pumps will probably have a higher effiency than the piston compressor pump, since they generate less friction. The total heat is then estimated as 
}\begin{align}
    \DIFadd{Q_\mathrm{heat} }&\DIFadd{= \epsilon P_\mathrm{draw}\qc
}\end{align}
\DIFadd{where $\epsilon$ is the effiency factor ranging between 0 and 1 and $P_\mathrm{draw}$ is the electrical power consumption of the pump. Note, that the operations schedule implies $P_\mathrm{draw}$ to vary with time.
}

\paragraph{\DIFadd{The heat generated by other electric components}} \DIFadd{is similarly calculated. Given in \tref{tab:power_consumption} is the total power draw from the BEXUS system, as well as the total consumption. This is already assuming an effiency factor, such that the difference between those values is assumed to be emitted as heat. Thus,
}\begin{align}
     \DIFadd{Q_\mathrm{heat} }&\DIFadd{= P_\mathrm{draw}- P_\mathrm{consumed}   
}\end{align}

\paragraph{\DIFadd{The heat absorbed through solar radiation}} \DIFadd{depends on the system material, surface, and dimensions, as well as on weather conditions like cloudiness and time/solar zenith angle and altitude. In general the total heat flux is comprised of three components: direct flux from the sun, secondary solar flux reflected from Earth, clouds, and atmosphere and the heat of the earth in the \mbox{%DIFAUXCMD
\gls{ir}}\hskip0pt%DIFAUXCMD
. Thus,
}\begin{align}
    \DIFadd{q }&\DIFadd{= q_\mathrm{sun}+q_\mathrm{reflected}+q_\mathrm{IR\,earth}}\\
    &\DIFadd{=\alpha_s (I_s \cos{\theta}+ I_s a F_{se} ) + \alpha_\mathrm{IR}(q_\mathrm{earth} F_{se})\qc
}\end{align}
\DIFadd{where $\alpha_s$ is the solar absorptivity of the surface, $I_s$ is the solar intensity received, $\theta$ is the angle between surface normal and sun, $a$ is the albedo coefficient of the Earth and atmosphere, $F_{se}$ is the view factor relating the amount of surface viewing in direction of the Earth, $\alpha_\mathrm{IR}$ is the \mbox{%DIFAUXCMD
\gls{ir} }\hskip0pt%DIFAUXCMD
absorptivity of the surface, and $q_\mathrm{earth}$ is the \mbox{%DIFAUXCMD
\gls{ir} }\hskip0pt%DIFAUXCMD
heat flux emitted by the Earth. Most of these values are or can be assumed to be roughly constant with respect to height. As the chassis and chamber are made out of aluminum the following analysis assumes $\alpha_{IR}=\dos{\epsilon}=0.1$ and $\alpha_s=0.4$.
The most dramatic change with altitude is seen in $I_s$.}\par 
\DIFadd{The atmosphere is absorbing part of the solar radiation incident on the top atmospheric layers $G_{sc}$. This can be modeled via the Lambert-Beer law as 
}\begin{align}
    \DIFadd{I_s(h) = G_{sc}\tau_0^{\frac{p(h)}{p_0}\frac{1}{\cos{\phi}+  0.505072 * (96.07995 - \phi) ^ {-1.6364}
}}\qc
}\end{align}
\DIFadd{where $\tau_0$ is the transmittance of solar light at ground level, $p_0$ is the atmospheric pressure at ground level, $p(h)$ is the pressure at a given altitude, and $\phi$ is the zenith angle. This formula assumes zero clouds. Since the minimum expected zenith angle in Kiruna in October is 77.11°, this value will be used by the following analysis.}\par 
\DIFadd{The final heat load is determined by multiplying the heat load for individual surfaces and summing over them, i.e.
}\begin{align}
    \DIFadd{Q_\mathrm{heat}=\sum_i q_i A_i.
}\end{align}

\paragraph{\DIFadd{The heat lost via convective cooling}} \DIFadd{via ambient air depends on the ambient air density and on system geometry. In general, 
}\begin{align}
    \DIFadd{Q_\mathrm{heat,A} = hA(T_\mathrm{system}-T_\mathrm{ambient})\qc
}\end{align}
\DIFadd{where $A$ is the area of the respective surface, $h$ is the convection coefficient and $T_\mathrm{ambient}$ is the temperature of the ambient air. The convection coefficient depends on the air density and can be modelled as 
}\begin{align}
    \DIFadd{h = h_0 \mbox{%DIFAUXCMD
\qty}\hskip0pt%DIFAUXCMD
(\frac{p}{p_0})^n\qc
}\end{align}
\DIFadd{where $h_0$ defines some reference convection coefficient on the ground and is usually $h_0=15\,\mathrm{W/m^2/K}$. 
$n$ is usually set to 0.5 for balloon applications with laminar flow (true during ascend and floating phase, if the wind speed is not extremely high). The total heat load is again defined by adding the individual heat loads of the surfaces together.
}

\paragraph{\DIFadd{The heat lost via blackbody radiation}} \DIFadd{is defined via the Stefan-Boltzmann law 
}\begin{align}
    \DIFadd{Q_\mathrm{heat,A}}&\DIFadd{= A\varepsilon\sigma (T_\mathrm{system}^4-T_\mathrm{ambient}^4)\qc
}\end{align}
\DIFadd{where $\varepsilon$ is the emissivity and $\sigma$ is the Stefan-Boltzmann constant. The total heat load is again defined by adding the individual heat loads of the surfaces together.
}

\paragraph{\DIFadd{The lumped system approach}} \DIFadd{is used to model subsystems with thermally distinct behaviours. Here, this means that the system is split thermally into two parts: 
}\begin{itemize}
    \item \DIFadd{the box/chassis containing the electronics and components, 
    }\item \DIFadd{the pressure chamber within,
    }\item \DIFadd{and potentially isolation if considered (see cold and hot case in \sref{sec: thermal cold case} and \sref{sec: thermal hot case}, respectively).
}\end{itemize}
\DIFadd{The division of the system into the lumps chassis and chamber is done as the pressure chamber is to be kept at a different temperature than the rest of the system. First, the temperature of the compartment/chassis containing all electronics and components is calculated based on a constant temperature of the pressure chamber. Then, the heat flux necessary to keep the pressure chamber at the desired temperature is determined. This means that the heat flux between chamber and chassis is added as a heat load to both. 
This transfer flux is defined as 
}\begin{align}
    \DIFadd{Q_\mathrm{transfer} = F_g (T_\mathrm{chamber}-T_\mathrm{chassis}) + \epsilon_\mathrm{chamber}\sigma A_\mathrm{chamber}\mbox{%DIFAUXCMD
\qty}\hskip0pt%DIFAUXCMD
(T_\mathrm{chamber}^4-T_\mathrm{chassis}^4)
}\end{align}
\DIFadd{Here, $F_g$ is a coefficient determining the efficiency of conductive and radiative heat transfer between chassis and chamber. While this coefficient depends on the exact geometry and material of conncecting bolts and surfaces, as well as on the air density inside the chassis, it is approximated at $F_g\approx 1.5$ W/K. The second term is responsible for modeling the radiative heat transfer between chamber and chassis. As the chassis has a much greater surface area than the chamber, it is modelled as a heat sink. This means that the efficiency of the radiative transfer is solely determined by chamber emissivity $\epsilon$ and surface area $A$.
}

\paragraph{\DIFadd{Insulation}} \DIFadd{is modelled as a white styrofoam box of 5 cm thickness, completely encapsulating the chassis. The chassis is no longer subject to solar radiation, air convection, and can no longer loose heat via blackbody radiation. Thus, $Q_\mathrm{rad\,in}=Q_\mathrm{rad\,out}=Q_\mathrm{conv}=0$. Instead, the chassis is interacting with the styrofoam box via 
}\begin{align}
    \DIFadd{Q_\mathrm{insulation}=\kappa_\mathrm{styro}\frac{A_\mathrm{styro}}{d_\mathrm{styro}}\mbox{%DIFAUXCMD
\qty}\hskip0pt%DIFAUXCMD
(T_\mathrm{styro}-T_\mathrm{chassis})\qc
}\end{align} 
\DIFadd{where $\kappa_\mathrm{styro}$ is the thermal conductivity and set to be 0.033 $\mathrm{W\,m^{-1}\,K^{-1}}$, $A_\mathrm{styro}$ is the insulation surface area, $d_\mathrm{styro}$ its thickness and $T_\mathrm{styro}$ its outer temperature. The outer temperature in turn is defined via 
}\begin{align}
    \DIFadd{Q_\mathrm{rad\,in} = Q_\mathrm{rad\,out}+Q_\mathrm{conv} + Q_\mathrm{insulation}.
}\end{align}
\DIFadd{This means that a steady-state temperature of the styrofoam is assumed. This holds true if the styrofoam's thermal capacity is small compared to that of the chassis. Since styrofoam is lightweight and the chassis is made out of aluminium this is justified.
The \mbox{%DIFAUXCMD
\gls{ir} }\hskip0pt%DIFAUXCMD
emissivity is set to $\epsilon=\alpha_\mathrm{IR}=0.6$ and the solar absorptivity to $\alpha_s=0.15$.}\par 

\DIFadd{All in all this means that for the steady state assumptions used in the cold and hot cases, the following two equations need to be solved:
}\begin{align}
    \DIFadd{Q_\mathrm{ele}+Q_\mathrm{pump}+Q_\mathrm{transfer} + Q_\mathrm{rad\,in} + Q_\mathrm{insulation} }&\DIFadd{= Q_\mathrm{conv} + Q_\mathrm{rad\,out} \label{eq:heat_balance_compartment}}\\ 
    \DIFadd{Q_\mathrm{thermal\,system}+Q_\mathrm{compressed\,air} }&\DIFadd{=  Q_\mathrm{transfer} \label{eq:heat_balance_chamber}
}\end{align} 
\DIFadd{\eqref{eq:heat_balance_compartment} describes the heat balance of the bigger chassis and \eqref{eq:heat_balance_chamber} describes the heat balance of the pressure chamber within. Note, that depending on the assumptions in the cold and hot cases, some of the terms might be 0.%DIF > %%%%%%%%%%%%%%%%%%%%%%%%%%%%%%%%%%%%%%%%%%%%%%%%%%%%%%%%%%%%%%%%%%%%%%%%%%%%%%
}\DIFaddend 

\end{document}
